\documentclass[10pt,a4paper]{article}
\usepackage[utf8]{inputenc}
\usepackage[T1]{fontenc}
\usepackage{lmodern}
\usepackage[margin=2.3cm]{geometry}
\usepackage{microtype}
\usepackage{booktabs}
\usepackage{longtable}
\usepackage{graphicx}
\usepackage{rotating}
\usepackage{array}
\usepackage{enumitem}
\usepackage{xcolor}
\usepackage{titlesec}
\usepackage{fancyhdr}
\usepackage{parskip}
\usepackage[numbers,sort&compress,square]{natbib}
\usepackage[hidelinks]{hyperref}
\usepackage{xurl}
\setlength{\emergencystretch}{3em}
\hypersetup{colorlinks=true, linkcolor=teal!60!black, citecolor=teal!55!black, urlcolor=blue!55!black,
  pdftitle={Earthquake Catalogues of Aotearoa New Zealand 1960-2026}, pdfauthor={Kenny Graham}}

\DeclareUnicodeCharacter{0100}{\=A}\DeclareUnicodeCharacter{0101}{\=a}
\DeclareUnicodeCharacter{0112}{\=E}\DeclareUnicodeCharacter{0113}{\=e}
\DeclareUnicodeCharacter{012A}{\=I}\DeclareUnicodeCharacter{012B}{\=\i}
\DeclareUnicodeCharacter{014C}{\=O}\DeclareUnicodeCharacter{014D}{\=o}
\DeclareUnicodeCharacter{016A}{\=U}\DeclareUnicodeCharacter{016B}{\=u}
\DeclareUnicodeCharacter{00F3}{\'o}\DeclareUnicodeCharacter{00E9}{\'e}\DeclareUnicodeCharacter{00FC}{\"u}
\DeclareUnicodeCharacter{2013}{--}\DeclareUnicodeCharacter{2014}{--}
\DeclareUnicodeCharacter{2018}{\textquoteleft}\DeclareUnicodeCharacter{2019}{\textquoteright}
\DeclareUnicodeCharacter{201C}{\textquotedblleft}\DeclareUnicodeCharacter{201D}{\textquotedblright}
\DeclareUnicodeCharacter{2265}{$\geq$}\DeclareUnicodeCharacter{2264}{$\leq$}
\DeclareUnicodeCharacter{00D7}{$\times$}\DeclareUnicodeCharacter{00B1}{$\pm$}
\DeclareUnicodeCharacter{00B0}{$^{\circ}$}\DeclareUnicodeCharacter{00B7}{$\cdot$}
\DeclareUnicodeCharacter{2248}{$\approx$}\DeclareUnicodeCharacter{2212}{$-$}\DeclareUnicodeCharacter{03BC}{$\mu$}

\definecolor{barblue}{RGB}{38,114,128}
\definecolor{rulegrey}{RGB}{120,120,120}
\newcolumntype{L}[1]{>{\raggedright\arraybackslash}p{#1}}
\titleformat{\section}{\Large\bfseries\color{teal!50!black}}{\thesection}{0.7em}{}
\titleformat{\subsection}{\large\bfseries}{\thesubsection}{0.6em}{}
\setlength{\parskip}{4pt}
\pagestyle{fancy}\fancyhf{}
\fancyhead[L]{\small\itshape Earthquake Catalogues of Aotearoa New Zealand, 1960--2026}
\fancyhead[R]{\small\thepage}
\renewcommand{\headrulewidth}{0.3pt}

\newcommand{\catentryrule}{\par\vspace{3pt}\noindent\textcolor{rulegrey}{\rule{\linewidth}{0.3pt}}\par\vspace{5pt}}
\newcommand{\kv}[2]{\item \textbf{#1:} #2}

\title{\vspace{-1.0cm}\Huge\bfseries Earthquake Catalogues of Aotearoa New Zealand, 1960--2026\\[6pt]
\Large A Curated Inventory of National, Regional, Temporary-Deployment, and Derived Catalogues\thanks{Each entry is a discovery record summarising what a catalogue contains and where to obtain it; confirm details against the cited source before quantitative use.}}
\author{Kenny Graham\\\small Earth Sciences New Zealand\\\small Aotearoa New Zealand FAIR Earthquake Catalogue Framework project}
\date{18 June 2026}

\begin{document}
\maketitle
\thispagestyle{fancy}

\begin{abstract}
\noindent New Zealand's earthquake-catalogue record is curated nationally by GeoNet (Earth Sciences New Zealand / GNS Science) but is complemented by a large, fragmented body of regional, temporary-deployment, volcano-seismic, slow-earthquake, and derived catalogues that are individually valuable yet difficult to discover, compare, and cite. This report presents a curated inventory of \textbf{81 distinct earthquake catalogues} that were developed for, or that cover, Aotearoa New Zealand and its offshore margins with coverage overlapping the period 1960--2026, with particular emphasis on \textbf{temporary seismic deployments} (passive-seismic experiments, ocean-bottom deployments, and aftershock rapid-response arrays). Each catalogue is profiled against the minimum-attributes \emph{Catalogue Submission Profile} proposed for a national \textquotedblleft catalogue of catalogues\textquotedblright, recording what it contains and how to obtain it. Of the 81 catalogues, 73 carry a resolvable DOI or stable dataset identifier. Catalogues from temporary deployments and targeted experiments make up a substantial share, reflecting both the diversity of New Zealand's regional seismology and the discoverability gap that a federated registry is intended to address. The inventory is offered as a working population of that registry and as a reproducible, citable reference for researchers, hazard modellers, and engineers.
\end{abstract}
\vspace{4pt}
{\small\noindent\textbf{Keywords:} earthquake catalogue; seismicity; temporary deployment; ocean-bottom seismometer; slow slip; FAIR data; Aotearoa New Zealand; GeoNet; metadata.}
\tableofcontents
\newpage
\section{Introduction and scope}\label{sec:intro}
An earthquake catalogue is the evidence base for seismicity research, seismic-hazard estimation, and engineering design. In Aotearoa New Zealand that evidence base is distributed across one authoritative national operational catalogue and many high-value but loosely coupled regional, temporary, and derived catalogues. Because these use different formats, identifiers, magnitude scales, and metadata conventions, combining them is manual and error-prone, a cost made explicit by national products such as the 2022 National Seismic Hazard Model, which required dedicated effort to standardise magnitudes across sources.

This report addresses the first practical step toward a federated capability: \emph{knowing which catalogues exist}. It inventories the earthquake catalogues developed for, or covering, New Zealand and its offshore margins (including the Hikurangi and Kermadec subduction systems and the subantarctic) whose coverage overlaps the period 1960--2026. Catalogues that begin earlier and extend into the period, and historical or derived catalogues used for New Zealand even where compiled recently, are included. Special attention is given to \textbf{temporary deployments}, which are scientifically rich yet routinely invisible to the national system.

Each catalogue is described using the \emph{Catalogue Submission Profile} (Appendix~\ref{app:schema}): a concise, controlled set of attributes covering identification and contact, coverage, content, classification, and availability. The profile records \emph{what} a catalogue contains and \emph{how} to obtain it, so that a reader can assess relevance before consulting the full documentation. Consistent with that scope, each entry here is a discovery record rather than a quality certification; deeper detail is carried by the catalogue's own dataset, DOI, or defining publication and is cited accordingly.
\section{Scope, sources, and compilation}\label{sec:methods}
This inventory was curated from peer-reviewed publications, dataset repositories (for example Zenodo, IRIS/EarthScope, and PANGAEA), FDSN network records, and the catalogues and reports of GeoNet / GNS Science and partner institutions. Each catalogue was profiled against the submission schema (Appendix~\ref{app:schema}) and reviewed against its defining publication or dataset; persistent identifiers and citations were checked against the original source, and fields that could not be established from the sources are left blank rather than inferred.

The inventory is a \emph{discovery} resource: each entry summarises what a catalogue contains and how to obtain it, and is not a certification of scientific adequacy; confirm details against the cited source before quantitative use. It is also not a census: earthquake-catalogue production in Aotearoa New Zealand is decentralised and continuous, so a residual long tail of specialised and thesis-derived catalogues remains (Section~\ref{sec:gaps}). Of the 81 catalogues listed, 73 carry a resolvable DOI or stable dataset identifier.
\section{The shape of New Zealand's catalogue ecosystem}\label{sec:overview}
New Zealand's earthquake-catalogue landscape has three tiers.

\textbf{(1) One authoritative national operational catalogue.} The GeoNet catalogue, produced by GNS Science / Earth Sciences New Zealand, provides continuous instrumental coverage (locations from the 1930s; descriptive events from $\sim$1460) and is underpinned by the permanent National Seismograph Network (FDSN code \texttt{NZ}). It is the parent dataset for nearly every derived New Zealand catalogue.

\textbf{(2) A large, scattered body of regional, temporary-deployment, volcano/geothermal, and special-purpose catalogues}, produced mostly by universities (Victoria University of Wellington, Otago, Auckland, Canterbury) and international partners. These deliver far higher \emph{local} precision (dense arrays, ocean-bottom and borehole sensors, and matched-filter / template-matching detection), but are typically published as one-off datasets discoverable only through the originating paper.

\textbf{(3) Derived products}: relocated and reprocessed catalogues, magnitude-homogenised and declustered hazard-model catalogues (NSHM 2010 and 2022), and physics-based synthetic (RSQSim) catalogues.

Catalogues produced by temporary deployments and targeted experiments make up a substantial part of this inventory (9 onshore and 4 offshore here, alongside much of the volcano-seismic and aftershock-response work), yet they are among the catalogues \emph{least} discoverable in the national system, and so are exactly the material a federated registry is designed to surface (their specific discoverability and heterogeneity issues are detailed in Section~\ref{sec:gaps}).
\section{Historical development, 1960--2026}\label{sec:timeline}
\begin{description}[leftmargin=0pt,style=unboxed,itemsep=3pt]
\item[1960s--1970s (analogue era).] The national catalogue is maintained by the DSIR Seismological Observatory; the \emph{Computer File of New Zealand Earthquakes} runs to 1974, while Eiby's \emph{Descriptive Catalogue} and \emph{Annotated List (1460--1965)} bridge the instrumental and historical records. Globally relocated coverage exists via ISC-GEM (from 1904) and later ISC-EHB (from 1964).
\item[1980s (digitisation).] Continuous digital waveforms are archived from September 1986. The 1987 Edgecumbe sequence is an early instrumented aftershock response.
\item[1990s (first large passive experiments).] SAPSE (Southern Alps Passive Seismic Experiment, 1995--96) is a landmark temporary deployment; Wellington-region relocation work begins; Downes' \emph{Atlas of Isoseismal Maps} (1995) systematises historical macroseismic intensities.
\item[2000s (targeted arrays and double-difference relocation).] A wave of temporary arrays (Marlborough \texttt{XB}, Tongariro, Taranaki, CNIPSE/NIGHT, SAHKE, ALFA08/09, and the SAMBA borehole array from 2008); double-difference relocation becomes standard; the 2003 Fiordland response.
\item[2010s (GeoNet modernises; sequences drive dense catalogues).] GeoNet moves to SeisComP with 3-D NonLinLoc (2012). The Canterbury sequence (Darfield 2010, Christchurch 2011) and Kaik\=oura 2016 drive dense aftershock and ocean-bottom catalogues (HOBITSS, 2014--15). Matched-filter / template-matching catalogues emerge, and slow-slip and tremor / low-frequency catalogues mature.
\item[2020s (machine-scale relocation, synthetics, multi-year geodetic catalogues).] High-precision matched-filter and nested-tomography catalogues (Taup\=o Volcanic Zone; Hikurangi); long RSQSim synthetic catalogues for NSHM 2022; multi-year GNSS slow-slip catalogues; dense onshore arrays (DWARFS, SALSA, SOSA); volcanic-unrest catalogues at Taup\=o (2019; 2022--23); and geothermal induced-seismicity monitoring (Rotokawa, Ngatamariki).
\end{description}
\section{Summary statistics}\label{sec:stats}
Table~\ref{tab:counts} gives the distribution of the 81 catalogues across categories; Table~\ref{tab:index} is a master index, numbered as in the detailed inventory (Section~\ref{sec:inventory}).

\begin{table}[h]
\centering
\caption{Catalogues by category.}\label{tab:counts}
\begin{tabular}{@{}L{9.5cm} r@{}}
\toprule
\textbf{Category} & \textbf{Count} \\
\midrule
National operational and network catalogues & 1 \\
Relocated and reprocessed catalogues & 11 \\
Historical and macroseismic (felt-report) catalogues & 6 \\
Aftershock-sequence rapid-response catalogues & 11 \\
Temporary onshore passive-seismic deployments & 9 \\
Temporary offshore / ocean-bottom (OBS) deployments & 4 \\
Volcano-seismic, geothermal and induced catalogues & 20 \\
Slow-slip, tremor, LFE and moment-tensor catalogues & 14 \\
Derived, homogenised and synthetic / hazard-model catalogues & 5 \\
\midrule
\textbf{Total} & \textbf{81} \\
\bottomrule
\end{tabular}
\end{table}

\begin{center}
{\small\textbf{Distribution by category}}\\[4pt]
\begin{tabular}{@{}L{6.6cm} l r@{}}
National operational and network & \textcolor{barblue}{\rule{0.45cm}{8pt}} & 1 \\
Relocated and reprocessed & \textcolor{barblue}{\rule{4.95cm}{8pt}} & 11 \\
Historical and macroseismic (felt-report) & \textcolor{barblue}{\rule{2.70cm}{8pt}} & 6 \\
Aftershock-sequence rapid-response & \textcolor{barblue}{\rule{4.95cm}{8pt}} & 11 \\
Temporary onshore passive-seismic deployments & \textcolor{barblue}{\rule{4.05cm}{8pt}} & 9 \\
Temporary offshore / ocean-bottom (OBS) deployments & \textcolor{barblue}{\rule{1.80cm}{8pt}} & 4 \\
Volcano-seismic, geothermal and induced & \textcolor{barblue}{\rule{9.00cm}{8pt}} & 20 \\
Slow-slip, tremor, LFE and moment-tensor & \textcolor{barblue}{\rule{6.30cm}{8pt}} & 14 \\
Derived, homogenised and synthetic / hazard-model & \textcolor{barblue}{\rule{2.25cm}{8pt}} & 5 \\
\end{tabular}
\end{center}

\paragraph{Availability.} 73 of the 81 catalogues carry a resolvable DOI or stable dataset identifier.

\begin{small}
\begin{longtable}{@{}r L{4.5cm} L{1.9cm} L{1.6cm} L{5.7cm}@{}}
\caption{Master index of catalogues (number matches the detailed entry in Section~\ref{sec:inventory}). The final column gives availability, the primary reference, and the DOI / dataset link where one exists.}\label{tab:index}\\
\toprule
\textbf{\#} & \textbf{Catalogue} & \textbf{Period} & \textbf{Type} & \textbf{Availability, ref \& DOI} \\
\midrule
\endfirsthead
\multicolumn{5}{l}{\small\itshape Table~\ref{tab:index} continued}\\
\toprule
\textbf{\#} & \textbf{Catalogue} & \textbf{Period} & \textbf{Type} & \textbf{Availability, ref \& DOI} \\
\midrule
\endhead
\bottomrule
\endfoot
1 & \hyperlink{c:1}{GeoNet Aotearoa New Zealand Earthquake Catalogue} & 1460\,--\,ongoing & background & open (DOI), \cite{GNSScience1970}, {\footnotesize \href{https://doi.org/10.21420/0S8P-TZ38}{\nolinkurl{10.21420/0S8P-TZ38}}} \\
2 & \hyperlink{c:2}{Eberhart-Phillips \& Reyners 2001-2011 NZ-Wide 3-D Relocated Catalogue} & 2001\,--\,2011 & background & open (DOI), \cite{Eberhart2022}, {\footnotesize \href{https://doi.org/10.1080/00288306.2022.2089171}{\nolinkurl{10.1080/00288306.2022.2089171}}} \\
3 & \hyperlink{c:3}{Hikurangi Subduction Zone Nested-Tomography Precise Relocation Catalogue (Aziz Zanjani et al. 2021)} & 2012\,--\,2016 & background & on request, \cite{AzizZanjani2021}, {\footnotesize \href{https://doi.org/10.1093/gji/ggab294}{\nolinkurl{10.1093/gji/ggab294}}} \\
4 & \hyperlink{c:4}{ISC-EHB Bulletin (New Zealand coverage)} & 1964\,--\,2021 & compilation & open (DOI), \cite{Weston2018}, {\footnotesize \href{https://doi.org/10.31905/PY08W6S3}{\nolinkurl{10.31905/PY08W6S3}}} \\
5 & \hyperlink{c:5}{ISC-GEM Global Instrumental Earthquake Catalogue (New Zealand coverage)} & 1904\,--\,2014 & large-events-only & open (DOI), \cite{DiGiacomo2018}, {\footnotesize \href{https://doi.org/10.31905/D808B825}{\nolinkurl{10.31905/D808B825}}} \\
6 & \hyperlink{c:6}{Taupo Volcanic Zone Consistent High-Precision Earthquake Catalogue (2007-2023)} & 2007\,--\,2023 & background & open (DOI), \cite{IllsleyKemp2025}, {\footnotesize \href{https://doi.org/10.26443/seismica.v4i1.1490}{\nolinkurl{10.26443/seismica.v4i1.1490}}} \\
7 & \hyperlink{c:7}{Wellington Region Double-Difference Relocated Catalogue (Du et al. 2004)} & 1990\,--\,2001 & background & not avail., \cite{Du2004}, {\footnotesize \href{https://doi.org/10.1111/j.1365-246X.2004.02366.x}{\nolinkurl{10.1111/j.1365-246X.2004.02366.x}}} \\
8 & \hyperlink{c:8}{Central Alpine Fault Relocated Microseismicity Catalogue (Michailos 2008-2017 / SAMBA)} & 2008\,--\,2017 & background & open (DOI), \cite{Michailos2019}, {\footnotesize \href{https://doi.org/10.5281/zenodo.3529755}{\nolinkurl{10.5281/zenodo.3529755}}} \\
9 & \hyperlink{c:9}{Central Southern Alps Matched-Filter Microseismicity Catalogue (Michailos 2009-2020)} & 2009\,--\,2020 & background & open (DOI), \cite{Michailos2025}, {\footnotesize \href{https://doi.org/10.5281/zenodo.15002912}{\nolinkurl{10.5281/zenodo.15002912}}} \\
10 & \hyperlink{c:10}{ISC Bulletin / Reviewed ISC Bulletin (New Zealand component)} & 1900\,--\,ongoing & compilation & open (DOI), \cite{InternationalSeismologicalCentre2026}, {\footnotesize \href{https://doi.org/10.31905/D808B830}{\nolinkurl{10.31905/D808B830}}} \\
11 & \hyperlink{c:11}{DWARFS Rupture-Limiting Section Boundary Relocated Microseismicity \& Focal-Mechanism Catalogue, Alpine Fault (Warren-Smith et al. 2022)} & 2019-04\,--\,2020-04 & background & open (DOI), \cite{WarrenSmith2022}, {\footnotesize \href{https://doi.org/10.5281/zenodo.6872209 (dataset); 10.1029/2022JB025219 (defining publication); 10.7914/SN/8N_2019 (FDSN 7S network)}{\nolinkurl{10.5281/zenodo.6872209 (dataset); 10.1029/2022JB025219 (defining publication); 10.7914/SN/8N_2019 (FDSN 7S network)}}} \\
12 & \hyperlink{c:12}{Consistent Earthquake Catalogue for New Zealand (Williams, Chamberlain \& Townend, 2001-2021)} & 2001-01\,--\,2021-01 & background & open (DOI), \cite{Williams2026}, {\footnotesize \href{https://doi.org/10.5281/zenodo.18820779}{\nolinkurl{10.5281/zenodo.18820779}}} \\
13 & \hyperlink{c:13}{Atlas of Isoseismal Maps of New Zealand Earthquakes (1st edition, Downes 1995)} & 1848\,--\,1990 & compilation & open (DOI), \cite{Downes1995}, {\footnotesize \href{https://www.naturalhazards.govt.nz/assets/Publications-Resources/7-Atlas-of-isoseismal-maps-of-New-Zealand-earthquakes-compressed-v2.pdf}{\nolinkurl{https://www.naturalhazards.govt.nz/assets/Publications-Resources/7-Atlas-of-isoseismal-maps-of-New-Zealand-earthquakes-compressed-v2.pdf}}} \\
14 & \hyperlink{c:14}{Atlas of Isoseismal Maps of New Zealand Earthquakes 1843-2003 (2nd edition, Downes \& Dowrick 2014)} & 1843\,--\,2003 & compilation & on request, \cite{Downes2014} \\
15 & \hyperlink{c:15}{A Descriptive Catalogue of New Zealand Earthquakes (Eiby 1968/1973)} & pre-1845 (Part I); 1846 (Part II)\,--\,1854 & historical & open (DOI), \cite{Eiby1968}, {\footnotesize \href{https://doi.org/10.1080/00288306.1968.10423671}{\nolinkurl{10.1080/00288306.1968.10423671}}} \\
16 & \hyperlink{c:16}{An Annotated List of New Zealand Earthquakes, 1460-1965 (Eiby 1968)} & 1460\,--\,1965 & compilation & open (DOI), \cite{Eiby1968a}, {\footnotesize \href{https://doi.org/10.1080/00288306.1968.10420275}{\nolinkurl{10.1080/00288306.1968.10420275}}} \\
17 & \hyperlink{c:17}{Magnitudes of New Zealand Earthquakes, 1901-1993 (Dowrick \& Rhoades 1998)} & 1901\,--\,1993 & large-events-only & open (DOI), \cite{Dowrick1998}, {\footnotesize \href{https://doi.org/10.5459/bnzsee.31.4.260-280}{\nolinkurl{10.5459/bnzsee.31.4.260-280}}} \\
18 & \hyperlink{c:18}{Pre-2010 Historical Seismicity near Christchurch (1869 Christchurch \& 1870 Lake Ellesmere)} & 1869\,--\,1870 & historical & open (DOI), \cite{Downes2012}, {\footnotesize \href{https://doi.org/10.1080/00288306.2012.690767}{\nolinkurl{10.1080/00288306.2012.690767}}} \\
19 & \hyperlink{c:19}{1987 Edgecumbe Earthquake Sequence (ML 6.3, Bay of Plenty)} & 1987-02\,--\,1987-03-18 & aftershock & on request, \cite{Smith1989}, {\footnotesize \href{https://doi.org/10.1080/00288306.1989.10421386}{\nolinkurl{10.1080/00288306.1989.10421386}}} \\
20 & \hyperlink{c:20}{2003 Mw 7.2 Fiordland (Secretary Island) Aftershock Relocations} & 2003-08\,--\,2003-10-31 & aftershock & on request, \cite{McGinty2007}, {\footnotesize \href{https://doi.org/10.1111/j.1365-246X.2007.03336.x}{\nolinkurl{10.1111/j.1365-246X.2007.03336.x}}} \\
21 & \hyperlink{c:21}{2010 Darfield (Canterbury) Mw 7.1 High-Resolution Aftershock Relocations} & 2010-09\,--\,2011 & aftershock & on request, \cite{Syracuse2013}, {\footnotesize \href{https://doi.org/10.1002/jgrb.50301}{\nolinkurl{10.1002/jgrb.50301}}} \\
22 & \hyperlink{c:22}{2011 Christchurch (Mw 6.2) Fine-Scale Aftershock Relocations} & 2011-02\,--\,? & aftershock & on request, \cite{Bannister2011}, {\footnotesize \href{https://doi.org/10.1785/gssrl.82.6.839}{\nolinkurl{10.1785/gssrl.82.6.839}}} \\
23 & \hyperlink{c:23}{Pegasus Bay Aftershock Sequence (2011-2012, Canterbury)} & 2011-12\,--\,2012 & aftershock & on request, \cite{Ristau2013}, {\footnotesize \href{https://doi.org/10.1093/gji/ggt222}{\nolinkurl{10.1093/gji/ggt222}}} \\
24 & \hyperlink{c:24}{2015 ML 6 Wanaka Aftershock Matched-Filter Catalogue} & 2015-05\,--\,2015-05 & aftershock & on request, \cite{WarrenSmith2017}, {\footnotesize \href{https://doi.org/10.1785/0220170016}{\nolinkurl{10.1785/0220170016}}} \\
25 & \hyperlink{c:25}{2016 Kaikoura (Mw 7.8) Aftershock Relocations (Lanza et al. 2019, STREWN + GeoNet)} & 2016-11\,--\,2017-05 & aftershock & on request, \cite{Lanza2019}, {\footnotesize \href{https://doi.org/10.1029/2019GL082780}{\nolinkurl{10.1029/2019GL082780}}} \\
26 & \hyperlink{c:26}{2016 Kaikoura 10-Year Matched-Filter Seismicity Catalogue (Chamberlain et al. 2021)} & 2009\,--\,2020 & background & on request, \cite{Chamberlain2021}, {\footnotesize \href{https://doi.org/10.1029/2021JB022304}{\nolinkurl{10.1029/2021JB022304}}} \\
27 & \hyperlink{c:27}{Kermadec 2021 Refined Earthquake Catalogue and Slip Models (Lythgoe et al.)} & 2021-03\,--\,2021 & aftershock & open (DOI), \cite{Lythgoe2023}, {\footnotesize \href{https://doi.org/10.7488/ds/7534}{\nolinkurl{10.7488/ds/7534}}} \\
28 & \hyperlink{c:28}{2016 Mw 7.1 Te Araroa (East Cape) Matched-Filter Foreshock \& Aftershock Catalogue (Warren-Smith et al. 2018)} & 2016-09\,--\,2016-11 & aftershock & open (DOI), \cite{WarrenSmith2018}, {\footnotesize \href{https://doi.org/10.1016/j.epsl.2017.11.020}{\nolinkurl{10.1016/j.epsl.2017.11.020}}} \\
29 & \hyperlink{c:29}{1994 Mw 6.7 Arthur's Pass Earthquake Finely Relocated Aftershock Catalogue (Bannister et al. 2006)} & 1994-06\,--\,? & aftershock & open (DOI), \cite{Bannister2006}, {\footnotesize \href{https://doi.org/10.1029/2006GL027462}{\nolinkurl{10.1029/2006GL027462}}} \\
30 & \hyperlink{c:30}{SAPSE, Southern Alps Passive Seismic Experiment (1995-96)} & 1995-11\,--\,1996-04 & background & open (DOI), \cite{Leitner2001}, {\footnotesize \href{https://doi.org/10.7914/SN/XC_1995}{\nolinkurl{10.7914/SN/XC_1995}}} \\
31 & \hyperlink{c:31}{DFDP-2 Real-Time / Matched-Filter Earthquake Catalogue, Whataroa (2014-2015)} & 2014-08\,--\,2015-01 & background & on request, \cite{Chamberlain2017} \\
32 & \hyperlink{c:32}{SOSA, Southland Otago Seismic Array Catalogue (2022-2023)} & 2022-10\,--\,2023-10 & background & open (DOI), \cite{Williams2025}, {\footnotesize \href{https://doi.org/10.7914/jr68-qq17}{\nolinkurl{10.7914/jr68-qq17}}} \\
33 & \hyperlink{c:33}{CNIPSE / NIGHT Passive, Central North Island Passive Seismic Experiment (FDSN XQ)} & 2001-01\,--\,2001-06 & background & open (DOI), \cite{Reyners2006}, {\footnotesize \href{https://doi.org/10.7914/SN/XQ_2001}{\nolinkurl{10.7914/SN/XQ_2001}}} \\
34 & \hyperlink{c:34}{Taranaki Dense Temporary Network Crustal Seismicity Catalogue (Sherburn \& White 2005)} & 2001-12\,--\,2002-09 & background & on request, \cite{Sherburn2005} \\
35 & \hyperlink{c:35}{Local Earthquake and Focal Mechanism Catalogue for the Northern Hikurangi Margin, December 2017-October 2018 (NZ3D onshore network)} & 2017-12\,--\,2018-10 & background & open (DOI), \cite{Woodward2026}, {\footnotesize \href{https://doi.org/10.5281/zenodo.17063327}{\nolinkurl{10.5281/zenodo.17063327}}} \\
36 & \hyperlink{c:36}{Southern Lakes / Northern Fiordland High-Resolution Microseismicity Catalogue (Central Otago Seismic Array, COSA)} & 2012-06\,--\,2013-03 & background & on request, \cite{WarrenSmith2017a} \\
37 & \hyperlink{c:37}{Dunedin Region Short-Term Broadband Array Microseismicity Catalogue (Akatore Fault)} & ?\,--\,? & background & on request, \cite{Todd2020} \\
38 & \hyperlink{c:38}{Hawke's Bay Region Microearthquake Catalogue (Bannister 1988; subducting Pacific plate fine structure)} & ?\,--\,? & background & not avail., \cite{Bannister1988} \\
39 & \hyperlink{c:39}{HOBITSS Microearthquake Catalogue (2014-2015, northern Hikurangi, manual)} & 2014-05\,--\,2015-06 & background & open (DOI), \cite{Yarce2019}, {\footnotesize \href{https://doi.org/10.5281/zenodo.2022405}{\nolinkurl{10.5281/zenodo.2022405}}} \\
40 & \hyperlink{c:40}{HOBITSS Automatic-Detection Earthquake Catalogue (2014 SSE, REST detector)} & 2014-09\,--\,2014-10 & background & open (DOI), \cite{Yarce2023}, {\footnotesize \href{https://doi.org/10.5281/zenodo.7274399}{\nolinkurl{10.5281/zenodo.7274399}}} \\
41 & \hyperlink{c:41}{Brothers Volcano Ocean-Bottom Hydrophone Seismicity Catalogue (south Kermadec arc)} & 2004\,--\,2005 & background & open (DOI), \cite{Dziak2008}, {\footnotesize \href{https://doi.org/10.1029/2007JB005533}{\nolinkurl{10.1029/2007JB005533}}} \\
42 & \hyperlink{c:42}{Dataset of Earthquakes in the Northern Hikurangi Subduction Zone During the 2019 Slow Slip Event (2018-2019 OBS deployment)} & 2018\,--\,2019 & background & open (DOI), \cite{Iwasaki2025}, {\footnotesize \href{https://doi.org/10.5281/zenodo.15713529}{\nolinkurl{10.5281/zenodo.15713529}}} \\
43 & \hyperlink{c:43}{Tongariro Volcanic Centre Local Earthquake Tomography Catalogue (Rowlands et al. 2005)} & 2001-01\,--\,2001-06 & background & on request, \cite{Rowlands2005} \\
44 & \hyperlink{c:44}{Okataina Volcanic Centre Double-Difference Relocated Catalogue (2010-2021)} & 2010\,--\,2021 & background & on request, \cite{Bannister2022}, {\footnotesize \href{https://doi.org/10.1016/j.jvolgeores.2022.107653}{\nolinkurl{10.1016/j.jvolgeores.2022.107653}}} \\
45 & \hyperlink{c:45}{Taupo Volcano Magmatic-Seismicity Catalogue (2019-2022 ECLIPSE temporary deployment)} & 2019\,--\,2022 & background & on request, \cite{Mestel2025}, {\footnotesize \href{https://doi.org/10.1029/2025JB031644}{\nolinkurl{10.1029/2025JB031644}}} \\
46 & \hyperlink{c:46}{Taupo 2019 Volcanic Unrest Earthquake Catalogue (Illsley-Kemp et al. 2021)} & 2019\,--\,2019 & background & on request, \cite{IllsleyKemp2021}, {\footnotesize \href{https://doi.org/10.1029/2021GC009803}{\nolinkurl{10.1029/2021GC009803}}} \\
47 & \hyperlink{c:47}{Taupo 2022-2023 Volcanic Unrest Seismicity Catalogue (Lamb et al. 2024)} & 2022-05\,--\,2023-05 & background & open (DOI), \cite{Lamb2024}, {\footnotesize \href{https://doi.org/10.26443/seismica.v3i2.1125}{\nolinkurl{10.26443/seismica.v3i2.1125}}} \\
48 & \hyperlink{c:48}{Cascading Earthquake Swarms, Northern Taupo Volcanic Zone (2019, Aber et al. 2025)} & 2019-03\,--\,2019-09 & background & on request, \cite{Aber2025}, {\footnotesize \href{https://doi.org/10.1029/2024GC012079}{\nolinkurl{10.1029/2024GC012079}}} \\
49 & \hyperlink{c:49}{Tarawera 2019 Dyke-Intrusion Earthquake Catalogue (Puhipuhi, Okataina; Benson et al. 2021)} & 2019-03\,--\,2019-03 & background & open (DOI), \cite{Benson2021}, {\footnotesize \href{https://doi.org/10.5281/zenodo.4035171}{\nolinkurl{10.5281/zenodo.4035171}}} \\
50 & \hyperlink{c:50}{Waimangu-Rotomahana-Mt Tarawera Geothermal Field Earthquake Swarm Catalogue (2004-2015)} & 2004\,--\,2015 & background & not avail., \cite{Bannister2016}, {\footnotesize \href{https://doi.org/10.1016/j.jvolgeores.2015.07.024}{\nolinkurl{10.1016/j.jvolgeores.2015.07.024}}} \\
51 & \hyperlink{c:51}{Mt Ruapehu Earthquake Swarm Catalogue (1990-2023)} & 1990\,--\,2023 & background & open (DOI), \cite{Mitchinson2024}, {\footnotesize \href{https://doi.org/10.3389/feart.2024.1343874}{\nolinkurl{10.3389/feart.2024.1343874}}} \\
52 & \hyperlink{c:52}{Whakaari/White Island Very-Long-Period (VLP) Earthquake Catalogue (2007-2019)} & 2007\,--\,2019 & background & on request, \cite{Park2020}, {\footnotesize \href{https://doi.org/10.1186/s40623-020-01224-z}{\nolinkurl{10.1186/s40623-020-01224-z}}} \\
53 & \hyperlink{c:53}{Ngatamariki Geothermal Injection-Induced Microseismicity Catalogue (2012-2015)} & 2012\,--\,2015 & induced & on request, \cite{Hopp2019}, {\footnotesize \href{https://doi.org/10.1029/2019GC008243}{\nolinkurl{10.1029/2019GC008243}}} \\
54 & \hyperlink{c:54}{Rotokawa Geothermal Field Induced Microseismicity Catalogue (2012-2015)} & 2012\,--\,2015 & induced & on request, \cite{Hopp2020}, {\footnotesize \href{https://doi.org/10.1016/j.geothermics.2019.101750}{\nolinkurl{10.1016/j.geothermics.2019.101750}}} \\
55 & \hyperlink{c:55}{Rotorua and Kawerau Geothermal Systems Relocated Seismicity Catalogue (1984-2004; Clarke et al. 2009)} & 1984\,--\,2004 & background & on request, \cite{Clarke2009}, {\footnotesize \href{https://doi.org/10.1016/j.jvolgeores.2008.11.004}{\nolinkurl{10.1016/j.jvolgeores.2008.11.004}}} \\
56 & \hyperlink{c:56}{Rotokawa Geothermal Field Microseismicity Catalogue 2008-2012 (Sherburn et al. 2015)} & 2008\,--\,2012 & induced & open (DOI), \cite{Sherburn2015}, {\footnotesize \href{https://doi.org/10.1016/j.geothermics.2014.11.001}{\nolinkurl{10.1016/j.geothermics.2014.11.001}}} \\
57 & \hyperlink{c:57}{Wairakei Geothermal Field Microseismicity Catalogue (Wairakei Seismic Network; Sepulveda et al. 2015)} & 2009-03\,--\,ongoing & induced & on request, \cite{Sepulveda2015} \\
58 & \hyperlink{c:58}{Mt Ruapehu 2022 Volcanic Unrest Drumbeat / Tremor Seismicity Catalogue (Bramwell et al. 2025)} & 2022-03\,--\,2022-07 & background & open (DOI), \cite{Bramwell2025}, {\footnotesize \href{https://doi.org/10.5281/zenodo.13910455}{\nolinkurl{10.5281/zenodo.13910455}}} \\
59 & \hyperlink{c:59}{Auckland Volcanic Field Matched-Filter Earthquake Catalogue (van Wijk et al. 2021, COVID-19 lockdown study)} & 2019-11\,--\,2020-06 & background & on request, \cite{vanWijk2021}, {\footnotesize \href{https://doi.org/10.5194/se-12-363-2021}{\nolinkurl{10.5194/se-12-363-2021}}} \\
60 & \hyperlink{c:60}{Monowai Submarine Volcano Long-Range Hydroacoustic (T-phase) Activity Catalogue} & 2003-07\,--\,2017-01 & induced & open (DOI), \cite{Metz2018}, {\footnotesize \href{https://doi.org/10.1029/2018JB015888}{\nolinkurl{10.1029/2018JB015888}}} \\
61 & \hyperlink{c:61}{Tongariro / Te Maari Relocated Microseismicity Catalogue 2010-2021 (Heise et al. 2024)} & 2010\,--\,2021 & background & on request, \cite{Heise2024}, {\footnotesize \href{https://doi.org/10.21420/0ky8-mx13}{\nolinkurl{10.21420/0ky8-mx13}}} \\
62 & \hyperlink{c:62}{Auckland Volcanic Field Self-Consistent Machine-Learning Earthquake Catalogue (Chamberlain et al. 2025; NHC report 114 'Sensing unrest in New Zealand's largest city')} & 2011\,--\,2022 & background & open (DOI), \cite{Chamberlain2025}, {\footnotesize \href{https://doi.org/10.5281/zenodo.15825886}{\nolinkurl{10.5281/zenodo.15825886}}} \\
63 & \hyperlink{c:63}{GeoNet New Zealand Earthquake Moment Tensor Solutions (MwNZ)} & 2003-08\,--\,ongoing & background & open (DOI), \cite{GNSScience2006}, {\footnotesize \href{https://doi.org/10.21420/MMJ9-CZ67}{\nolinkurl{10.21420/MMJ9-CZ67}}} \\
64 & \hyperlink{c:64}{Hikurangi Slow Slip Event Inventory (Wallace \& Beavan 2010)} & 2002\,--\,2010 & compilation & on request, \cite{Wallace2010}, {\footnotesize \href{https://doi.org/10.1029/2010JB007717}{\nolinkurl{10.1029/2010JB007717}}} \\
65 & \hyperlink{c:65}{Northern Hikurangi Tectonic Tremor Catalogue 2010-2015 (Todd \& Schwartz 2016)} & 2010\,--\,2015 & compilation & on request, \cite{Todd2016}, {\footnotesize \href{https://doi.org/10.1002/2016JB013480}{\nolinkurl{10.1002/2016JB013480}}} \\
66 & \hyperlink{c:66}{Ambient Tectonic Tremor Catalogue: Manawatu, Cape Turnagain, Marlborough \& Puysegur (Romanet \& Ide 2019)} & 2005\,--\,2016 & compilation & on request, \cite{Romanet2019}, {\footnotesize \href{https://doi.org/10.1186/s40623-019-1039-1}{\nolinkurl{10.1186/s40623-019-1039-1}}} \\
67 & \hyperlink{c:67}{Alpine Fault Low-Frequency Earthquake Catalogue (Chamberlain et al. 2014)} & 2009\,--\,2012 & background & on request, \cite{Chamberlain2014}, {\footnotesize \href{https://doi.org/10.1002/2014GC005436}{\nolinkurl{10.1002/2014GC005436}}} \\
68 & \hyperlink{c:68}{Alpine Fault Low-Frequency Earthquake Catalogue with Focal Mechanisms (Baratin et al. 2018)} & 2009-03\,--\,2016 & background & on request, \cite{Baratin2018}, {\footnotesize \href{https://doi.org/10.1016/j.epsl.2017.12.021}{\nolinkurl{10.1016/j.epsl.2017.12.021}}} \\
69 & \hyperlink{c:69}{Kaimanawa (Deep Hikurangi) Low-Frequency Earthquake Catalogue (Aden-Antoniow et al. 2024)} & 2010\,--\,2022 & background & on request, \cite{AdenAntoniow2024}, {\footnotesize \href{https://doi.org/10.1029/2023JB027971}{\nolinkurl{10.1029/2023JB027971}}} \\
70 & \hyperlink{c:70}{Raukumara Peninsula Repeating Earthquake Catalogue 2003-2020 (Hughes et al. 2021)} & 2003\,--\,2020 & background & on request, \cite{Hughes2021}, {\footnotesize \href{https://doi.org/10.1029/2021GC009670}{\nolinkurl{10.1029/2021GC009670}}} \\
71 & \hyperlink{c:71}{New Zealand-Wide Earthquake Focal Mechanism Catalogue for Tectonic Stress (Townend et al. 2012)} & 2004-01\,--\,2011-02 & compilation & on request, \cite{Townend2012}, {\footnotesize \href{https://doi.org/10.1016/j.epsl.2012.08.003}{\nolinkurl{10.1016/j.epsl.2012.08.003}}} \\
72 & \hyperlink{c:72}{Southeastern South Island / Otago Earthquake Focal Mechanism Catalogue (Williams et al. 2026)} & 2001\,--\,2020 & compilation & open (DOI), \cite{Williams2026a}, {\footnotesize \href{https://doi.org/10.5281/zenodo.17228270}{\nolinkurl{10.5281/zenodo.17228270}}} \\
73 & \hyperlink{c:73}{HOBITSS Ocean-Bottom Focal-Mechanism / Stress-Tensor Cycling Catalogue, Northern Hikurangi (Warren-Smith et al. 2019)} & 2014-05\,--\,2015-06 & background & on request, \cite{WarrenSmith2019}, {\footnotesize \href{https://doi.org/10.1038/s41561-019-0367-x}{\nolinkurl{10.1038/s41561-019-0367-x}}} \\
74 & \hyperlink{c:74}{Hikurangi Trench Earthquake-Swarm Catalogue (ETAS-based, Nishikawa et al. 2021)} & 1997\,--\,2015 & compilation & \cite{Nishikawa2021}, {\footnotesize \href{https://doi.org/10.1029/2020JB020618}{\nolinkurl{10.1029/2020JB020618}}} \\
75 & \hyperlink{c:75}{Manawatu Deep Short-Term Slow-Slip and Tremor Catalogue (Fasola et al. 2023)} & 2005\,--\,2016 & background & on request, \cite{Fasola2023}, {\footnotesize \href{https://doi.org/10.1029/2023GL105428}{\nolinkurl{10.1029/2023GL105428}}} \\
76 & \hyperlink{c:76}{Global Centroid Moment Tensor (GCMT) catalogue (New Zealand component)} & 1976\,--\,ongoing & compilation & open (DOI), \cite{Ekstrm2012}, {\footnotesize \href{https://doi.org/10.1016/j.pepi.2012.04.002}{\nolinkurl{10.1016/j.pepi.2012.04.002}}} \\
77 & \hyperlink{c:77}{NSHM 2010 Earthquake Catalogue (Stirling et al. 2012)} & 1840\,--\,2009 & background & open (DOI), \cite{Stirling2012}, {\footnotesize \href{https://doi.org/10.1785/0120110170}{\nolinkurl{10.1785/0120110170}}} \\
78 & \hyperlink{c:78}{Mw-Homogenised NZ Earthquake Catalogue for NSHM 2022 (Christophersen et al. 2022)} & 1931\,--\,2020 & compilation & open (DOI), \cite{Christophersen2022}, {\footnotesize \href{https://doi.org/10.21420/A2SN-XM76}{\nolinkurl{10.21420/A2SN-XM76}}} \\
79 & \hyperlink{c:79}{Integrated / Augmented Earthquake Catalogue for Aotearoa New Zealand (NSHM 2022; Rollins et al.)} & 1840\,--\,2020 & compilation & open (DOI), \cite{Rollins2022}, {\footnotesize \href{https://doi.org/10.21420/XT4Y-WY45}{\nolinkurl{10.21420/XT4Y-WY45}}} \\
80 & \hyperlink{c:80}{RSQSim Synthetic Earthquake Catalogue for New Zealand (NSHM2012 fault system, 276 kyr)} & synthetic (276,000-year simulated span; published 2021)\,--\,synthetic (276,000-year simulated span) & background & open (DOI), \cite{Shaw2021}, {\footnotesize \href{https://doi.org/10.5281/zenodo.5534462}{\nolinkurl{10.5281/zenodo.5534462}}} \\
81 & \hyperlink{c:81}{RSQSim Synthetic Earthquake Catalogue + Ground Motions for New Zealand (NZNSHM22 fault system)} & synthetic (RSQSim simulated span; published 2025)\,--\,synthetic (RSQSim simulated span) & background & open (DOI), \cite{Shaw2025}, {\footnotesize \href{https://doi.org/10.5281/zenodo.14532399}{\nolinkurl{10.5281/zenodo.14532399}}} \\
\end{longtable}
\end{small}
\section{Spatial coverage}\label{sec:spatial}
Figure~\ref{fig:coverage} maps where the inventory's catalogues sit. Panel~(a) shows the footprint of each regional catalogue (a representative bounding box) over the New Zealand coastline and active-fault network, coloured by type; the 18 national / NZ-wide catalogues are not drawn because they span the whole country, and two Kermadec-arc catalogues lie north of the map. Panel~(b) counts how many regional catalogues overlap each $0.1^\circ$ cell, which shows where catalogues concentrate, the Taup\=o Volcanic Zone, the Hikurangi margin, Canterbury, and the central Alpine Fault, while large offshore tracts, Fiordland, and the far south remain thinly catalogued. This spatial unevenness is the visual counterpart of the gaps discussed in Section~\ref{sec:gaps}.

\begin{sidewaysfigure}
\centering
\includegraphics[width=0.9\textheight]{figures/nz_catalogue_coverage.pdf}
\caption{Spatial coverage of the Aotearoa New Zealand earthquake-catalogue inventory. \textbf{(a)} Regional catalogue footprints (representative bounding boxes) coloured by category, over the coastline and active-fault network; national / NZ-wide and Kermadec-arc catalogues are handled in the annotations. \textbf{(b)} Number of regional catalogues overlapping each $0.1^\circ$ grid cell. Footprints are indicative extents for discovery, not precise data boundaries.}
\label{fig:coverage}
\end{sidewaysfigure}
\section{Coverage gaps and FAIR weaknesses}\label{sec:gaps}
\begin{itemize}[leftmargin=1.4em,itemsep=3pt]
\item \textbf{Not a census.} Because earthquake-catalogue production is decentralised and continuous, a residual long tail remains: likely-missing material includes further PhD-thesis catalogues, additional dedicated catalogues for New Zealand $M_w\!\geq\!6$ aftershock responses since 1960, more geothermal fields (Ohaaki, Mokai, Ng\=awh\=a), and the SISIE/MOANA/SHIRE offshore experiments not yet retained as distinct event catalogues.
\item \textbf{Discoverability is the core problem.} Temporary-deployment and regional catalogues live across journal supplements, Zenodo, IRIS/EarthScope, and FDSN network DOIs: the fragmentation a registry exists to fix.
\item \textbf{Heterogeneity.} Magnitude scales vary ($M_L$; $M_w$ via the GeoNet moment-tensor catalogue from 2003; the homogenised $M_{LNZ20}\!\rightarrow\!M_w$ used for NSHM 2022); the magnitude of completeness $M_c$ is rarely stated and is spatially variable; depths span 0--$\sim$600\,km for subduction seismicity, while many regional catalogues are shallow-crustal only.
\end{itemize}
\section{The inventory}\label{sec:inventory}
Each entry follows the submission-profile schema (Appendix~\ref{app:schema}). The italic header line gives \emph{catalogue type $\cdot$ focus $\cdot$ detection method $\cdot$ review status}. References are cited to the bibliography at the end of the report.
\subsection{National operational and network catalogues}
\textit{The authoritative national record and the permanent-network data products that underpin it.}\par\vspace{4pt}
\subsubsection*{\hypertarget{c:1}{1.}~ GeoNet Aotearoa New Zealand Earthquake Catalogue}
{\small\itshape background $\cdot$ geographical $\cdot$ automatic $\cdot$ partially reviewed}\par\vspace{2pt}
\begin{itemize}[leftmargin=1.3em,itemsep=0.6pt,topsep=1pt,parsep=0pt]
\kv{Also known as}{GeoNet earthquake catalogue; NZ Earthquake Catalogue; QuakeSearch; FDSN network NZ catalogue product; nz3drx 3-D NonLinLoc locations; NEID catalogue component}
\kv{Coverage}{1460 to ongoing}
\kv{Region}{Aotearoa New Zealand (national), incl. Hikurangi margin, Kermadec, Southern Alps, offshore regions}
\kv{Bounding box}{New Zealand and offshore margins, approx 34S-48S, 165E-180E/-175W (incl. Kermadec and Hikurangi margins and offshore regions)}
\kv{Depth range}{0-\textasciitilde{}600 km (subduction-zone seismicity to deep events)}
\kv{Magnitude}{Roughly M0 to M8+ instrumental; magnitude of completeness varies spatially and over time (network-dependent), Mc not stated on landing page}
\kv{Events}{\textasciitilde{}700,000+ (approx; \textasciitilde{}20,000 new events added per year)}
\kv{Content}{hypocentral locations; phase arrival-time picks; amplitude picks; quality information; historical observations}
\kv{Producer / contact}{GNS Science (Earth Sciences New Zealand) / GeoNet; Geohazards Information Manager (info@\allowbreak{}geonet.org.nz)}
\kv{Data source}{Continuous waveforms from the New Zealand National Seismograph Network (FDSN code NZ) plus historical felt/observational records; routinely processed by GeoNet (SeisComP since 2012). Successor to the DSIR Seismological Observatory / IGNS national catalogue lineage.}
\kv{Availability}{available (DOI) $\cdot$ DOI \href{https://doi.org/10.21420/0S8P-TZ38}{\nolinkurl{10.21420/0S8P-TZ38}}}
\end{itemize}
\par\vspace{2pt}The authoritative national operational earthquake catalogue for Aotearoa New Zealand, produced by GNS Science/GeoNet to provide a continuous public record of NZ seismicity for monitoring, research, and hazard assessment. It spans historical (descriptive) events from \textasciitilde{}1460 and instrumental locations from the 1930s to the present, and is the parent dataset for nearly every derived/relocated/hazard NZ catalogue. Location methods evolved over time (LOCAL+nz1d pre-1987; CUSP/GROPE+nz1dr \textasciitilde{}1987-2012; SeisComP from 2012 with NonLinLoc 'nz3drx' 3-D and LOCSAT/iasp91, LOCSAT routine since the New Generation Magnitude Calibration Dec 2018). The 3-D NonLinLoc 'nz3drx' is a location method within this single DOI, not a separate catalogue.
\par\vspace{2pt}\textbf{Key references:} \cite{GNSScience1970,WarrenSmith2024}.
\catentryrule
\subsection{Relocated and reprocessed catalogues}
\textit{Catalogues that re-pick, relocate, or re-time existing events with improved velocity models or double-difference / matched-filter methods, including the New Zealand components of global relocation efforts.}\par\vspace{4pt}
\subsubsection*{\hypertarget{c:2}{2.}~ Eberhart-Phillips \& Reyners 2001-2011 NZ-Wide 3-D Relocated Catalogue}
{\small\itshape background $\cdot$ geographical $\cdot$ manual $\cdot$ reviewed}\par\vspace{2pt}
\begin{itemize}[leftmargin=1.3em,itemsep=0.6pt,topsep=1pt,parsep=0pt]
\kv{Also known as}{NZ-Wide 3-D relocated catalogue 2001-2011}
\kv{Coverage}{2001 to 2011}
\kv{Region}{New Zealand (nationwide); southern South Island comparison}
\kv{Bounding box}{New Zealand (nationwide), approx. 34S-48S, 166E-179E, plus southern South Island comparison region}
\kv{Content}{hypocentral locations; relative relocations; quality information}
\kv{Producer / contact}{Donna Eberhart-Phillips \& Martin Reyners, GNS Science (Te Pū Ao)}
\kv{Data source}{GeoNet / New Zealand National Seismograph Network catalogue earthquakes (2001-2011), relocated using the NZ-Wide 3-D seismic velocity model (Eberhart-Phillips et al. NZ-Wide v2.2/v2.3); includes comparison with 2019-2020 auto-detected earthquakes in the southern South Island.}
\kv{Availability}{available (DOI) $\cdot$ DOI \href{https://doi.org/10.1080/00288306.2022.2089171}{\nolinkurl{10.1080/00288306.2022.2089171}}}
\end{itemize}
\par\vspace{2pt}Produced to improve hypocentral accuracy over routine 1-D GeoNet locations by relocating the 2001-2011 national earthquake catalogue within the NZ-Wide 3-D velocity model, giving locations consistent with the velocity structure for use in seismic-hazard and tectonic studies. A secondary aim was to assess detection/location performance in the sparsely instrumented southern South Island by comparing relocated events with 2019-2020 auto-detected seismicity. A derived child of the GeoNet catalogue.
\par\vspace{2pt}\textbf{Key references:} \cite{Eberhart2022,EberhartPhillips2022,EberhartPhillips2020}.
\catentryrule
\subsubsection*{\hypertarget{c:3}{3.}~ Hikurangi Subduction Zone Nested-Tomography Precise Relocation Catalogue (Aziz Zanjani et al. 2021)}
{\small\itshape background $\cdot$ geographical $\cdot$ manual $\cdot$ reviewed}\par\vspace{2pt}
\begin{itemize}[leftmargin=1.3em,itemsep=0.6pt,topsep=1pt,parsep=0pt]
\kv{Coverage}{2012 to 2016}
\kv{Region}{Hikurangi margin / Hikurangi subduction zone}
\kv{Bounding box}{Hikurangi subduction zone, North Island and northern South Island, New Zealand}
\kv{Depth range}{0-250}
\kv{Magnitude}{Approx. M2-5 (local/regional events used in relocation)}
\kv{Events}{57233 relocated (of 61654 catalogued); 43025 (75\%) further refined with cross-correlation}
\kv{Content}{hypocentral locations; relative relocations; phase arrival-time picks; other}
\kv{Producer / contact}{Farzaneh Aziz Zanjani (University of Miami) with Guoqing Lin (Univ. Miami) \& Clifford H. Thurber (Univ. Wisconsin-Madison)}
\kv{Data source}{GeoNet local/regional P-wave arrival times (2012-2016) combined with ISC-EHB teleseismic P-wave data (1964-2016); relocated with teletomoDD nested regional-global double-difference tomography plus waveform cross-correlation differential times.}
\kv{Availability}{available on request $\cdot$ DOI \href{https://doi.org/10.1093/gji/ggab294}{\nolinkurl{10.1093/gji/ggab294}}}
\end{itemize}
\par\vspace{2pt}Produced to obtain precise hypocentres and a nested regional-global velocity model along the Hikurangi subduction zone, sharpening imaging of the subducting slab and the double seismic zone. The teletomoDD inversion relocated the 2012-2016 GeoNet seismicity, with a large subset further refined using waveform cross-correlation to reduce relative location errors. A derived child of the GeoNet catalogue augmented by ISC-EHB teleseismic data.
\par\vspace{2pt}\textbf{Key references:} \cite{AzizZanjani2021}.
\catentryrule
\subsubsection*{\hypertarget{c:4}{4.}~ ISC-EHB Bulletin (New Zealand coverage)}
{\small\itshape compilation $\cdot$ geographical $\cdot$ manual $\cdot$ reviewed}\par\vspace{2pt}
\begin{itemize}[leftmargin=1.3em,itemsep=0.6pt,topsep=1pt,parsep=0pt]
\kv{Also known as}{EHB; Engdahl-van der Hilst-Buland Bulletin; ISC-EHB}
\kv{Coverage}{1964 to 2021}
\kv{Region}{New Zealand / Kermadec / Hikurangi subset of global coverage}
\kv{Bounding box}{Global; NZ/Kermadec/Hikurangi subset approx. 25S-55S, 160E-170W}
\kv{Magnitude}{Teleseismically well-constrained events (typically \textasciitilde{}M4.5+); magnitudes from ISC scheme}
\kv{Events}{\textasciitilde{}208,480 events globally (NZ/Kermadec/Hikurangi subset not separately reported)}
\kv{Content}{hypocentral locations; phase arrival-time picks; quality information; other}
\kv{Producer / contact}{International Seismological Centre (ISC), Thatcham, UK (E.R. Engdahl; D. Di Giacomo; D.A. Storchak)}
\kv{Data source}{Groomed subset of the ISC Bulletin: teleseismically well-recorded events relocated with the EHB algorithm using primary and depth-phase arrival times; NZ/Kermadec/Hikurangi events are the qualifying subset.}
\kv{Availability}{available (DOI) $\cdot$ DOI \href{https://doi.org/10.31905/PY08W6S3}{\nolinkurl{10.31905/PY08W6S3}}}
\end{itemize}
\par\vspace{2pt}A global, internally consistent 'groomed' dataset of teleseismically well-constrained earthquakes (1964-2021; the 1964-2008 portion rebuilt for this release) relocated with the EHB procedure to minimise location and especially depth bias. The qualifying New Zealand / Kermadec / Hikurangi subset provides homogeneous high-quality teleseismic hypocentres and depths suitable for slab and structure studies. It is a global product, not a standalone NZ catalogue; the NZ events are a regional subset.
\par{\small\itshape Related: One of three ISC global products with NZ coverage in this inventory (with ISC-GEM and the ISC Bulletin); ISC-EHB applies the EHB relocation procedure to well-recorded events.}
\par\vspace{2pt}\textbf{Key references:} \cite{Weston2018,Engdahl2020,InternationalSeismologicalCentre2022}.
\catentryrule
\subsubsection*{\hypertarget{c:5}{5.}~ ISC-GEM Global Instrumental Earthquake Catalogue (New Zealand coverage)}
{\small\itshape large-events-only $\cdot$ geographical $\cdot$ manual $\cdot$ reviewed}\par\vspace{2pt}
\begin{itemize}[leftmargin=1.3em,itemsep=0.6pt,topsep=1pt,parsep=0pt]
\kv{Also known as}{ISC-GEM}
\kv{Coverage}{1904 to 2014}
\kv{Region}{New Zealand / Kermadec / Hikurangi subset of global coverage}
\kv{Bounding box}{Global; NZ/Kermadec/Hikurangi subset approx. 25S-55S, 160E-170W}
\kv{Magnitude}{Large instrumental events: Mw\textgreater{}=7.5 (pre-1918), \textgreater{}=6.25 (1918-1959), \textgreater{}=5.5 (1960 onward); recomputed Mw}
\kv{Events}{NZ subset not separately reported (global catalogue, large-events-only)}
\kv{Content}{hypocentral locations; quality information; other}
\kv{Producer / contact}{International Seismological Centre (ISC) / Global Earthquake Model (GEM) Foundation (D. Di Giacomo; E.R. Engdahl; D.A. Storchak)}
\kv{Data source}{Reprocessed historical/instrumental station parametric data (1904 onward); hypocentres relocated and Mw recomputed under consistent methodology; NZ/Kermadec/Hikurangi large events are the qualifying subset.}
\kv{Availability}{available (DOI) $\cdot$ DOI \href{https://doi.org/10.31905/D808B825}{\nolinkurl{10.31905/D808B825}}}
\end{itemize}
\par\vspace{2pt}A homogeneous global catalogue of large instrumental earthquakes (1904-2014, extended toward 2018+) with relocated hypocentres and recomputed moment magnitudes (Mw), selected by time-dependent magnitude thresholds. The qualifying large New Zealand / Kermadec / Hikurangi events provide a long-period, large-event source/benchmark used in NZ seismic-hazard work (e.g. NSHM 2022). It is a global product, not NZ-specific; NZ events are a subset.
\par{\small\itshape Related: Companion ISC global products are ISC-EHB and the full ISC Bulletin; ISC-GEM is the magnitude-thresholded, relocated instrumental catalogue (1904-2014).}
\par\vspace{2pt}\textbf{Key references:} \cite{DiGiacomo2018,Storchak2015,InternationalSeismologicalCentre2021}.
\catentryrule
\subsubsection*{\hypertarget{c:6}{6.}~ Taupo Volcanic Zone Consistent High-Precision Earthquake Catalogue (2007-2023)}
{\small\itshape background $\cdot$ geographical $\cdot$ machine-learning $\cdot$ reviewed}\par\vspace{2pt}
\begin{itemize}[leftmargin=1.3em,itemsep=0.6pt,topsep=1pt,parsep=0pt]
\kv{Also known as}{TVZ catalogue; Illsley-Kemp \& Mestel 2025}
\kv{Coverage}{2007 to 2023}
\kv{Region}{Taupō Volcanic Zone (TVZ)}
\kv{Bounding box}{Taupō Volcanic Zone, central North Island, New Zealand (\textasciitilde{}300 km long zone)}
\kv{Events}{86579 events (\textasciitilde{}65\% high-precision relocated)}
\kv{Content}{hypocentral locations; relative relocations; phase arrival-time picks; quality information}
\kv{Producer / contact}{Finnigan Illsley-Kemp \& Eleanor Mestel, Te Herenga Waka - Victoria University of Wellington}
\kv{Data source}{GeoNet permanent network waveforms integrated with temporary TVZ deployments (incl. HADES array, Bannister 2009; back-arc rifting arrays); EQTransformer ML phase picking, 3-D velocity model location, GrowClust3D relative relocation.}
\kv{Availability}{available (DOI) $\cdot$ DOI \href{https://doi.org/10.26443/seismica.v4i1.1490}{\nolinkurl{10.26443/seismica.v4i1.1490}}}
\end{itemize}
\par\vspace{2pt}Produced to create a single internally consistent, high-precision earthquake catalogue spanning the whole \textasciitilde{}300 km Taupō Volcanic Zone over 2007-2023, overcoming the heterogeneity of routine locations across many separate deployments. Uses EQTransformer machine-learning phase picking, location in a 3-D velocity model, and GrowClust3D relative relocation. An ML + 3-D reprocessing effort merging permanent GeoNet data with temporary TVZ arrays.
\par\vspace{2pt}\textbf{Key references:} \cite{IllsleyKemp2025}.
\catentryrule
\subsubsection*{\hypertarget{c:7}{7.}~ Wellington Region Double-Difference Relocated Catalogue (Du et al. 2004)}
{\small\itshape background $\cdot$ geographical $\cdot$ manual $\cdot$ reviewed}\par\vspace{2pt}
\begin{itemize}[leftmargin=1.3em,itemsep=0.6pt,topsep=1pt,parsep=0pt]
\kv{Also known as}{Du et al. 2004 Wellington hypoDD catalogue}
\kv{Coverage}{1990 to 2001}
\kv{Region}{Wellington region / SW Hikurangi margin}
\kv{Bounding box}{Wellington region, southern North Island, New Zealand}
\kv{Magnitude}{Approx. M1.8-6.5}
\kv{Events}{6592 relocated of 6825 (96.5\%)}
\kv{Content}{hypocentral locations; relative relocations; phase arrival-time picks}
\kv{Producer / contact}{Wen-xuan Du \& Clifford H. Thurber (Univ. Wisconsin-Madison); with M. Reyners \& D. Eberhart-Phillips (GNS Science) and H. Zhang}
\kv{Data source}{Wellington-region seismicity 1990-2001 from the NZ National Seismograph Network plus some temporary stations (47 stations within 200 km; 22 primary control stations); catalogue and waveform-based differential times relocated with hypoDD double-difference.}
\kv{Availability}{not available $\cdot$ DOI \href{https://doi.org/10.1111/j.1365-246X.2004.02366.x}{\nolinkurl{10.1111/j.1365-246X.2004.02366.x}}}
\end{itemize}
\par\vspace{2pt}An early double-difference (hypoDD) relocation of Wellington-region earthquakes (1990-2001) to obtain precise relative hypocentres and resolve fine seismicity structure. The relocations sharpened both planes of the double seismic zone to \textless{}10 km thickness (\textasciitilde{}20 km apart) and revealed intraslab faulting at the SW end of the Hikurangi margin. A derived child of the national catalogue using permanent NZNSN plus some temporary stations.
\par\vspace{2pt}\textbf{Key references:} \cite{Du2004}.
\catentryrule
\subsubsection*{\hypertarget{c:8}{8.}~ Central Alpine Fault Relocated Microseismicity Catalogue (Michailos 2008-2017 / SAMBA)}
{\small\itshape background $\cdot$ geographical $\cdot$ manual $\cdot$ reviewed}\par\vspace{2pt}
\begin{itemize}[leftmargin=1.3em,itemsep=0.6pt,topsep=1pt,parsep=0pt]
\kv{Also known as}{SAMBA microseismicity catalogue; Michailos et al. 2019}
\kv{Coverage}{2008 to 2017}
\kv{Region}{Central Southern Alps / Alpine Fault (Ka Tiritiri o te Moana)}
\kv{Bounding box}{Central Southern Alps / Alpine Fault, South Island, New Zealand}
\kv{Magnitude}{ML -1.2 to 4.6; magnitude of completeness Mc = ML 1.1}
\kv{Events}{9,111 earthquakes (complete above ML 1.1)}
\kv{Content}{hypocentral locations; relative relocations; phase arrival-time picks; quality information}
\kv{Producer / contact}{Konstantinos Michailos with Euan G.C. Smith, Calum J. Chamberlain, Martha K. Savage \& John Townend, Victoria University of Wellington}
\kv{Data source}{GeoNet permanent stations plus temporary Southern Alps arrays (SAMBA, ALFA08/09, WIZARD, DFDP); local-earthquake locations relocated with hypoDD double-difference; QuakeML catalogue deposited on Zenodo.}
\kv{Availability}{available (DOI) $\cdot$ DOI \href{https://doi.org/10.5281/zenodo.3529755}{\nolinkurl{10.5281/zenodo.3529755}}}
\end{itemize}
\par\vspace{2pt}A decade-long (2008-2017) double-difference (hypoDD) relocated microseismicity catalogue of the central Southern Alps / Alpine Fault, produced to map along-strike variations in seismogenic thickness and the fault's microseismic behaviour. Built by combining permanent GeoNet data with multiple temporary Southern Alps deployments. Predecessor to the 2025 matched-filter catalogue.
\par{\small\itshape Related: The 2008-2017 double-difference catalogue; extended by the central Southern Alps matched-filter catalogue (Michailos et al. 2025, 2009-2020).}
\par\vspace{2pt}\textbf{Key references:} \cite{Michailos2019,Michailos2019a}.
\catentryrule
\subsubsection*{\hypertarget{c:9}{9.}~ Central Southern Alps Matched-Filter Microseismicity Catalogue (Michailos 2009-2020)}
{\small\itshape background $\cdot$ geographical $\cdot$ template-matching $\cdot$ reviewed}\par\vspace{2pt}
\begin{itemize}[leftmargin=1.3em,itemsep=0.6pt,topsep=1pt,parsep=0pt]
\kv{Also known as}{SAMBA matched-filter catalogue; Southern Alps/Ka Tiritiri o te Moana matched-filter catalog}
\kv{Coverage}{2009 to 2020}
\kv{Region}{Central Southern Alps / Alpine Fault (Ka Tiritiri o te Moana)}
\kv{Bounding box}{Central Southern Alps / Alpine Fault, South Island, New Zealand}
\kv{Magnitude}{Microseismicity; local magnitudes provided per event}
\kv{Content}{hypocentral locations; relative relocations; phase arrival-time picks; quality information; other}
\kv{Producer / contact}{Konstantinos Michailos with Calum J. Chamberlain, Simon C. Cox, John Townend et al. (Victoria University of Wellington / Australian National University)}
\kv{Data source}{Templates from the 2008-2017 SAMBA hypoDD catalogue applied via matched-filter (template-matching) detection to continuous GeoNet + Southern Alps array waveforms (2009-2020); QuakeML + CSV catalogue on Zenodo.}
\kv{Availability}{available (DOI) $\cdot$ DOI \href{https://doi.org/10.5281/zenodo.15002912}{\nolinkurl{10.5281/zenodo.15002912}}}
\end{itemize}
\par\vspace{2pt}A template-matching (matched-filter) enhanced microseismicity catalogue for the central Southern Alps / Alpine Fault (2009-2020) that substantially densifies the 2019 SAMBA catalogue with many additional small events. Produced to study temporal seismicity variations, including rainfall/snowmelt-triggered seismicity. A distinct child of the Michailos 2008-2017 catalogue rather than a duplicate.
\par{\small\itshape Related: Extends the 2008-2017 double-difference catalogue (Michailos et al. 2019) to 2009-2020 using template matching.}
\par\vspace{2pt}\textbf{Key references:} \cite{Michailos2025,Michailos2025a}.
\catentryrule
\subsubsection*{\hypertarget{c:10}{10.}~ ISC Bulletin / Reviewed ISC Bulletin (New Zealand component)}
{\small\itshape compilation $\cdot$ geographical $\cdot$ manual $\cdot$ reviewed}\par\vspace{2pt}
\begin{itemize}[leftmargin=1.3em,itemsep=0.6pt,topsep=1pt,parsep=0pt]
\kv{Also known as}{International Seismological Centre On-line Bulletin; ISC Reviewed Bulletin; reporting agency code ISC (NZ contributing agencies include GNS/WEL and historically DSIR)}
\kv{Coverage}{1900 to ongoing}
\kv{Region}{Global; New Zealand component covers onshore NZ, Hikurangi-Kermadec subduction margin, and subantarctic region}
\kv{Bounding box}{Global (NZ component covers approx. -50 to -33 lat, 165 to -176 lon, including Hikurangi/Kermadec margin and subantarctic region)}
\kv{Depth range}{0 to \textasciitilde{}700 km (full crustal-to-deep subduction range; global product)}
\kv{Magnitude}{Global product; Reviewed Bulletin applies manual review to all events with M\textgreater{}=3.5 (M = maximum reported network magnitude), and selectively for 2.5\textless{}=M\textless{}3.5. No single Mc defined for the NZ component (Mc varies with era and reporting completeness).}
\kv{Content}{hypocentral locations, phase arrival-time picks, amplitude picks, focal mechanisms, quality information, other}
\kv{Producer / contact}{International Seismological Centre (ISC), Thatcham, United Kingdom (Director: D.A. Storchak). NZ data contributed primarily by GNS Science / GeoNet (and historically DSIR Geophysics Division).}
\kv{Data source}{Parametric earthquake data (hypocentres, phase arrival-time picks, amplitude/magnitude readings, collected focal-mechanism solutions) reported to the ISC by \textasciitilde{}150 seismological agencies and networks worldwide, grouped and relocated by the ISC. The NZ component derives mainly from GNS Science / GeoNet (and historically DSIR) station and bulletin contributions, supplemented by globally collected phases.}
\kv{Availability}{available (DOI) $\cdot$ DOI \href{https://doi.org/10.31905/D808B830}{\nolinkurl{10.31905/D808B830}}}
\end{itemize}
\par\vspace{2pt}The ISC Bulletin is regarded as among the most complete long-term parametric records of global seismicity, compiling hypocentres, phase arrival times, amplitudes and magnitudes, and collected focal-mechanism solutions. The 'Reviewed' product applies manual analyst checking and ISC relocation to all events with M\textgreater{}=3.5 (and selectively for 2.5\textless{}=M\textless{}3.5), producing the authoritative reviewed location/magnitude record for moderate-to-large NZ events. It is distinct from the ISC-EHB and ISC-GEM derivative products and is widely cited for NZ regional/teleseismic locations.
\par{\small\itshape Related: The full ISC global bulletin; the relocated ISC-EHB and large-event ISC-GEM products are listed separately.}
\par\vspace{2pt}\textbf{Key references:} \cite{InternationalSeismologicalCentre2026,Bondr2011,Storchak2017,Storchak2020}.
\catentryrule
\subsubsection*{\hypertarget{c:11}{11.}~ DWARFS Rupture-Limiting Section Boundary Relocated Microseismicity \& Focal-Mechanism Catalogue, Alpine Fault (Warren-Smith et al. 2022)}
{\small\itshape background $\cdot$ particular-event-type $\cdot$ template-matching $\cdot$ reviewed}\par\vspace{2pt}
\begin{itemize}[leftmargin=1.3em,itemsep=0.6pt,topsep=1pt,parsep=0pt]
\kv{Also known as}{DWARFS (Dense Westland Arrays Researching Fault Segmentation); Earthquake Catalogues for DWARFS; DWARFS South / DWARFS North; FDSN 7S (2019-2021); network DOI identifier 8N (10.7914/SN/8N\_2019)}
\kv{Coverage}{2019-04 to 2020-04}
\kv{Region}{South Westland / central Alpine Fault, South Island, Aotearoa New Zealand}
\kv{Bounding box}{approx. -44.3 to -42.4 S, 168.7 to 171.5 E (DWARFS South near Martyr River, South Westland; DWARFS North near Inchbonnie / Hope-Alpine Fault intersection); not stated precisely}
\kv{Depth range}{crustal; upper \textasciitilde{}25 km (precise range not stated in record)}
\kv{Magnitude}{MLv -0.7 to 4.2}
\kv{Events}{\textasciitilde{}7,500 located earthquakes; \textasciitilde{}800 focal mechanisms (\textasciitilde{}1 year of data)}
\kv{Content}{hypocentral locations; relative relocations; phase arrival-time picks; focal mechanisms; quality information}
\kv{Producer / contact}{Emily Warren-Smith (GNS Science); contributors John Townend and Calum J. Chamberlain (Victoria University of Wellington / Te Herenga Waka)}
\kv{Data source}{Temporary dense broadband seismometer arrays (DWARFS South \textasciitilde{}10 stations near the South Westland-Central section boundary, Martyr River; DWARFS North \textasciitilde{}9 stations near the North Westland-Central boundary / Hope-Alpine Fault intersection at Inchbonnie), FDSN network 7S (2019-2021), operated by GNS Science; waveforms archived at IRIS/PASSCAL. Derived event catalogue archived on Zenodo.}
\kv{Availability}{available (DOI) $\cdot$ DOI \href{https://doi.org/10.5281/zenodo.6872209 (dataset); 10.1029/2022JB025219 (defining publication); 10.7914/SN/8N_2019 (FDSN 7S network)}{\nolinkurl{10.5281/zenodo.6872209 (dataset); 10.1029/2022JB025219 (defining publication); 10.7914/SN/8N_2019 (FDSN 7S network)}}}
\end{itemize}
\par\vspace{2pt}Relocated hypocentre catalogues and focal-mechanism solutions for microseismicity recorded by two dense temporary broadband sub-arrays straddling rupture-limiting section boundaries on the late-interseismic Alpine Fault, South Island. Produced to map subsurface fault geometry, structural heterogeneity, and stress state near section boundaries that may control earthquake rupture segmentation. Hypocentres derived primarily via HypoDD double-difference relocation with NonLinLoc origins where needed; all magnitudes are local magnitudes (MLv) from vertical-channel displacements. The catalogue is distinct from the listed Central Southern Alps / SAMBA Alpine Fault catalogues (different sub-array footprints, section boundaries, and a dedicated focal-mechanism set).
\par\vspace{2pt}\textbf{Key references:} \cite{WarrenSmith2022,WarrenSmith2022a,GNSScience2019}.
\catentryrule
\subsubsection*{\hypertarget{c:12}{12.}~ Consistent Earthquake Catalogue for New Zealand (Williams, Chamberlain \& Townend, 2001-2021)}
{\small\itshape background $\cdot$ geographical $\cdot$ machine-learning $\cdot$ partially reviewed}\par\vspace{2pt}
\begin{itemize}[leftmargin=1.3em,itemsep=0.6pt,topsep=1pt,parsep=0pt]
\kv{Also known as}{Consistent (Enhanced) Earthquake Catalogue for Aotearoa New Zealand 2001-2021; Williams et al. consistent NZ catalogue; Enhanced Earthquake Catalogue for Aotearoa New Zealand}
\kv{Coverage}{2001-01 to 2021-01}
\kv{Region}{Aotearoa New Zealand (nationwide, incl. offshore margins)}
\kv{Bounding box}{Aotearoa New Zealand (NZ-wide; approx. 165E-180E, 48S-33S, incl. offshore margins / Hikurangi / Kermadec) - not explicitly stated in dataset record}
\kv{Magnitude}{Local magnitudes (ML) computed consistently; explicit magnitude range / completeness magnitude (Mc) not stated in the dataset record}
\kv{Events}{407553}
\kv{Content}{hypocentral locations, relative relocations, phase arrival-time picks, amplitude picks, quality information (location uncertainties)}
\kv{Producer / contact}{Codee-Leigh Williams (lead/researcher), with Calum J. Chamberlain \& John Townend (supervisors), Te Herenga Waka-Victoria University of Wellington (School of Geography, Environment and Earth Sciences)}
\kv{Data source}{GeoNet seismic waveform archive (events reassessed from the GeoNet catalogue 2001-2021); automated re-picking and relocation workflow by VUW}
\kv{Availability}{available (DOI) $\cdot$ DOI \href{https://doi.org/10.5281/zenodo.18820779}{\nolinkurl{10.5281/zenodo.18820779}}}
\end{itemize}
\par\vspace{2pt}A modern, NZ-wide earthquake EVENT catalogue produced with a uniform automated workflow: EQTransformer deep-learning phase picking applied to GeoNet waveforms, location/relocation in the NZWide2.3 3-D velocity model (Eberhart-Phillips et al. 2022) with depth-fixing removed, and consistent local-magnitude calculation. It reassesses 407,553 earthquakes from the GeoNet catalogue spanning 1 Jan 2001 - 1 Jan 2021, providing freshly picked and located hypocentres, location uncertainties, local magnitudes, phase arrival-time picks and amplitude picks. It is a freshly re-picked/located catalogue rather than a magnitude-homogenised compilation of pre-existing solutions, making it distinct from NSHM 'consistent magnitudes' compilations and the GeoNet operational catalogue, and it is NZ-wide rather than a single-region relocation study.
\par\vspace{2pt}\textbf{Key references:} \cite{Williams2026,Williams2024,EberhartPhillips2022a}.
\catentryrule
\subsection{Historical and macroseismic (felt-report) catalogues}
\textit{Pre-instrumental and felt-report catalogues that extend the record through the nineteenth century and quantify shaking via Modified Mercalli / macroseismic intensity.}\par\vspace{4pt}
\subsubsection*{\hypertarget{c:13}{13.}~ Atlas of Isoseismal Maps of New Zealand Earthquakes (1st edition, Downes 1995)}
{\small\itshape compilation $\cdot$ geographical $\cdot$ felt-reports $\cdot$ reviewed}\par\vspace{2pt}
\begin{itemize}[leftmargin=1.3em,itemsep=0.6pt,topsep=1pt,parsep=0pt]
\kv{Also known as}{Downes 1995; IGNS Monograph 11; EQC Paper 7 (EQC 91/40)}
\kv{Coverage}{1848 to 1990}
\kv{Region}{New Zealand (national)}
\kv{Bounding box}{New Zealand mainland and near-offshore; one event south of the Kermadec Islands}
\kv{Magnitude}{Smallest event ML 3.5 (1963 Mangonui); largest the 1855 Wairarapa earthquake (\textasciitilde{}Mw 8.0-8.2); intensities in NZ-adapted Modified Mercalli (MM)}
\kv{Events}{123 earthquakes / 133 isoseismal maps}
\kv{Content}{historical observations; hypocentral locations; quality information; other}
\kv{Producer / contact}{G.L. Downes, Institute of Geological \& Nuclear Sciences (IGNS, now GNS Science); funded by EQC (Earthquake Commission, now Natural Hazards Commission Toka Tu Ake)}
\kv{Data source}{Compilation of historical felt/macroseismic observations, newspaper and documentary records, and prior catalogues (incl. Eiby) reassessed into NZ-adapted Modified Mercalli intensities; isoseismal maps drawn for each event.}
\kv{Availability}{available (DOI) $\cdot$ DOI \href{https://www.naturalhazards.govt.nz/assets/Publications-Resources/7-Atlas-of-isoseismal-maps-of-New-Zealand-earthquakes-compressed-v2.pdf}{\nolinkurl{https://www.naturalhazards.govt.nz/assets/Publications-Resources/7-Atlas-of-isoseismal-maps-of-New-Zealand-earthquakes-compressed-v2.pdf}}}
\end{itemize}
\par\vspace{2pt}Foundational national atlas of isoseismal (felt-intensity) maps for New Zealand earthquakes, produced to systematise macroseismic intensity data for hazard, engineering and historical seismicity studies using the NZ-adapted Modified Mercalli scale. Provides 133 isoseismal maps for the felt distributions of 123 earthquakes (122 in/near the NZ mainland plus one Mw\textasciitilde{}6.7 event south of the Kermadecs), each with revised hypocentral parameters and a brief descriptive account. Free PDF via the Natural Hazards Commission/EQC; superseded and extended by the 2014 2nd edition (Monograph 25).
\par\vspace{2pt}\textbf{Key references:} \cite{Downes1995}.
\catentryrule
\subsubsection*{\hypertarget{c:14}{14.}~ Atlas of Isoseismal Maps of New Zealand Earthquakes 1843-2003 (2nd edition, Downes \& Dowrick 2014)}
{\small\itshape compilation $\cdot$ geographical $\cdot$ felt-reports $\cdot$ reviewed}\par\vspace{2pt}
\begin{itemize}[leftmargin=1.3em,itemsep=0.6pt,topsep=1pt,parsep=0pt]
\kv{Also known as}{Downes \& Dowrick 2014; GNS Science Monograph 25; ISBN 978-0-478-19663-4}
\kv{Coverage}{1843 to 2003}
\kv{Region}{New Zealand (national)}
\kv{Bounding box}{New Zealand mainland and near-offshore margins}
\kv{Magnitude}{Spans small felt events through the largest historical NZ shocks (incl. 1855 Wairarapa \textasciitilde{}Mw 8.0-8.2); intensities in NZ-adapted Modified Mercalli (MM)}
\kv{Content}{historical observations; hypocentral locations; quality information; other}
\kv{Producer / contact}{G.L. Downes \& D.J. Dowrick, GNS Science, Lower Hutt}
\kv{Data source}{Revision and extension of the 1995 1st edition (Monograph 11), incorporating reassessed Modified Mercalli intensities, magnitudes and hypocentral locations from additional documentary and instrumental sources.}
\kv{Availability}{available on request}
\end{itemize}
\par\vspace{2pt}Revised and expanded national isoseismal atlas (797 pp.) covering NZ earthquakes 1843-2003, extending the 1995 1st edition with reassessed MM intensities, magnitudes and locations and additional events/maps. Produced to provide an updated authoritative macroseismic reference for seismic hazard, engineering and historical-seismicity work in New Zealand. A distinct successor edition (GNS Science Monograph 25), originally distributed on CD and available on request rather than as an open download.
\par\vspace{2pt}\textbf{Key references:} \cite{Downes2014}.
\catentryrule
\subsubsection*{\hypertarget{c:15}{15.}~ A Descriptive Catalogue of New Zealand Earthquakes (Eiby 1968/1973)}
{\small\itshape historical $\cdot$ geographical $\cdot$ felt-reports $\cdot$ reviewed}\par\vspace{2pt}
\begin{itemize}[leftmargin=1.3em,itemsep=0.6pt,topsep=1pt,parsep=0pt]
\kv{Also known as}{Eiby 1968 Part I; Eiby 1973 Part II}
\kv{Coverage}{pre-1845 (Part I); 1846 (Part II) to 1854}
\kv{Region}{New Zealand (national)}
\kv{Bounding box}{New Zealand (national)}
\kv{Magnitude}{No instrumental magnitudes; events assigned to four felt-based magnitude classes}
\kv{Events}{\textgreater{}80 (Part I, pre-1845); \textasciitilde{}300 further entries (Part II, 1846-1854)}
\kv{Content}{historical observations; quality information; other}
\kv{Producer / contact}{G.A. Eiby, Geophysics Division, DSIR (New Zealand)}
\kv{Data source}{Documentary/historical records: published and previously unpublished accounts, newspapers, settlers' and observers' reports of felt earthquakes; full source references given per event.}
\kv{Availability}{available (DOI) $\cdot$ DOI \href{https://doi.org/10.1080/00288306.1968.10423671}{\nolinkurl{10.1080/00288306.1968.10423671}}}
\end{itemize}
\par\vspace{2pt}Foundational documentary chronological catalogue of pre-instrumental New Zealand earthquakes compiled from felt/historical evidence, with events assigned to four felt-based magnitude classes and full source references. Part I (1968) lists \textgreater{}80 events felt before the end of 1845; Part II (1973) adds \textasciitilde{}300 further entries for shocks felt 1846-1854 (incl. correcting the mis-dated/mis-located 1848 Marlborough event) with isoseismal map(s). Underpins later Downes historical-macroseismic work and W.D. Smith's 1976 computer file.
\par\vspace{2pt}\textbf{Key references:} \cite{Eiby1968,Eiby1973}.
\catentryrule
\subsubsection*{\hypertarget{c:16}{16.}~ An Annotated List of New Zealand Earthquakes, 1460-1965 (Eiby 1968)}
{\small\itshape compilation $\cdot$ geographical $\cdot$ felt-reports $\cdot$ reviewed}\par\vspace{2pt}
\begin{itemize}[leftmargin=1.3em,itemsep=0.6pt,topsep=1pt,parsep=0pt]
\kv{Also known as}{Eiby 1968 annotated list}
\kv{Coverage}{1460 to 1965}
\kv{Region}{New Zealand (national)}
\kv{Bounding box}{New Zealand (national)}
\kv{Magnitude}{Mix of felt-based magnitude classes (early/pre-instrumental) and early instrumental magnitude estimates}
\kv{Content}{historical observations; hypocentral locations; quality information; other}
\kv{Producer / contact}{G.A. Eiby, Geophysics Division, DSIR (New Zealand)}
\kv{Data source}{Merged pre-instrumental documentary/felt records and early instrumental observatory data into a single long-span annotated chronological list.}
\kv{Availability}{available (DOI) $\cdot$ DOI \href{https://doi.org/10.1080/00288306.1968.10420275}{\nolinkurl{10.1080/00288306.1968.10420275}}}
\end{itemize}
\par\vspace{2pt}Long-span (1460-1965) annotated chronological list of New Zealand earthquakes that merges pre-instrumental documentary/felt records with early instrumental observatory determinations, serving as a companion to Eiby's descriptive catalogue and a precursor to W.D. Smith's 1976 computer file. Distinct in scope and span from the descriptive catalogue (which stops at 1854). Published in NZJGG 11(3):630-647.
\par\vspace{2pt}\textbf{Key references:} \cite{Eiby1968a}.
\catentryrule
\subsubsection*{\hypertarget{c:17}{17.}~ Magnitudes of New Zealand Earthquakes, 1901-1993 (Dowrick \& Rhoades 1998)}
{\small\itshape large-events-only $\cdot$ parameter-testing $\cdot$ other $\cdot$ reviewed}\par\vspace{2pt}
\begin{itemize}[leftmargin=1.3em,itemsep=0.6pt,topsep=1pt,parsep=0pt]
\kv{Also known as}{Dowrick \& Rhoades 1998 Ms/Mw catalogue}
\kv{Coverage}{1901 to 1993}
\kv{Region}{New Zealand (national)}
\kv{Bounding box}{New Zealand and near-offshore}
\kv{Magnitude}{202 larger earthquakes; Ms determinations with Mw and ML comparisons (Mw-from-Ms standard error \textasciitilde{}0.15)}
\kv{Events}{202}
\kv{Content}{hypocentral locations; quality information; other}
\kv{Producer / contact}{D.J. Dowrick \& D.A. Rhoades, GNS Science (Institute of Geological \& Nuclear Sciences)}
\kv{Data source}{Reassessment of surface-wave (Ms) and other magnitudes for larger NZ earthquakes from observatory/bulletin amplitude data, with revision of earlier observatory magnitudes and locations; conversions to Mw.}
\kv{Availability}{available (DOI) $\cdot$ DOI \href{https://doi.org/10.5459/bnzsee.31.4.260-280}{\nolinkurl{10.5459/bnzsee.31.4.260-280}}}
\end{itemize}
\par\vspace{2pt}Reassessment catalogue of magnitudes for 202 larger New Zealand earthquakes over 1901-1993, deriving surface-wave magnitudes (Ms) and conversions to moment magnitude (Mw) and ML, and revising earlier observatory magnitudes and locations. Provides a consistent magnitude backbone for NZ historical seismicity alongside the Downes isoseismal atlas. Published in the Bulletin of the NZSEE.
\par\vspace{2pt}\textbf{Key references:} \cite{Dowrick1998}.
\catentryrule
\subsubsection*{\hypertarget{c:18}{18.}~ Pre-2010 Historical Seismicity near Christchurch (1869 Christchurch \& 1870 Lake Ellesmere)}
{\small\itshape historical $\cdot$ geographical $\cdot$ felt-reports $\cdot$ reviewed}\par\vspace{2pt}
\begin{itemize}[leftmargin=1.3em,itemsep=0.6pt,topsep=1pt,parsep=0pt]
\kv{Also known as}{Downes \& Yetton 2012 Canterbury historical reassessment}
\kv{Coverage}{1869 to 1870}
\kv{Region}{Canterbury}
\kv{Bounding box}{Canterbury, Christchurch and Lake Ellesmere (Te Waihora) area, South Island NZ}
\kv{Magnitude}{1869 Christchurch Mw 4.7-4.9; 1870 Lake Ellesmere Mw 5.6-5.8 (estimated from MM intensities)}
\kv{Events}{2}
\kv{Content}{historical observations; hypocentral locations; quality information; other}
\kv{Producer / contact}{G. Downes (GNS Science) \& M. Yetton (Geotech Consulting)}
\kv{Data source}{Documentary felt/MM evidence (newspaper and historical accounts) reassessed to relocate and re-magnitude two 19th-century Canterbury earthquakes, produced after the 2010-2011 Canterbury sequence.}
\kv{Availability}{available (DOI) $\cdot$ DOI \href{https://doi.org/10.1080/00288306.2012.690767}{\nolinkurl{10.1080/00288306.2012.690767}}}
\end{itemize}
\par\vspace{2pt}Focused regional historical reassessment relocating and re-magnituding the 1869 Christchurch (Mw 4.7-4.9) and 1870 Lake Ellesmere (Mw 5.6-5.8) earthquakes from documentary felt/MM evidence, undertaken after the 2010-2011 Canterbury earthquake sequence to clarify pre-2010 Canterbury seismicity. A focused two-event historical study rather than a broad catalogue. Published in NZJGG 55(3):199-205.
\par\vspace{2pt}\textbf{Key references:} \cite{Downes2012}.
\catentryrule
\subsection{Aftershock-sequence rapid-response catalogues}
\textit{Sequence-specific catalogues, often built from rapid-response temporary instrumentation deployed around a significant mainshock.}\par\vspace{4pt}
\subsubsection*{\hypertarget{c:19}{19.}~ 1987 Edgecumbe Earthquake Sequence (ML 6.3, Bay of Plenty)}
{\small\itshape aftershock $\cdot$ geographical $\cdot$ manual $\cdot$ reviewed}\par\vspace{2pt}
\begin{itemize}[leftmargin=1.3em,itemsep=0.6pt,topsep=1pt,parsep=0pt]
\kv{Also known as}{Edgecumbe aftershock sequence; 1987 Edgecumbe earthquake}
\kv{Coverage}{1987-02 to 1987-03-18}
\kv{Region}{Bay of Plenty (Rangitaiki Plains / Whakatane Graben)}
\kv{Bounding box}{Rangitaiki Plains / Whakatane Graben, Bay of Plenty, North Island, New Zealand}
\kv{Depth range}{\textasciitilde{}4-8}
\kv{Magnitude}{Mainshock ML \textasciitilde{}6.3 (Mw \textasciitilde{}6.5); largest aftershock \textasciitilde{}ML 5.6; sequence reported to include hundreds of shocks ML\textgreater{}=3.0}
\kv{Events}{\textgreater{}600 shocks ML\textgreater{}=3.0 (\textasciitilde{}130 before the main shock)}
\kv{Content}{hypocentral locations; phase arrival-time picks; fault-geometry}
\kv{Producer / contact}{Euan G. C. Smith \& Clive M. M. Oppenheimer (Victoria University of Wellington) / DSIR Geophysics Division (now GNS Science)}
\kv{Data source}{New Zealand National Seismograph Network (NZNSN) plus rapidly deployed portable seismographs on the Rangitaiki Plains}
\kv{Availability}{available on request $\cdot$ DOI \href{https://doi.org/10.1080/00288306.1989.10421386}{\nolinkurl{10.1080/00288306.1989.10421386}}}
\end{itemize}
\par\vspace{2pt}Foreshock-mainshock-aftershock catalogue of the 2 March 1987 Edgecumbe earthquake (NZDT; 1 March UTC), which produced surface rupture on the Edgecumbe Fault across the Rangitaiki Plains. Produced to characterise the spatial and temporal evolution of the sequence and the geometry of the causative faulting in the Whakatane Graben. The defining paper is Smith \& Oppenheimer (1989), covering 1987 February 21 to March 18.
\par\vspace{2pt}\textbf{Key references:} \cite{Smith1989}.
\catentryrule
\subsubsection*{\hypertarget{c:20}{20.}~ 2003 Mw 7.2 Fiordland (Secretary Island) Aftershock Relocations}
{\small\itshape aftershock $\cdot$ geographical $\cdot$ manual $\cdot$ reviewed}\par\vspace{2pt}
\begin{itemize}[leftmargin=1.3em,itemsep=0.6pt,topsep=1pt,parsep=0pt]
\kv{Also known as}{2003 Fiordland earthquake; Secretary Island / Te Anau 2003}
\kv{Coverage}{2003-08 to 2003-10-31}
\kv{Region}{Fiordland (subduction interface; overlying Alpine Fault)}
\kv{Bounding box}{Fiordland / Secretary Island region, southwest South Island, New Zealand}
\kv{Magnitude}{Aftershocks ML 1.0-6.1; mainshock ML 7.0 / Mw 7.2}
\kv{Events}{\textasciitilde{}8,025 recorded; 1,771 relocated (377 best-quality)}
\kv{Content}{hypocentral locations; relative relocations; phase arrival-time picks; fault-geometry}
\kv{Producer / contact}{Peter McGinty \& Russell Robinson (GNS Science)}
\kv{Data source}{NZNSN plus \textasciitilde{}10 helicopter-deployed portable seismographs deployed within 3 days of the mainshock}
\kv{Availability}{available on request $\cdot$ DOI \href{https://doi.org/10.1111/j.1365-246X.2007.03336.x}{\nolinkurl{10.1111/j.1365-246X.2007.03336.x}}}
\end{itemize}
\par\vspace{2pt}Double-difference relocated aftershock catalogue of the 21 August 2003 Mw 7.2 (ML 7.0) Fiordland subduction-interface earthquake. Produced to map the aftershock distribution, constrain the main shock fault plane, and assess static stress changes on the overlying Alpine Fault. \textasciitilde{}8,025 aftershocks recorded (ML 1.0-6.1), 1,771 relocated by double-difference, 377 best-quality.
\par\vspace{2pt}\textbf{Key references:} \cite{McGinty2007}.
\catentryrule
\subsubsection*{\hypertarget{c:21}{21.}~ 2010 Darfield (Canterbury) Mw 7.1 High-Resolution Aftershock Relocations}
{\small\itshape aftershock $\cdot$ geographical $\cdot$ manual $\cdot$ reviewed}\par\vspace{2pt}
\begin{itemize}[leftmargin=1.3em,itemsep=0.6pt,topsep=1pt,parsep=0pt]
\kv{Also known as}{Darfield aftershock relocations; Syracuse et al. 2013; D-RAD / FDSN 4A deployment}
\kv{Coverage}{2010-09 to 2011}
\kv{Region}{Canterbury (Greendale Fault / Darfield)}
\kv{Bounding box}{Canterbury Plains west of Christchurch, South Island, New Zealand}
\kv{Magnitude}{Mainshock Mw 7.1}
\kv{Content}{hypocentral locations; relative relocations; phase arrival-time picks; fault-geometry}
\kv{Producer / contact}{Ellen M. Syracuse \& Clifford H. Thurber (University of Wisconsin-Madison); co-authors C. J. Rawles, M. K. Savage (VUW), S. Bannister (GNS Science)}
\kv{Data source}{Rapid-response 10-station PASSCAL/IRIS RAMP array (Darfield RAMP Aftershock Deployment, FDSN 4A\_2010, DOI 10.7914/SN/4A\_2010, operated by University of Wisconsin-Madison) plus temporary GNS and GeoNet stations}
\kv{Availability}{available on request $\cdot$ DOI \href{https://doi.org/10.1002/jgrb.50301}{\nolinkurl{10.1002/jgrb.50301}}}
\end{itemize}
\par\vspace{2pt}High-resolution double-difference tomography relocation of Darfield aftershocks following the 4 September 2010 Mw 7.1 earthquake. Produced to delineate buried/blind faults (including the Greendale Fault) and derive a 3-D velocity model, clarifying fault activity in the Canterbury sequence. Distinct from the later Christchurch and Pegasus Bay sub-sequence relocations.
\par\vspace{2pt}\textbf{Key references:} \cite{Syracuse2013,DarfieldRAMPAftershockDeploymentDRAD2010}.
\catentryrule
\subsubsection*{\hypertarget{c:22}{22.}~ 2011 Christchurch (Mw 6.2) Fine-Scale Aftershock Relocations}
{\small\itshape aftershock $\cdot$ geographical $\cdot$ manual $\cdot$ reviewed}\par\vspace{2pt}
\begin{itemize}[leftmargin=1.3em,itemsep=0.6pt,topsep=1pt,parsep=0pt]
\kv{Also known as}{Christchurch 22 February 2011 aftershock relocations; Bannister et al. 2011}
\kv{Coverage}{2011-02 to unknown}
\kv{Region}{Canterbury (Christchurch / Port Hills fault)}
\kv{Bounding box}{Christchurch / Port Hills, Canterbury, South Island, New Zealand}
\kv{Magnitude}{Mainshock Mw 6.2}
\kv{Content}{hypocentral locations; relative relocations; phase arrival-time picks; fault-geometry}
\kv{Producer / contact}{Stephen Bannister, Bill Fry, Martin Reyners, John Ristau, Haijiang Zhang (GNS Science; Zhang at MIT/USTC)}
\kv{Data source}{Dense GeoNet permanent network plus CanNet and temporary portable recorders around Christchurch}
\kv{Availability}{available on request $\cdot$ DOI \href{https://doi.org/10.1785/gssrl.82.6.839}{\nolinkurl{10.1785/gssrl.82.6.839}}}
\end{itemize}
\par\vspace{2pt}Fine-scale double-difference tomography relocation of aftershocks of the 22 February 2011 Mw 6.2 Christchurch earthquake. Produced to resolve the near-surface (Port Hills) fault responsible for the high ground motions that caused widespread damage and fatalities in Christchurch. Part of the broader Canterbury earthquake sequence work, distinct from the Darfield (Syracuse 2013) and Pegasus Bay (Ristau 2013) relocations.
\par\vspace{2pt}\textbf{Key references:} \cite{Bannister2011}.
\catentryrule
\subsubsection*{\hypertarget{c:23}{23.}~ Pegasus Bay Aftershock Sequence (2011-2012, Canterbury)}
{\small\itshape aftershock $\cdot$ geographical $\cdot$ manual $\cdot$ reviewed}\par\vspace{2pt}
\begin{itemize}[leftmargin=1.3em,itemsep=0.6pt,topsep=1pt,parsep=0pt]
\kv{Also known as}{Pegasus Bay sequence; offshore Christchurch 2011-2012 aftershocks; Ristau et al. 2013}
\kv{Coverage}{2011-12 to 2012}
\kv{Region}{Canterbury (offshore Pegasus Bay)}
\kv{Bounding box}{Pegasus Bay, offshore east of Christchurch, Canterbury, New Zealand}
\kv{Magnitude}{Mainshock (parent) Mw 7.1 Darfield; Pegasus Bay sub-sequence aftershock magnitude range not verified}
\kv{Events}{1,075 events located; 53 regional moment tensors}
\kv{Content}{hypocentral locations; relative relocations; focal mechanisms; phase arrival-time picks}
\kv{Producer / contact}{John Ristau, Caroline Holden, Anna Kaiser, Charles Williams, Stephen Bannister, Bill Fry (GNS Science)}
\kv{Data source}{GeoNet network plus temporary low-cost accelerometers; regional moment tensor analysis}
\kv{Availability}{available on request $\cdot$ DOI \href{https://doi.org/10.1093/gji/ggt222}{\nolinkurl{10.1093/gji/ggt222}}}
\end{itemize}
\par\vspace{2pt}Relocated double-difference tomography catalogue of the offshore Pegasus Bay sub-sequence of the Mw 7.1 Darfield (Canterbury) earthquake, which began \textasciitilde{}December 2011 east of Christchurch in Pegasus Bay. Produced to locate the offshore aftershock activity and characterise its faulting style; 1,075 events located and 53 regional moment tensors derived (predominantly reverse mechanisms).
\par\vspace{2pt}\textbf{Key references:} \cite{Ristau2013}.
\catentryrule
\subsubsection*{\hypertarget{c:24}{24.}~ 2015 ML 6 Wanaka Aftershock Matched-Filter Catalogue}
{\small\itshape aftershock $\cdot$ geographical $\cdot$ template-matching $\cdot$ reviewed}\par\vspace{2pt}
\begin{itemize}[leftmargin=1.3em,itemsep=0.6pt,topsep=1pt,parsep=0pt]
\kv{Also known as}{4 May 2015 Wanaka earthquake aftershock catalogue; Warren-Smith et al. 2017}
\kv{Coverage}{2015-05 to 2015-05}
\kv{Region}{Central Otago / Southern Alps (Wanaka)}
\kv{Bounding box}{Wanaka / Central Otago, Southern Alps, South Island, New Zealand}
\kv{Magnitude}{Mainshock ML 6.0; aftershock catalogue extends below GeoNet Mc via matched-filter detection}
\kv{Events}{2,544 aftershocks detected (over \textasciitilde{}26 days)}
\kv{Content}{hypocentral locations; relative relocations; phase arrival-time picks; fault-geometry}
\kv{Producer / contact}{Emily Warren-Smith, Calum J. Chamberlain, Simon Lamb, John Townend (Victoria University of Wellington; Warren-Smith later GNS Science)}
\kv{Data source}{Temporary broadband network plus GeoNet national network; matched-filter (template-matching) detection using mainshock plus aftershock templates}
\kv{Availability}{available on request $\cdot$ DOI \href{https://doi.org/10.1785/0220170016}{\nolinkurl{10.1785/0220170016}}}
\end{itemize}
\par\vspace{2pt}High-precision matched-filter detected and relocated aftershock catalogue of the 4 May 2015 ML 6.0 Wanaka earthquake, the largest shallow Central Otago / Southern Alps event in the record. Produced to characterise the aftershock sequence and identify a NW-dipping blind fault using template-matching detection and precise relative relocation. Detected 2,544 aftershocks over the \textasciitilde{}26 days following the mainshock.
\par\vspace{2pt}\textbf{Key references:} \cite{WarrenSmith2017}.
\catentryrule
\subsubsection*{\hypertarget{c:25}{25.}~ 2016 Kaikoura (Mw 7.8) Aftershock Relocations (Lanza et al. 2019, STREWN + GeoNet)}
{\small\itshape aftershock $\cdot$ geographical $\cdot$ manual $\cdot$ reviewed}\par\vspace{2pt}
\begin{itemize}[leftmargin=1.3em,itemsep=0.6pt,topsep=1pt,parsep=0pt]
\kv{Also known as}{Kaikoura aftershock relocations; Lanza et al. 2019; STREWN FDSN Z1 (2016-2017)}
\kv{Coverage}{2016-11 to 2017-05}
\kv{Region}{North Canterbury / Marlborough (Kaikoura rupture zone)}
\kv{Bounding box}{Northeast South Island / Marlborough to Kaikoura / Cook Strait, New Zealand}
\kv{Magnitude}{\textasciitilde{}2,700 events ML\textgreater{}=3; mainshock Mw 7.8}
\kv{Events}{\textasciitilde{}2,700 (ML\textgreater{}=3)}
\kv{Content}{hypocentral locations; relative relocations; phase arrival-time picks; fault-geometry}
\kv{Producer / contact}{Federica Lanza (lead; later ETH Zurich) with C. J. Chamberlain, K. Jacobs, E. Warren-Smith, H. J. Godfrey, M. Kortink, C. H. Thurber, M. K. Savage, J. Townend, S. Roecker, D. Eberhart-Phillips}
\kv{Data source}{Rapid-response STREWN temporary network (Seismic Triggering Response for Earthquakes around Wellington NZ; FDSN Z1\_2016, DOI 10.7914/SN/Z1\_2016, University of Wisconsin-Madison) plus GeoNet}
\kv{Availability}{available on request $\cdot$ DOI \href{https://doi.org/10.1029/2019GL082780}{\nolinkurl{10.1029/2019GL082780}}}
\end{itemize}
\par\vspace{2pt}Relocated aftershock catalogue of the 14 November 2016 Mw 7.8 Kaikoura earthquake (\textasciitilde{}2,700 ML\textgreater{}=3 events, Nov 2016-May 2017), repicked and relocated in the Eberhart-Phillips NZ-Wide 3-D velocity model. Produced to constrain crustal fault connectivity across the \textgreater{}20 faults that ruptured in the Kaikoura earthquake. Distinct from the longer Chamberlain et al. 2021 10-year matched-filter catalogue.
\par\vspace{2pt}\textbf{Key references:} \cite{Lanza2019,STREWNSeismicTriggeringResponseforEarthquakesaroundWellingtonNZ2016}.
\catentryrule
\subsubsection*{\hypertarget{c:26}{26.}~ 2016 Kaikoura 10-Year Matched-Filter Seismicity Catalogue (Chamberlain et al. 2021)}
{\small\itshape background $\cdot$ geographical $\cdot$ template-matching $\cdot$ reviewed}\par\vspace{2pt}
\begin{itemize}[leftmargin=1.3em,itemsep=0.6pt,topsep=1pt,parsep=0pt]
\kv{Also known as}{Chamberlain et al. 2021 Kaikoura catalogue; pre/co/post-seismic Kaikoura catalogue}
\kv{Coverage}{2009 to 2020}
\kv{Region}{North Canterbury / Marlborough (Kaikoura rupture faults)}
\kv{Bounding box}{Northeast South Island / Marlborough to Kaikoura rupture zone, New Zealand}
\kv{Magnitude}{Extends well below GeoNet Mc via matched-filter detection; mainshock Mw 7.8}
\kv{Events}{33,328 events; 1,755 focal mechanisms / template events}
\kv{Content}{hypocentral locations; relative relocations; phase arrival-time picks; focal mechanisms}
\kv{Producer / contact}{Calum J. Chamberlain (Victoria University of Wellington) with W. B. Frank (MIT), F. Lanza, J. Townend, E. Warren-Smith}
\kv{Data source}{GeoNet network plus temporary deployments; matched-filter (template-matching) detection (EQcorrscan) using 1,755 template events}
\kv{Availability}{available on request $\cdot$ DOI \href{https://doi.org/10.1029/2021JB022304}{\nolinkurl{10.1029/2021JB022304}}}
\end{itemize}
\par\vspace{2pt}Dense matched-filter catalogue of 33,328 earthquakes (2009-2020) on and adjacent to the faults that ruptured in the 2016 Mw 7.8 Kaikoura earthquake, spanning the pre-seismic, co-seismic and post-seismic phases and extending well below the GeoNet magnitude of completeness. Focal mechanisms were computed for 1,755 template events. Produced to illuminate the full seismic cycle around the Kaikoura rupture; distinct from the Lanza et al. 2019 aftershock-only relocation catalogue.
\par\vspace{2pt}\textbf{Key references:} \cite{Chamberlain2021}.
\catentryrule
\subsubsection*{\hypertarget{c:27}{27.}~ Kermadec 2021 Refined Earthquake Catalogue and Slip Models (Lythgoe et al.)}
{\small\itshape aftershock $\cdot$ particular-event-type $\cdot$ manual $\cdot$ reviewed}\par\vspace{2pt}
\begin{itemize}[leftmargin=1.3em,itemsep=0.6pt,topsep=1pt,parsep=0pt]
\kv{Also known as}{Kermadec 2021 refined earthquake catalogue and slip models; Lythgoe et al. 2023 Kermadec doublet dataset; Edinburgh DataShare 10283/8556}
\kv{Coverage}{2021-03 to 2021}
\kv{Region}{Kermadec subduction zone (New Zealand offshore northern margin / Kermadec-Tonga)}
\kv{Magnitude}{Mainshocks Mw 7.4 and Mw 8.1 (5 March 2021); compared to 1976 Mw \textasciitilde{}7.9 doublet.}
\kv{Content}{hypocentral locations, relative relocations, focal mechanisms, fault-geometry, other (finite-fault slip models)}
\kv{Producer / contact}{Karen H. Lythgoe (Earth Observatory of Singapore, Nanyang Technological University; previously associated via University of Edinburgh DataShare deposit). Co-authors: K. Bradley, H. Zeng, S. Wei.}
\kv{Data source}{Refined relocations, focal mechanisms and finite-fault slip models derived for the 5 March 2021 Kermadec Mw 7.4 and Mw 8.1 doublet sequence (and re-analysis of the 1976 doublet at the same location). Based on regional/teleseismic and aftershock data analysed in Lythgoe et al. (2023, EPSL). Dataset deposited on University of Edinburgh DataShare.}
\kv{Availability}{available (DOI) $\cdot$ DOI \href{https://doi.org/10.7488/ds/7534}{\nolinkurl{10.7488/ds/7534}}}
\end{itemize}
\par\vspace{2pt}A standalone dataset accompanying Lythgoe et al. (2023) that provides a refined earthquake catalogue, focal mechanisms and slip models for the 2021 Kermadec Mw 7.4 + Mw 8.1 megathrust doublet, comparing it to the 1976 doublet that ruptured the same asperity. Produced to investigate why large earthquakes recur on the same plate-boundary segment with differing slip distributions, attributing persistent asperities to forearc structure and an isolated sedimentary basin. Focus is the Kermadec subduction zone along New Zealand's offshore northern margin.
\par\vspace{2pt}\textbf{Key references:} \cite{Lythgoe2023,Lythgoe2023a}.
\catentryrule
\subsubsection*{\hypertarget{c:28}{28.}~ 2016 Mw 7.1 Te Araroa (East Cape) Matched-Filter Foreshock \& Aftershock Catalogue (Warren-Smith et al. 2018)}
{\small\itshape aftershock $\cdot$ particular-event-type $\cdot$ template-matching $\cdot$ non-reviewed}\par\vspace{2pt}
\begin{itemize}[leftmargin=1.3em,itemsep=0.6pt,topsep=1pt,parsep=0pt]
\kv{Also known as}{Te Araroa 2016 matched-filter catalogue; East Cape 2 September 2016 Mw 7.1 foreshock/aftershock catalogue; Warren-Smith et al. 2018 EQcorrscan catalogue}
\kv{Coverage}{2016-09 to 2016-11}
\kv{Region}{Offshore East Cape / northern Hikurangi margin, North Island, New Zealand}
\kv{Magnitude}{Mainshock Mw 7.1; largest aftershock \textasciitilde{}M6.2; foreshock Mw 5.7. Matched-filter detections down to \textasciitilde{}ML 2.5-3.5 reported.}
\kv{Events}{\textasciitilde{}8,000+ earthquakes over \textasciitilde{}66 days (from 582 template events)}
\kv{Content}{hypocentral locations, phase arrival-time picks, other (template/matched-filter detections; relative timing). Relative relocation not confirmed.}
\kv{Producer / contact}{Emily Warren-Smith (lead; GNS Science). Co-authors: Bill Fry (GNS Science), Yoshihiro Kaneko, Calum J. Chamberlain (Victoria University of Wellington).}
\kv{Data source}{Matched-filter (template-matching) detection applied to GeoNet continuous waveforms using 582 well-located GeoNet template events; recorded on broadband and short-period sensors (\textasciitilde{}27 stations). Built on the GeoNet catalogue and continuous archive.}
\kv{Availability}{available (DOI) $\cdot$ DOI \href{https://doi.org/10.1016/j.epsl.2017.11.020}{\nolinkurl{10.1016/j.epsl.2017.11.020}}}
\end{itemize}
\par\vspace{2pt}A dense matched-filter foreshock and aftershock catalogue for the 2 September 2016 Mw 7.1 Te Araroa earthquake offshore East Cape, North Island. Produced to study precursory/foreshock slip, delayed triggering, and the dynamic reinvigoration of the aftershock sequence by the November 2016 Mw 7.8 Kaikoura earthquake. Detects highly correlated events beginning \textasciitilde{}12 minutes after a Mw 5.7 foreshock and documents bilateral migration near the eventual mainshock hypocentre.
\par\vspace{2pt}\textbf{Key references:} \cite{WarrenSmith2018}.
\catentryrule
\subsubsection*{\hypertarget{c:29}{29.}~ 1994 Mw 6.7 Arthur's Pass Earthquake Finely Relocated Aftershock Catalogue (Bannister et al. 2006)}
{\small\itshape aftershock $\cdot$ particular-event-type $\cdot$ manual $\cdot$ reviewed}\par\vspace{2pt}
\begin{itemize}[leftmargin=1.3em,itemsep=0.6pt,topsep=1pt,parsep=0pt]
\kv{Also known as}{Arthur's Pass 1994 finely relocated aftershocks; Bannister, Thurber \& Louie 2006 double-difference tomography catalogue}
\kv{Coverage}{1994-06 to unknown}
\kv{Region}{Arthur's Pass / central Southern Alps, South Island, New Zealand}
\kv{Depth range}{\textasciitilde{}0-10 km (\textgreater{}98\% of relocated events above \textasciitilde{}10 km; brittle-ductile transition near 10 km)}
\kv{Magnitude}{Mainshock Mw 6.7.}
\kv{Events}{\textgreater{}3,500 relocated (of \textgreater{}6,000 recorded aftershocks)}
\kv{Content}{hypocentral locations, relative relocations, phase arrival-time picks, fault-geometry, other (3-D Vp and Vp/Vs velocity model)}
\kv{Producer / contact}{Stephen Bannister (lead; GNS Science). Co-authors: Clifford Thurber (University of Wisconsin-Madison), John Louie (University of Nevada, Reno).}
\kv{Data source}{Aftershock arrival-time and waveform data following the 18 June 1994 Arthur's Pass Mw 6.7 earthquake (Southern Alps), relocated with double-difference tomography (tomoDD) using waveform cross-correlation and catalogue differential times.}
\kv{Availability}{available (DOI) $\cdot$ DOI \href{https://doi.org/10.1029/2006GL027462}{\nolinkurl{10.1029/2006GL027462}}}
\end{itemize}
\par\vspace{2pt}A finely relocated aftershock catalogue for the 18 June 1994 Mw 6.7 Arthur's Pass earthquake in the Southern Alps, South Island. More than 6,000 aftershocks were recorded and over 3,500 relocated, resolving two parallel \textasciitilde{}N60E (+/-10 deg), 70-80 deg-dipping fault clusters and a 3-D Vp / Vp/Vs model around the aftershock volume. Produced to illuminate the detailed fault structure of an unusual, complex earthquake sequence.
\par\vspace{2pt}\textbf{Key references:} \cite{Bannister2006}.
\catentryrule
\subsection{Temporary onshore passive-seismic deployments}
\textit{Published earthquake catalogues from time-limited onshore passive-seismic experiments and dense arrays; among the least-discoverable parts of the national record.}\par\vspace{4pt}
\subsubsection*{\hypertarget{c:30}{30.}~ SAPSE, Southern Alps Passive Seismic Experiment (1995-96)}
{\small\itshape background $\cdot$ geographical $\cdot$ manual $\cdot$ reviewed}\par\vspace{2pt}
\begin{itemize}[leftmargin=1.3em,itemsep=0.6pt,topsep=1pt,parsep=0pt]
\kv{Also known as}{SAPSE; FDSN XC (1995-1996); Cooperative Project on South Island New Zealand (NZ Passive)}
\kv{Coverage}{1995-11 to 1996-04}
\kv{Region}{Southern Alps / central Alpine Fault}
\kv{Bounding box}{Central South Island / Southern Alps, New Zealand (Alpine Fault region around Mt Cook / Aoraki)}
\kv{Depth range}{0-\textasciitilde{}12 km seismogenic (relocated events)}
\kv{Magnitude}{Local microseismicity; \textasciitilde{}5491 events triggered during the deployment, 195 selected and precisely relocated. Mc not formally documented in sources read.}
\kv{Events}{\textasciitilde{}5491 recorded; 195 precisely relocated}
\kv{Content}{hypocentral locations; phase arrival-time picks; focal mechanisms; quality information}
\kv{Producer / contact}{IRIS/PASSCAL with Institute of Geological \& Nuclear Sciences (GNS) and Oregon State University; lead analysis Beate Leitner / Donna Eberhart-Phillips}
\kv{Data source}{Temporary US-NZ passive deployment of \textasciitilde{}26 portable broadband/short-period seismographs across the central South Island (IRIS/PASSCAL + IGNS + Oregon State); passive partner to the active-source SIGHT transect. Network XC continuous waveforms.}
\kv{Availability}{available (DOI) $\cdot$ DOI \href{https://doi.org/10.7914/SN/XC_1995}{\nolinkurl{10.7914/SN/XC_1995}}}
\end{itemize}
\par\vspace{2pt}A focused dense temporary deployment to image microseismicity, focal mechanisms and stress along the central Alpine Fault and Southern Alps, Nov 1995-Apr 1996. Produced an early high-resolution Alpine Fault microseismicity and stress dataset and constrained the seismogenic thickness (\textasciitilde{}12 km). Served as the passive seismology component complementing the SIGHT active-source crustal transect.
\par\vspace{2pt}\textbf{Key references:} \cite{Leitner2001,IRISPASSCAL1995}.
\catentryrule
\subsubsection*{\hypertarget{c:31}{31.}~ DFDP-2 Real-Time / Matched-Filter Earthquake Catalogue, Whataroa (2014-2015)}
{\small\itshape background $\cdot$ particular-event-type $\cdot$ template-matching $\cdot$ reviewed}\par\vspace{2pt}
\begin{itemize}[leftmargin=1.3em,itemsep=0.6pt,topsep=1pt,parsep=0pt]
\kv{Also known as}{DFDP-2 monitoring catalogue; Deep Fault Drilling Project Phase 2}
\kv{Coverage}{2014-08 to 2015-01}
\kv{Region}{Central Alpine Fault / Whataroa, Westland}
\kv{Bounding box}{Within \textasciitilde{}30 km of Whataroa, central Alpine Fault, Westland, South Island, New Zealand}
\kv{Magnitude}{ML 0.6-4.2; 493 located earthquakes}
\kv{Events}{493}
\kv{Content}{hypocentral locations; phase arrival-time picks; quality information; other}
\kv{Producer / contact}{Calum J. Chamberlain \& co-authors (VUW, GNS, Univ. Auckland, Univ. Wisconsin, Tohoku Univ.)}
\kv{Data source}{Real-time monitoring network of borehole + surface seismometers within \textasciitilde{}30 km of the DFDP-2 drill site at Whataroa, central Alpine Fault, operated around-the-clock Aug 2014-Jan 2015; real-time picking plus offline matched-filter detection.}
\kv{Availability}{available on request}
\end{itemize}
\par\vspace{2pt}An earthquake catalogue produced from real-time seismic monitoring and a traffic-light response protocol during DFDP-2 scientific drilling on the central Alpine Fault. The network detected and located 493 earthquakes (ML 0.6-4.2) near Whataroa; no event occurred within 3 km of the drill site and none required changes to drilling. Demonstrates real-time + matched-filter monitoring around an active scientific borehole.
\par\vspace{2pt}\textbf{Key references:} \cite{Chamberlain2017}.
\catentryrule
\subsubsection*{\hypertarget{c:32}{32.}~ SOSA, Southland Otago Seismic Array Catalogue (2022-2023)}
{\small\itshape background $\cdot$ geographical $\cdot$ automatic $\cdot$ reviewed}\par\vspace{2pt}
\begin{itemize}[leftmargin=1.3em,itemsep=0.6pt,topsep=1pt,parsep=0pt]
\kv{Also known as}{SOSA; FDSN 6Y (2022-2023)}
\kv{Coverage}{2022-10 to 2023-10}
\kv{Region}{Southland / Otago / SE Fiordland}
\kv{Bounding box}{Western Catlins to eastern Fiordland, Southland/Otago, South Island, New Zealand}
\kv{Depth range}{0-40 km (microseismicity nucleation depths reported)}
\kv{Magnitude}{ML \textasciitilde{}0.2-3.1; 85 located earthquakes (\textasciitilde{}37\% in swarms)}
\kv{Events}{85}
\kv{Content}{hypocentral locations; relative relocations; phase arrival-time picks; quality information}
\kv{Producer / contact}{University of Otago; Jack N. Williams, Mark W. Stirling et al.}
\kv{Data source}{University of Otago temporary array (FDSN 6Y): 19 short-period stations between the western Catlins and eastern Fiordland, Oct 2022-Oct 2023, densifying sparse GeoNet coverage (\textasciitilde{}10-30 km spacing vs \textasciitilde{}100 km).}
\kv{Availability}{available (DOI) $\cdot$ DOI \href{https://doi.org/10.7914/jr68-qq17}{\nolinkurl{10.7914/jr68-qq17}}}
\end{itemize}
\par\vspace{2pt}A temporary dense array deployed to image microseismicity of the low-coverage southeastern South Island near the transpressive Australian-Pacific plate boundary. The resulting catalogue of 85 earthquakes reveals deeper-than-expected and clustered microseismicity (depths 0-40 km), including Fiordland swarms. A priority temporary deployment improving on sparse permanent coverage.
\par\vspace{2pt}\textbf{Key references:} \cite{Williams2025,UniversityofOtago2022}.
\catentryrule
\subsubsection*{\hypertarget{c:33}{33.}~ CNIPSE / NIGHT Passive, Central North Island Passive Seismic Experiment (FDSN XQ)}
{\small\itshape background $\cdot$ geographical $\cdot$ manual $\cdot$ reviewed}\par\vspace{2pt}
\begin{itemize}[leftmargin=1.3em,itemsep=0.6pt,topsep=1pt,parsep=0pt]
\kv{Also known as}{CNIPSE; NIGHT Passive; FDSN XQ (2001-2002); North Island Geophysical Transect Passive}
\kv{Coverage}{2001-01 to 2001-06}
\kv{Region}{Central North Island / Hikurangi subduction zone}
\kv{Bounding box}{Central North Island, \textasciitilde{}Bennydale to Tarawera (Hikurangi subduction system), New Zealand}
\kv{Depth range}{0-\textasciitilde{}300 km (subduction zone to \textasciitilde{}300 km depth imaged)}
\kv{Magnitude}{Local-to-regional seismicity; \textgreater{}4700 events located, 1239 used in tomography. Mc not separately documented in sources read.}
\kv{Events}{\textgreater{}4700 located (1239 used in tomography)}
\kv{Content}{hypocentral locations; phase arrival-time picks; other}
\kv{Producer / contact}{IRIS/PASSCAL with GNS Science and VUW; analysis Martin Reyners, Donna Eberhart-Phillips et al.}
\kv{Data source}{Dense temporary deployment (FDSN XQ) of \textasciitilde{}74 portable seismographs (32 broadband) across the central North Island, \textasciitilde{}Jan-Jun 2001 (IRIS/PASSCAL + GNS + VUW). Continuous waveforms; \textgreater{}4700 located earthquakes.}
\kv{Availability}{available (DOI) $\cdot$ DOI \href{https://doi.org/10.7914/SN/XQ_2001}{\nolinkurl{10.7914/SN/XQ_2001}}}
\end{itemize}
\par\vspace{2pt}A dense \textasciitilde{}6-month passive deployment to produce a high-resolution snapshot of Hikurangi subduction-zone seismicity beneath the central North Island and to support 3-D Vp/Vp-Vs tomography from the trench to \textasciitilde{}300 km depth. Over 4700 earthquakes were located within or near the network; a subset (1239 events) was used in the defining tomographic inversion.
\par\vspace{2pt}\textbf{Key references:} \cite{Reyners2006,IRISPASSCAL2001}.
\catentryrule
\subsubsection*{\hypertarget{c:34}{34.}~ Taranaki Dense Temporary Network Crustal Seismicity Catalogue (Sherburn \& White 2005)}
{\small\itshape background $\cdot$ geographical $\cdot$ manual $\cdot$ reviewed}\par\vspace{2pt}
\begin{itemize}[leftmargin=1.3em,itemsep=0.6pt,topsep=1pt,parsep=0pt]
\kv{Also known as}{Taranaki dense network; Cape Egmont fault zone deployment}
\kv{Coverage}{2001-12 to 2002-09}
\kv{Region}{Taranaki / Cape Egmont Fault Zone}
\kv{Bounding box}{Taranaki region around Mt Taranaki/Egmont and Cape Egmont Fault Zone, North Island, New Zealand}
\kv{Magnitude}{ML 1.4-4.0; 389 located earthquakes; \textgreater{}15,000 phase picks}
\kv{Events}{389}
\kv{Content}{hypocentral locations; phase arrival-time picks; quality information}
\kv{Producer / contact}{Steven Sherburn (GNS Science) \& Robert S. White (University of Cambridge)}
\kv{Data source}{Dense temporary deployment of 75 three-component broadband seismographs around Mt Taranaki/Egmont, Dec 2001-Sep 2002 (GNS Science + University of Cambridge). \textgreater{}15,000 P/S phase picks; 389 located local earthquakes.}
\kv{Availability}{available on request}
\end{itemize}
\par\vspace{2pt}A dense temporary network deployed to obtain accurate hypocentres of crustal seismicity in the Taranaki region. It located 389 local earthquakes (ML 1.4-4.0), found the Cape Egmont Fault Zone most active, and notably found no volcanic earthquakes beneath Mt Taranaki. A focused regional crustal-seismicity catalogue.
\par\vspace{2pt}\textbf{Key references:} \cite{Sherburn2005}.
\catentryrule
\subsubsection*{\hypertarget{c:35}{35.}~ Local Earthquake and Focal Mechanism Catalogue for the Northern Hikurangi Margin, December 2017-October 2018 (NZ3D onshore network)}
{\small\itshape background $\cdot$ geographical $\cdot$ manual $\cdot$ reviewed}\par\vspace{2pt}
\begin{itemize}[leftmargin=1.3em,itemsep=0.6pt,topsep=1pt,parsep=0pt]
\kv{Also known as}{NZ3D onshore earthquake/focal-mechanism catalogue; Woodward et al. (2026) Northern Hikurangi catalogue; FDSN network 3C (2017-2018) 'NZ3D FWI'}
\kv{Coverage}{2017-12 to 2018-10}
\kv{Region}{Northern Hikurangi subduction margin, near Gisborne, eastern North Island, New Zealand}
\kv{Events}{3071 high-quality local earthquakes (reported in the defining study; plus associated focal mechanisms)}
\kv{Content}{hypocentral locations, phase arrival-time picks, focal mechanisms, quality information}
\kv{Producer / contact}{Amy Woodward (lead author), Ian D. Bastow, Rebecca Bell, Imperial College London; deployment PI Rebecca Bell (Imperial College London)}
\kv{Data source}{Onshore temporary NZ3D-FWI seismograph network (FDSN code 3C, 2017-2018): 49 Guralp CMG-6TD broadband instruments (NERC GEF) deployed Dec 2017-Oct 2018, supplemented by \textasciitilde{}119 short-period DATACUBE3 (GIPP Potsdam) and 25 GSX3 (ERI Tokyo) short-period instruments (Dec 2017-Feb 2018), in a \textasciitilde{}30 km x 15 km array near Gisborne above the 10-30 km portion of the northern Hikurangi slow-slip region. Hypocentres located and focal mechanisms derived from these continuous recordings; v2 of the dataset removed the active-source airgun shots present in v1.}
\kv{Availability}{available (DOI) $\cdot$ DOI \href{https://doi.org/10.5281/zenodo.17063327}{\nolinkurl{10.5281/zenodo.17063327}}}
\end{itemize}
\par\vspace{2pt}Produced to seismologically characterise the deeper (10-30 km) portion of the northern Hikurangi margin slow-slip region near Gisborne, associating microseismicity and faulting with normal faults, subducting seamounts and fluid seep sites. The catalogue resolves well-located earthquakes and focal mechanisms (including deep normal-faulting events within the subducting slab) to trace fluid pathways rising from the slab mantle to surface seeps. It fills a gap left by the offshore HOBITSS catalogues and the NZ3D-FWI active-source velocity-model work, which did not include an onshore-network-derived hypocentre+focal-mechanism event catalogue.
\par\vspace{2pt}\textbf{Key references:} \cite{Woodward2026,Woodward2025,ImperialCollegeLondonIFDSN2017}.
\catentryrule
\subsubsection*{\hypertarget{c:36}{36.}~ Southern Lakes / Northern Fiordland High-Resolution Microseismicity Catalogue (Central Otago Seismic Array, COSA)}
{\small\itshape background $\cdot$ geographical $\cdot$ manual $\cdot$ reviewed}\par\vspace{2pt}
\begin{itemize}[leftmargin=1.3em,itemsep=0.6pt,topsep=1pt,parsep=0pt]
\kv{Also known as}{COSA microearthquake catalogue; Warren-Smith et al. (2017) Southern South Island microseismicity catalogue; FDSN network 6S (Central Otago Seismic Array)}
\kv{Coverage}{2012-06 to 2013-03}
\kv{Region}{Southern Lakes / northern Fiordland, southern South Island, New Zealand (zone adjacent to the southern Alpine Fault)}
\kv{Content}{hypocentral locations, relative relocations, phase arrival-time picks, focal mechanisms, quality information}
\kv{Producer / contact}{Emily Warren-Smith (lead author); co-authors Simon Lamb, Tim A. Stern, Euan Smith, Victoria University of Wellington (VUW). COSA network operated by Victoria University of Wellington.}
\kv{Data source}{Central Otago Seismic Array (COSA), FDSN network code 6S, \textasciitilde{}8 temporary broadband stations operated by Victoria University of Wellington supplementing the permanent GeoNet network; stations installed mid-2012 to early 2013 (\textasciitilde{}June 2012-March 2013). A new 15-month high-resolution microearthquake catalogue derived from these recordings, including relocations and (in the companion paper) focal mechanisms and stress inversions.}
\kv{Availability}{available on request}
\end{itemize}
\par\vspace{2pt}Built to image diffuse crustal-to-mantle microseismicity in the Southern Lakes and northern Fiordland region adjacent to the southern Alpine Fault, and to test mechanisms of crustal deformation in an obliquely convergent (transpressive) Australian-Pacific plate boundary zone. The catalogue documents shallow (\textless{}25 km) distributed crustal seismicity across a \textasciitilde{}150 km-wide zone plus deeper/upper-mantle events, interpreted as distributed shear and crustal thickening. A companion paper (Warren-Smith et al., 2017) derives the stress field and kinematics from the same dataset.
\par\vspace{2pt}\textbf{Key references:} \cite{WarrenSmith2017a,WarrenSmith2017b,Victoria2012}.
\catentryrule
\subsubsection*{\hypertarget{c:37}{37.}~ Dunedin Region Short-Term Broadband Array Microseismicity Catalogue (Akatore Fault)}
{\small\itshape background $\cdot$ geographical $\cdot$ template-matching $\cdot$ reviewed}\par\vspace{2pt}
\begin{itemize}[leftmargin=1.3em,itemsep=0.6pt,topsep=1pt,parsep=0pt]
\kv{Also known as}{Todd et al. (2020) Dunedin microseismicity catalogue; Dunedin short-term broadband seismic array catalogue; Akatore Fault microseismicity cluster}
\kv{Coverage}{unknown to unknown}
\kv{Region}{Dunedin region / Akatore Fault, eastern Otago, southern South Island, New Zealand}
\kv{Content}{hypocentral locations, phase arrival-time picks, quality information, other (template/matched-filter detections)}
\kv{Producer / contact}{Erin K. Todd (lead author); co-authors Mark W. Stirling, Bill Fry, Jerome Salichon, Pilar Villamor, University of Otago (Stirling) / GNS Science (Fry, Salichon, Villamor)}
\kv{Data source}{Short-term temporary broadband seismic array deployments in/around Dunedin (a low-seismicity region), supplementing the permanent GeoNet national network. Matched-filter (template-matching) analysis using a master event detected events absent from the GeoNet national catalogue, including a microearthquake cluster at the northern end of the Akatore Fault (\textless{}15 km from Dunedin city centre).}
\kv{Availability}{available on request}
\end{itemize}
\par\vspace{2pt}Produced to characterise microseismicity in a low-seismicity region (Dunedin) where the sparse permanent network under-detects small events, and to demonstrate the value of short-term dense broadband arrays for hazard assessment. Using template/matched-filter detection it revealed otherwise-undetected microseismicity, notably a cluster at the northern Akatore Fault, consistent with paleoseismic evidence that the fault is in an elevated-activity state after \textasciitilde{}110 ka of quiescence.
\par\vspace{2pt}\textbf{Key references:} \cite{Todd2020}.
\catentryrule
\subsubsection*{\hypertarget{c:38}{38.}~ Hawke's Bay Region Microearthquake Catalogue (Bannister 1988; subducting Pacific plate fine structure)}
{\small\itshape background $\cdot$ geographical $\cdot$ manual $\cdot$ reviewed}\par\vspace{2pt}
\begin{itemize}[leftmargin=1.3em,itemsep=0.6pt,topsep=1pt,parsep=0pt]
\kv{Also known as}{Bannister (1988) Hawke's Bay microseismicity catalogue; Hawke's Bay 11-station temporary network microearthquake dataset}
\kv{Coverage}{unknown to unknown}
\kv{Region}{Hawke's Bay, eastern North Island, New Zealand (Hikurangi subduction forearc)}
\kv{Events}{more than 350 micro-earthquakes (recorded over two months)}
\kv{Content}{hypocentral locations, phase arrival-time picks, other (velocity-structure inversion)}
\kv{Data source}{Temporary dense onshore network of 11 seismograph stations in the Hawke's Bay region; \textgreater{}350 micro-earthquakes recorded over a two-month deployment (mid-1980s). Hypocentres located and the data inverted for P-wave velocity/structure of the shallow subducting Pacific plate.}
\kv{Availability}{not available}
\end{itemize}
\par\vspace{2pt}Produced to image the fine structure of the shallow subducting Pacific plate beneath Hawke's Bay using a short, dense temporary seismograph deployment. From \textgreater{}350 microearthquakes the study resolved a thin (\textasciitilde{}11-12 km) low P-velocity band interpreted as subducted sediments, with the slab dipping \textasciitilde{}6$^{\circ}$ beneath Hawke Bay and steepening (\textgreater{}20$^{\circ}$) westward. It provides an early, citable regional microseismicity product for a North Island east-coast area lacking inventoried 1980s dense-array catalogues.
\par\vspace{2pt}\textbf{Key references:} \cite{Bannister1988}.
\catentryrule
\subsection{Temporary offshore / ocean-bottom (OBS) deployments}
\textit{Ocean-bottom seismometer deployments along the offshore Hikurangi margin and the Kermadec volcanic arc.}\par\vspace{4pt}
\subsubsection*{\hypertarget{c:39}{39.}~ HOBITSS Microearthquake Catalogue (2014-2015, northern Hikurangi, manual)}
{\small\itshape background $\cdot$ particular-event-type $\cdot$ manual $\cdot$ reviewed}\par\vspace{2pt}
\begin{itemize}[leftmargin=1.3em,itemsep=0.6pt,topsep=1pt,parsep=0pt]
\kv{Also known as}{HOBITSS; Hikurangi Ocean Bottom Investigation of Tremor and Slow Slip; YH (2014); Yarce et al. 2019 catalogue}
\kv{Coverage}{2014-05 to 2015-06}
\kv{Region}{Northern Hikurangi subduction margin (offshore Gisborne/Poverty Bay)}
\kv{Bounding box}{approx 37.5S-39.3S, 176.5E-179.5E (offshore Gisborne / northern Hikurangi margin, North Island)}
\kv{Magnitude}{approx M0.5-4.7 (local magnitude)}
\kv{Events}{2,313 earthquakes}
\kv{Content}{hypocentral locations; phase arrival-time picks}
\kv{Producer / contact}{Jefferson Yarce / Anne F. Sheehan (CIRES, University of Colorado Boulder); HOBITSS PIs incl. Laura Wallace, Susan Schwartz, Spahr Webb, Kimihiro Mochizuki, Stuart Henrys; network operated by University of Texas at Austin}
\kv{Data source}{HOBITSS ocean-bottom seismometer array (FDSN network YH, 2014) deployed on the northern Hikurangi margin seafloor, supplemented by onshore GeoNet permanent stations; manual P- and S-wave picking and hypocentre location}
\kv{Availability}{available (DOI) $\cdot$ DOI \href{https://doi.org/10.5281/zenodo.2022405}{\nolinkurl{10.5281/zenodo.2022405}}}
\end{itemize}
\par\vspace{2pt}Built to investigate whether a spatial/temporal relationship exists between microseismicity and the shallow \textasciitilde{}2-week September-October 2014 slow slip event (SSE) targeted by the year-long HOBITSS deployment above the northern Hikurangi plate interface. Earthquakes were detected and located by hand-picking P/S arrivals on OBS plus onshore GeoNet stations. The catalogue reveals a 'microseismicity gap' bordering the 2014 slow-slip patch, later reused by Yarce et al. (2023).
\par{\small\itshape Related: Companion products from the same 2014-2015 HOBITSS ocean-bottom deployment: the automatic-detection earthquake catalogue (Yarce et al. 2023) and the focal-mechanism / stress-tensor catalogue (Warren-Smith et al. 2019). This is the manually picked microearthquake catalogue.}
\par\vspace{2pt}\textbf{Key references:} \cite{Yarce2019,Yarce2018,UniversityofTexasatAustin2014}.
\catentryrule
\subsubsection*{\hypertarget{c:40}{40.}~ HOBITSS Automatic-Detection Earthquake Catalogue (2014 SSE, REST detector)}
{\small\itshape background $\cdot$ particular-event-type $\cdot$ automatic $\cdot$ non-reviewed}\par\vspace{2pt}
\begin{itemize}[leftmargin=1.3em,itemsep=0.6pt,topsep=1pt,parsep=0pt]
\kv{Also known as}{HOBITSS auto-detection catalogue; Yarce et al. 2023; REST (Regressive ESTimator) catalogue}
\kv{Coverage}{2014-09 to 2014-10}
\kv{Region}{Northern Hikurangi subduction margin (offshore Gisborne)}
\kv{Bounding box}{approx 37.5S-39.3S, 176.5E-179.5E (offshore northern Hikurangi margin; 3D velocity model box of Yarce et al. 2021 down to \textasciitilde{}94 km)}
\kv{Depth range}{approx 0-15 km plate interface focus; full catalogue depth extent not separately stated (velocity model to \textasciitilde{}94 km)}
\kv{Magnitude}{Magnitude of completeness Mc = 1.7 and Gutenberg-Richter b-value = 1.3, derived from the 451 events with REST-based magnitude estimates (subset of catalogue); local magnitude ML scale}
\kv{Events}{854 events (per Zenodo dataset record); paper text reports \textgreater{}3x the manual analyst count and 451 events with magnitude estimates}
\kv{Content}{hypocentral locations; phase arrival-time picks; quality information}
\kv{Producer / contact}{Jefferson Yarce / Anne F. Sheehan (CIRES \& Dept. of Geological Sciences, University of Colorado Boulder); Steven Roecker (Rensselaer Polytechnic Institute)}
\kv{Data source}{Continuous waveform data from the HOBITSS ocean-bottom seismometers (FDSN network YH, 2014) plus onshore GeoNet land stations; events detected and located automatically using the REST (Regressive ESTimator) autoregressive phase-arrival-time detector and onset estimator}
\kv{Availability}{available (DOI) $\cdot$ DOI \href{https://doi.org/10.5281/zenodo.7274399}{\nolinkurl{10.5281/zenodo.7274399}}}
\end{itemize}
\par\vspace{2pt}An automatically generated earthquake catalogue for the two months around the September-October 2014 northern Hikurangi slow slip event, created to test whether automated detection increases catalogue completeness and to examine the temporal relationship between SSEs and microseismicity. REST applied to HOBITSS OBS + GeoNet data found more than three times as many events as manual analysis, and shows seismicity increasing during and remaining elevated for \textasciitilde{}2 weeks after the 2014 SSE. A larger, distinct product from the same OBS deployment as the manual Yarce et al. (2019) catalogue.
\par{\small\itshape Related: The automatic-detection counterpart (REST detector, focused on the 2014 slow-slip event) to the manually picked HOBITSS microearthquake catalogue (Yarce et al. 2019).}
\par\vspace{2pt}\textbf{Key references:} \cite{Yarce2023,Yarce2022,UniversityofTexasatAustin2014}.
\catentryrule
\subsubsection*{\hypertarget{c:41}{41.}~ Brothers Volcano Ocean-Bottom Hydrophone Seismicity Catalogue (south Kermadec arc)}
{\small\itshape background $\cdot$ geographical $\cdot$ other $\cdot$ reviewed}\par\vspace{2pt}
\begin{itemize}[leftmargin=1.3em,itemsep=0.6pt,topsep=1pt,parsep=0pt]
\kv{Also known as}{Brothers OBH array; Kermadec arc OBH T-wave catalogue}
\kv{Coverage}{2004 to 2005}
\kv{Region}{Southern Kermadec arc (Brothers submarine volcano)}
\kv{Bounding box}{Brothers volcano caldera, southern Kermadec arc, approx 34.9S, 179.1E (\textasciitilde{}350 km NE of New Zealand / \textasciitilde{}340 km NE of Whakaari/White Island)}
\kv{Events}{964 regional earthquakes (plus local harmonic tremor)}
\kv{Content}{hypocentral locations; other}
\kv{Producer / contact}{Robert P. Dziak / Joseph H. Haxel / Haru Matsumoto (NOAA Pacific Marine Environmental Laboratory \& Oregon State University, Hatfield Marine Science Center); with Cornel E. J. de Ronde and collaborators (GNS Science / NOAA Kermadec-arc program)}
\kv{Data source}{Four ocean-bottom hydrophones (OBH) deployed on the caldera floor of Brothers submarine volcano for \textasciitilde{}7 months (2004-2005); low-frequency (0.5-110 Hz) hydroacoustic T-wave arrivals recorded on three of the four OBH used to locate earthquakes and characterize harmonic tremor}
\kv{Availability}{available (DOI) $\cdot$ DOI \href{https://doi.org/10.1029/2007JB005533}{\nolinkurl{10.1029/2007JB005533}}}
\end{itemize}
\par\vspace{2pt}A hydroacoustic (T-wave) catalogue produced to characterize regional and local seismicity and volcanic harmonic tremor at Brothers volcano in the southern Kermadec intraoceanic arc (\textasciitilde{}350 km NE of New Zealand). T-wave-derived locations for 964 regional earthquakes show clustering beneath the dacite cone and along the east flank of the caldera, while local harmonic tremor reflects subseafloor hydrothermal/magmatic processes. It is the only OBS/OBH-based seismicity catalogue identified for the NZ Kermadec arc.
\par\vspace{2pt}\textbf{Key references:} \cite{Dziak2008}.
\catentryrule
\subsubsection*{\hypertarget{c:42}{42.}~ Dataset of Earthquakes in the Northern Hikurangi Subduction Zone During the 2019 Slow Slip Event (2018-2019 OBS deployment)}
{\small\itshape background $\cdot$ particular-event-type $\cdot$ template-matching $\cdot$ reviewed}\par\vspace{2pt}
\begin{itemize}[leftmargin=1.3em,itemsep=0.6pt,topsep=1pt,parsep=0pt]
\kv{Also known as}{Iwasaki et al. (2025) northern Hikurangi 2019 SSE OBS catalogue; Zenodo 15713529}
\kv{Coverage}{2018 to 2019}
\kv{Region}{Northern Hikurangi subduction zone, offshore Gisborne, Aotearoa New Zealand}
\kv{Events}{964 located events}
\kv{Content}{hypocentral locations, phase arrival-time picks, other (event waveforms, instrument responses, station coordinates)}
\kv{Producer / contact}{Y. (Yuriko) Iwasaki (lead author); co-authors S. Schwartz, K. Mochizuki, H. Crume, D. Bassett, T. Yamada, D. Morad}
\kv{Data source}{Ocean-bottom seismometer (OBS) network deployed 2018-2019 directly above the main slip region of the 2019 northern Hikurangi slow slip event, offshore Gisborne; high-resolution catalogue built with a detailed velocity model plus machine-learning and template-matching detection. Archive includes the earthquake catalogue, three-component OBS event waveforms, instrument (pole-zero) responses, and station coordinates.}
\kv{Availability}{available (DOI) $\cdot$ DOI \href{https://doi.org/10.5281/zenodo.15713529}{\nolinkurl{10.5281/zenodo.15713529}}}
\end{itemize}
\par\vspace{2pt}An OBS earthquake catalogue from a deployment directly above the main slip region of the large 2019 northern Hikurangi slow slip event, offshore Gisborne, New Zealand. Produced to characterise how the SSE modulated nearby seismicity, revealing intraslab seismicity near subducted seamounts reactivated by the slow slip and potentially influenced by fluid migration. Distinct from the earlier HOBITSS 2014-2015 manual (Yarce et al. 2019) and auto-detection (Yarce et al. 2023) catalogues - this covers the later 2018-2019 OBS deployment and the 2019 SSE.
\par\vspace{2pt}\textbf{Key references:} \cite{Iwasaki2025,Iwasaki2025a}.
\catentryrule
\subsection{Volcano-seismic, geothermal and induced catalogues}
\textit{Volcano-monitoring, magmatic-unrest, and geothermal injection-induced microseismicity catalogues, typically from dense local or borehole networks.}\par\vspace{4pt}
\subsubsection*{\hypertarget{c:43}{43.}~ Tongariro Volcanic Centre Local Earthquake Tomography Catalogue (Rowlands et al. 2005)}
{\small\itshape background $\cdot$ particular-event-type $\cdot$ manual $\cdot$ reviewed}\par\vspace{2pt}
\begin{itemize}[leftmargin=1.3em,itemsep=0.6pt,topsep=1pt,parsep=0pt]
\kv{Also known as}{TgVC tomography; Rowlands, White \& Haines 2005}
\kv{Coverage}{2001-01 to 2001-06}
\kv{Region}{Tongariro Volcanic Centre (TVZ)}
\kv{Bounding box}{Tongariro Volcanic Centre, central North Island, New Zealand}
\kv{Magnitude}{Local volcanic-tectonic earthquakes; 155 selected events. Mc not separately documented in sources read.}
\kv{Events}{155 (selected for tomography)}
\kv{Content}{hypocentral locations; phase arrival-time picks; quality information}
\kv{Producer / contact}{D. P. Rowlands, R. S. White (University of Cambridge) \& A. J. Haines (GNS Science)}
\kv{Data source}{Temporary deployment of 23 broadband seismometers in the Tongariro Volcanic Centre, Jan-Jun 2001 (University of Cambridge / GNS), part of the same 2001 central North Island campaign as CNIPSE. \textasciitilde{}3000 P + \textasciitilde{}1100 S picks; 155 selected local earthquakes.}
\kv{Availability}{available on request}
\end{itemize}
\par\vspace{2pt}A temporary local network deployed to obtain accurate earthquake depths and 3-D Vp / Vp-Vs structure of the Tongariro Volcanic Centre. The best 155 local earthquakes (with quality criteria on picks and azimuthal gap) were selected for local-earthquake tomography of the volcanic system. A focused volcanic-zone seismicity/tomography catalogue.
\par\vspace{2pt}\textbf{Key references:} \cite{Rowlands2005}.
\catentryrule
\subsubsection*{\hypertarget{c:44}{44.}~ Okataina Volcanic Centre Double-Difference Relocated Catalogue (2010-2021)}
{\small\itshape background $\cdot$ geographical $\cdot$ manual $\cdot$ reviewed}\par\vspace{2pt}
\begin{itemize}[leftmargin=1.3em,itemsep=0.6pt,topsep=1pt,parsep=0pt]
\kv{Also known as}{OVC double-difference catalogue; Bannister et al. 2022 Okataina tomography catalogue}
\kv{Coverage}{2010 to 2021}
\kv{Region}{Okataina Volcanic Centre, Taupo Volcanic Zone}
\kv{Bounding box}{Okataina Volcanic Centre / eastern Taupo Volcanic Zone, Rotorua region, North Island, New Zealand}
\kv{Depth range}{low-velocity magmatic zone resolved at \textasciitilde{}5-8 km depth}
\kv{Events}{6989}
\kv{Content}{hypocentral locations; relative relocations; phase arrival-time picks; other}
\kv{Producer / contact}{Stephen Bannister, GNS Science}
\kv{Data source}{Temporary array of broadband and short-period seismometers (\textasciitilde{}4+ years) supplementing the permanent GeoNet New Zealand National Seismograph Network; double-difference tomography with waveform cross-correlation}
\kv{Availability}{available on request $\cdot$ DOI \href{https://doi.org/10.1016/j.jvolgeores.2022.107653}{\nolinkurl{10.1016/j.jvolgeores.2022.107653}}}
\end{itemize}
\par\vspace{2pt}Produced to image sub-caldera structure and seismicity beneath the Okataina Volcanic Centre (TVZ, Rotorua region) using a 3-D P-wave velocity model derived by double-difference seismic tomography. The study relocated 6989 earthquakes (2010-2021) to resolve spatial clusters (e.g. west of Lake Rotomahana, beneath Haroharo/Makatiti) and a \textasciitilde{}5-8 km low-velocity zone interpreted as partial melt and/or fluid-volatile pathways.
\par\vspace{2pt}\textbf{Key references:} \cite{Bannister2022}.
\catentryrule
\subsubsection*{\hypertarget{c:45}{45.}~ Taupo Volcano Magmatic-Seismicity Catalogue (2019-2022 ECLIPSE temporary deployment)}
{\small\itshape background $\cdot$ geographical $\cdot$ manual $\cdot$ reviewed}\par\vspace{2pt}
\begin{itemize}[leftmargin=1.3em,itemsep=0.6pt,topsep=1pt,parsep=0pt]
\kv{Also known as}{Mestel et al. 2025 Taupo catalogue; ECLIPSE Taupo deployment}
\kv{Coverage}{2019 to 2022}
\kv{Region}{Taupo Volcano / Taupo Volcanic Zone}
\kv{Bounding box}{Taupo caldera / Lake Taupo, central Taupo Volcanic Zone, North Island, New Zealand}
\kv{Depth range}{arcuate ring-fault structure resolved at \textasciitilde{}6 km depth}
\kv{Content}{hypocentral locations; phase arrival-time picks; fault-geometry; other}
\kv{Producer / contact}{Eleanor R. H. Mestel / Finnigan Illsley-Kemp, Victoria University of Wellington (ECLIPSE programme)}
\kv{Data source}{Dedicated temporary broadband deployment of \textasciitilde{}13 seismometers (ECLIPSE programme) supplementing the GeoNet permanent network around Lake Taupo}
\kv{Availability}{available on request $\cdot$ DOI \href{https://doi.org/10.1029/2025JB031644}{\nolinkurl{10.1029/2025JB031644}}}
\end{itemize}
\par\vspace{2pt}Built from a dedicated temporary broadband deployment (ECLIPSE programme) to improve location control under Lake Taupo where the permanent network gives poor coverage, and to separate magmatic from tectonic seismicity through the 2022-2023 unrest. Resolves an arcuate caldera-ring-fault structure at roughly 6 km depth. A priority temporary-deployment catalogue.
\par{\small\itshape Related: Overlaps the 2019 Taupo unrest catalogue (Illsley-Kemp et al. 2021) and the 2019 northern-TVZ cascading-swarm catalogue (Aber et al. 2025); this entry spans the longer 2019-2022 ECLIPSE deployment.}
\par\vspace{2pt}\textbf{Key references:} \cite{Mestel2025}.
\catentryrule
\subsubsection*{\hypertarget{c:46}{46.}~ Taupo 2019 Volcanic Unrest Earthquake Catalogue (Illsley-Kemp et al. 2021)}
{\small\itshape background $\cdot$ particular-event-type $\cdot$ template-matching $\cdot$ reviewed}\par\vspace{2pt}
\begin{itemize}[leftmargin=1.3em,itemsep=0.6pt,topsep=1pt,parsep=0pt]
\kv{Also known as}{Illsley-Kemp et al. 2021 Taupo unrest catalogue}
\kv{Coverage}{2019 to 2019}
\kv{Region}{Taupo Volcano / Taupo Volcanic Zone}
\kv{Bounding box}{Taupo caldera / Lake Taupo, central Taupo Volcanic Zone, North Island, New Zealand}
\kv{Magnitude}{up to \textasciitilde{}M5.2 (largest event of the 2019 unrest)}
\kv{Events}{\textgreater{}1,100 (caldera earthquakes)}
\kv{Content}{hypocentral locations; phase arrival-time picks; other}
\kv{Producer / contact}{Finnigan Illsley-Kemp, Victoria University of Wellington (VUW / EQcorrscan team)}
\kv{Data source}{GeoNet network data analysed by the VUW group; integrated with deformation (GNSS/InSAR) and gas observations}
\kv{Availability}{available on request $\cdot$ DOI \href{https://doi.org/10.1029/2021GC009803}{\nolinkurl{10.1029/2021GC009803}}}
\end{itemize}
\par\vspace{2pt}Catalogue of the first instrumentally recorded Taupo caldera unrest episode (2019), documenting \textgreater{}1,100 caldera earthquakes culminating in an M5.2 event, interpreted as related to a magma-mush margin. Integrates seismicity with ground deformation and gas data to assess causes, mechanisms and implications of the unrest.
\par{\small\itshape Related: The 2019 events also appear in the longer 2019-2022 magmatic-seismicity catalogue (Mestel et al. 2025); northern-TVZ swarms of the same year are catalogued separately (Aber et al. 2025).}
\par\vspace{2pt}\textbf{Key references:} \cite{IllsleyKemp2021}.
\catentryrule
\subsubsection*{\hypertarget{c:47}{47.}~ Taupo 2022-2023 Volcanic Unrest Seismicity Catalogue (Lamb et al. 2024)}
{\small\itshape background $\cdot$ particular-event-type $\cdot$ manual $\cdot$ reviewed}\par\vspace{2pt}
\begin{itemize}[leftmargin=1.3em,itemsep=0.6pt,topsep=1pt,parsep=0pt]
\kv{Also known as}{Lamb et al. 2024 Taupo unrest catalogue}
\kv{Coverage}{2022-05 to 2023-05}
\kv{Region}{Taupo Volcano / Taupo Volcanic Zone}
\kv{Bounding box}{Taupo caldera / Lake Taupo, central Taupo Volcanic Zone, North Island, New Zealand}
\kv{Magnitude}{largest event M5.7 (largest beneath Lake Taupo in \textasciitilde{}50 years)}
\kv{Content}{hypocentral locations; phase arrival-time picks; focal mechanisms; other}
\kv{Producer / contact}{Oliver Lamb / Stephen Bannister, GNS Science / GeoNet}
\kv{Data source}{GeoNet network data; moment-tensor and velocity-model analysis by GNS Science}
\kv{Availability}{available (DOI) $\cdot$ DOI \href{https://doi.org/10.26443/seismica.v3i2.1125}{\nolinkurl{10.26443/seismica.v3i2.1125}}}
\end{itemize}
\par\vspace{2pt}Characterises the May 2022-May 2023 Taupo unrest episode, which produced the highest earthquake count of any instrumented Taupo unrest and the largest earthquake beneath the lake in \textasciitilde{}50 years (M5.7). Moment-tensor analysis of the largest event shows an inflationary isotropic component, leading the authors to infer magmatic intrusion at depth triggering fault reactivation.
\par\vspace{2pt}\textbf{Key references:} \cite{Lamb2024}.
\catentryrule
\subsubsection*{\hypertarget{c:48}{48.}~ Cascading Earthquake Swarms, Northern Taupo Volcanic Zone (2019, Aber et al. 2025)}
{\small\itshape background $\cdot$ geographical $\cdot$ other $\cdot$ reviewed}\par\vspace{2pt}
\begin{itemize}[leftmargin=1.3em,itemsep=0.6pt,topsep=1pt,parsep=0pt]
\kv{Also known as}{Aber et al. 2025 cascading swarms catalogue}
\kv{Coverage}{2019-03 to 2019-09}
\kv{Region}{Northern/central Taupo Volcanic Zone (Okataina, Rotorua, toward Whakaari)}
\kv{Bounding box}{\textasciitilde{}175 km of the northern and central Taupo Volcanic Zone, North Island, New Zealand}
\kv{Events}{10 swarms (number of individual events not verified)}
\kv{Content}{hypocentral locations; other}
\kv{Producer / contact}{S. Aber / C. J. Ebinger (with VUW and GNS co-authors)}
\kv{Data source}{GeoNet network data and temporary deployments across the northern/central TVZ; swarm detection and relocation}
\kv{Availability}{available on request $\cdot$ DOI \href{https://doi.org/10.1029/2024GC012079}{\nolinkurl{10.1029/2024GC012079}}}
\end{itemize}
\par\vspace{2pt}Documents and analyses 10 spatio-temporal earthquake swarms cascading along \textasciitilde{}175 km of the northern and central Taupo Volcanic Zone (March-September 2019), near Okataina and Rotorua and extending toward Whakaari/White Island, interpreted as fault-magma interaction. Treats the same 2019 episode as Illsley-Kemp et al. 2021 but focuses on the broader northern-TVZ cascading sequence rather than the Taupo caldera.
\par{\small\itshape Related: Covers the 2019 northern-TVZ swarms; the Taupo caldera unrest of the same year is catalogued by Illsley-Kemp et al. (2021) and Mestel et al. (2025).}
\par\vspace{2pt}\textbf{Key references:} \cite{Aber2025}.
\catentryrule
\subsubsection*{\hypertarget{c:49}{49.}~ Tarawera 2019 Dyke-Intrusion Earthquake Catalogue (Puhipuhi, Okataina; Benson et al. 2021)}
{\small\itshape background $\cdot$ particular-event-type $\cdot$ template-matching $\cdot$ reviewed}\par\vspace{2pt}
\begin{itemize}[leftmargin=1.3em,itemsep=0.6pt,topsep=1pt,parsep=0pt]
\kv{Also known as}{Benson et al. 2021 Tarawera/Puhipuhi swarm catalogue}
\kv{Coverage}{2019-03 to 2019-03}
\kv{Region}{Tarawera / Okataina Volcanic Centre, Taupo Volcanic Zone}
\kv{Bounding box}{Puhipuhi embayment, NE of Mt Tarawera, Okataina Volcanic Centre, North Island, New Zealand}
\kv{Depth range}{\textasciitilde{}8-10.5 km}
\kv{Magnitude}{ML -1.12 to 2.54}
\kv{Events}{348 detected; 94 relocated}
\kv{Content}{hypocentral locations; relative relocations; phase arrival-time picks; other}
\kv{Producer / contact}{Thomas W. Benson / Finnigan Illsley-Kemp, Victoria University of Wellington}
\kv{Data source}{GeoNet network data, analysed with ObsPy and GrowClust (matched-filter detection + relative relocation)}
\kv{Availability}{available (DOI) $\cdot$ DOI \href{https://doi.org/10.5281/zenodo.4035171}{\nolinkurl{10.5281/zenodo.4035171}}}
\end{itemize}
\par\vspace{2pt}Catalogue of a March 2019 earthquake swarm in the Puhipuhi embayment NE of Tarawera (Okataina Volcanic Centre), interpreted as a dyke intrusion. Detected 348 events and relocated 94 of them (ML -1.12 to 2.54, \textasciitilde{}8-10.5 km depth), demonstrating magmatic intrusion during the 2019 activity.
\par\vspace{2pt}\textbf{Key references:} \cite{Benson2021,Benson2020}.
\catentryrule
\subsubsection*{\hypertarget{c:50}{50.}~ Waimangu-Rotomahana-Mt Tarawera Geothermal Field Earthquake Swarm Catalogue (2004-2015)}
{\small\itshape background $\cdot$ geographical $\cdot$ manual $\cdot$ reviewed}\par\vspace{2pt}
\begin{itemize}[leftmargin=1.3em,itemsep=0.6pt,topsep=1pt,parsep=0pt]
\kv{Also known as}{Bannister et al. 2016 WRT swarm catalogue}
\kv{Coverage}{2004 to 2015}
\kv{Region}{Waimangu-Rotomahana-Tarawera, Okataina / Taupo Volcanic Zone}
\kv{Bounding box}{Waimangu-Rotomahana-Mt Tarawera geothermal field, SW Okataina Volcanic Centre, North Island, New Zealand}
\kv{Events}{582 (relocated)}
\kv{Content}{hypocentral locations; relative relocations; phase arrival-time picks}
\kv{Producer / contact}{Stephen Bannister, GNS Science}
\kv{Data source}{GeoNet / GNS network data; earthquake relocation of swarm seismicity}
\kv{Availability}{not available $\cdot$ DOI \href{https://doi.org/10.1016/j.jvolgeores.2015.07.024}{\nolinkurl{10.1016/j.jvolgeores.2015.07.024}}}
\end{itemize}
\par\vspace{2pt}Catalogue of 582 relocated earthquakes (2004-2015) in the Waimangu-Rotomahana-Mt Tarawera (WRT) geothermal field on the SW flank of the Okataina Volcanic Centre. The swarm activity highlights crustal faulting associated with the geothermal/volcanic system in the Taupo Volcanic Zone.
\par\vspace{2pt}\textbf{Key references:} \cite{Bannister2016}.
\catentryrule
\subsubsection*{\hypertarget{c:51}{51.}~ Mt Ruapehu Earthquake Swarm Catalogue (1990-2023)}
{\small\itshape background $\cdot$ geographical $\cdot$ machine-learning $\cdot$ reviewed}\par\vspace{2pt}
\begin{itemize}[leftmargin=1.3em,itemsep=0.6pt,topsep=1pt,parsep=0pt]
\kv{Also known as}{Mitchinson et al. 2024 Ruapehu ML catalogue}
\kv{Coverage}{1990 to 2023}
\kv{Region}{Mt Ruapehu, Taupo Volcanic Zone}
\kv{Bounding box}{Mt Ruapehu region, southern Taupo Volcanic Zone, North Island, New Zealand}
\kv{Magnitude}{ML \textasciitilde{}1.58-4.21; Mc \textasciitilde{}1.6}
\kv{Events}{2,795 swarm events (from 28,522 catalogued)}
\kv{Content}{hypocentral locations; other}
\kv{Producer / contact}{Sam Mitchinson / Jessica H. Johnson, University of East Anglia}
\kv{Data source}{GeoNet QuakeSearch earthquake catalogue (33 years of locations); unsupervised machine-learning clustering (HDBSCAN/DBSCAN)}
\kv{Availability}{available (DOI) $\cdot$ DOI \href{https://doi.org/10.3389/feart.2024.1343874}{\nolinkurl{10.3389/feart.2024.1343874}}}
\end{itemize}
\par\vspace{2pt}A derived catalogue isolating earthquake swarms at Mt Ruapehu from decades of GeoNet QuakeSearch data using unsupervised machine-learning clustering (HDBSCAN/DBSCAN). From 28,522 catalogued events it identifies 2,795 swarm-related events in 8 sequences (ML \textasciitilde{}1.58-4.21, Mc \textasciitilde{}1.6) to test whether the 'Erua' swarm was distinctive in space, time and magnitude.
\par\vspace{2pt}\textbf{Key references:} \cite{Mitchinson2024}.
\catentryrule
\subsubsection*{\hypertarget{c:52}{52.}~ Whakaari/White Island Very-Long-Period (VLP) Earthquake Catalogue (2007-2019)}
{\small\itshape background $\cdot$ particular-event-type $\cdot$ other $\cdot$ reviewed}\par\vspace{2pt}
\begin{itemize}[leftmargin=1.3em,itemsep=0.6pt,topsep=1pt,parsep=0pt]
\kv{Also known as}{Park et al. 2020 Whakaari VLP families F1/F2}
\kv{Coverage}{2007 to 2019}
\kv{Region}{Whakaari/White Island, Taupo Volcanic Zone (offshore)}
\kv{Bounding box}{Whakaari/White Island, Bay of Plenty, offshore North Island, New Zealand}
\kv{Content}{hypocentral locations; quality information; other}
\kv{Producer / contact}{Iseul Park / Arthur Jolly, University of Canterbury / GNS Science}
\kv{Data source}{GeoNet seismic stations on Whakaari/White Island; waveform-semblance detection and clustering of VLP events}
\kv{Availability}{available on request $\cdot$ DOI \href{https://doi.org/10.1186/s40623-020-01224-z}{\nolinkurl{10.1186/s40623-020-01224-z}}}
\end{itemize}
\par\vspace{2pt}A long-term classification catalogue of very-long-period (VLP) volcanic earthquakes at Whakaari/White Island over 2007 to end-2019, detected by waveform semblance and grouped into two families (F1, F2). F2 swarms mark the onset of unrest, making the catalogue relevant to eruption forecasting at the volcano.
\par\vspace{2pt}\textbf{Key references:} \cite{Park2020}.
\catentryrule
\subsubsection*{\hypertarget{c:53}{53.}~ Ngatamariki Geothermal Injection-Induced Microseismicity Catalogue (2012-2015)}
{\small\itshape induced $\cdot$ particular-event-type $\cdot$ template-matching $\cdot$ reviewed}\par\vspace{2pt}
\begin{itemize}[leftmargin=1.3em,itemsep=0.6pt,topsep=1pt,parsep=0pt]
\kv{Also known as}{Hopp et al. 2019 Ngatamariki matched-filter catalogue}
\kv{Coverage}{2012 to 2015}
\kv{Region}{Ngatamariki geothermal field, Taupo Volcanic Zone}
\kv{Bounding box}{Ngatamariki geothermal field, central Taupo Volcanic Zone, North Island, New Zealand}
\kv{Depth range}{\textasciitilde{}1-3 km (reservoir)}
\kv{Magnitude}{ML \textasciitilde{}0.26-3.17}
\kv{Events}{\textasciitilde{}77,457 total (templates + detections); \textasciitilde{}41,114 located (approximate)}
\kv{Content}{hypocentral locations; relative relocations; phase arrival-time picks; focal mechanisms; other}
\kv{Producer / contact}{Chet Hopp / John Townend \& Martha Savage, Victoria University of Wellington (with Mercury NZ)}
\kv{Data source}{Mercury NZ operator temporary seismic network at Ngatamariki plus GeoNet; matched-filter (template-matching) detection}
\kv{Availability}{available on request $\cdot$ DOI \href{https://doi.org/10.1029/2019GC008243}{\nolinkurl{10.1029/2019GC008243}}}
\end{itemize}
\par\vspace{2pt}Catalogue of injection/stimulation-induced microseismicity at the Ngatamariki geothermal field tied to 2012-2015 well stimulation and injection, built by matched-filter detection. Reports very large numbers of detections (templates plus matched-filter detections), a located subset in a \textasciitilde{}1-3 km reservoir (ML \textasciitilde{}0.26-3.17), and includes focal mechanisms, examining the relationship between injection and induced seismicity.
\par\vspace{2pt}\textbf{Key references:} \cite{Hopp2019}.
\catentryrule
\subsubsection*{\hypertarget{c:54}{54.}~ Rotokawa Geothermal Field Induced Microseismicity Catalogue (2012-2015)}
{\small\itshape induced $\cdot$ particular-event-type $\cdot$ template-matching $\cdot$ reviewed}\par\vspace{2pt}
\begin{itemize}[leftmargin=1.3em,itemsep=0.6pt,topsep=1pt,parsep=0pt]
\kv{Also known as}{Hopp et al. 2020 Rotokawa matched-filter catalogue}
\kv{Coverage}{2012 to 2015}
\kv{Region}{Rotokawa geothermal field, Taupo Volcanic Zone}
\kv{Bounding box}{Rotokawa geothermal field, central Taupo Volcanic Zone, North Island, New Zealand}
\kv{Depth range}{\textasciitilde{}1-3 km (reservoir)}
\kv{Magnitude}{ML \textasciitilde{}0.37-3.52; Mc \textasciitilde{}0.74}
\kv{Events}{\textasciitilde{}25,148 events (2,665 templates); \textasciitilde{}6,500 relocated (approximate)}
\kv{Content}{hypocentral locations; relative relocations; phase arrival-time picks; other}
\kv{Producer / contact}{Chet Hopp / John Townend \& Martha Savage, Victoria University of Wellington (with Mercury NZ)}
\kv{Data source}{Mercury NZ operator temporary seismic network at Rotokawa plus GeoNet; matched-filter (template-matching) detection and double-difference relocation}
\kv{Availability}{available on request $\cdot$ DOI \href{https://doi.org/10.1016/j.geothermics.2019.101750}{\nolinkurl{10.1016/j.geothermics.2019.101750}}}
\end{itemize}
\par\vspace{2pt}A \textasciitilde{}four-year matched-filter catalogue of induced/produced seismicity at the Rotokawa geothermal field, examining the seismic response to evolving injection. Reports thousands of templates yielding tens of thousands of detected events, with a relocated subset in a \textasciitilde{}1-3 km reservoir (ML \textasciitilde{}0.37-3.52, Mc \textasciitilde{}0.74).
\par\vspace{2pt}\textbf{Key references:} \cite{Hopp2020}.
\catentryrule
\subsubsection*{\hypertarget{c:55}{55.}~ Rotorua and Kawerau Geothermal Systems Relocated Seismicity Catalogue (1984-2004; Clarke et al. 2009)}
{\small\itshape background $\cdot$ geographical $\cdot$ manual $\cdot$ reviewed}\par\vspace{2pt}
\begin{itemize}[leftmargin=1.3em,itemsep=0.6pt,topsep=1pt,parsep=0pt]
\kv{Also known as}{Clarke et al. 2009 Rotorua-Kawerau relocated seismicity; Rotorua and Kawerau geothermal seismicity (TVZ)}
\kv{Coverage}{1984 to 2004}
\kv{Region}{Rotorua and Kawerau geothermal fields, central/eastern Taupo Volcanic Zone, North Island, New Zealand}
\kv{Bounding box}{Two sub-areas in the central/eastern Taupo Volcanic Zone, North Island, New Zealand: Rotorua geothermal system (\textasciitilde{}38.1S, 176.25E) and Kawerau geothermal field (\textasciitilde{}38.07S, 176.74E).}
\kv{Depth range}{Shallow crustal, approximately 0-20 km (datasets defined as events \textless{}=20 km depth)}
\kv{Magnitude}{Local magnitude (ML); low-magnitude microseismicity to small earthquakes.}
\kv{Events}{Input datasets: 504 shallow earthquakes around Rotorua and 1875 around Kawerau (1984-2004); relocated hypocentres computed for 155 (Rotorua) and 400 (Kawerau) earthquakes using combined cross-correlation and catalogue arrival times.}
\kv{Content}{hypocentral locations, relative relocations, phase arrival-time picks, quality information, focal mechanisms}
\kv{Producer / contact}{D. Clarke (lead author); J. Townend, M.K. Savage (Victoria University of Wellington); S. Bannister (GNS Science). Underlying picks/waveforms: GeoNet / GNS Science.}
\kv{Data source}{Earthquakes recorded by the New Zealand national seismograph network and local stations around the Rotorua and Kawerau geothermal areas; parent catalogue arrival-time picks plus newly computed cross-correlation differential travel times. New 1-D P/S velocity models and station corrections derived with VELEST; relocation with HypoDD double-difference algorithm.}
\kv{Availability}{available on request $\cdot$ DOI \href{https://doi.org/10.1016/j.jvolgeores.2008.11.004}{\nolinkurl{10.1016/j.jvolgeores.2008.11.004}}}
\end{itemize}
\par\vspace{2pt}A dedicated double-difference (HypoDD) relocated earthquake catalogue for two Taupo Volcanic Zone geothermal systems (Rotorua city / Mt Ngongotaha and the Kawerau geothermal field), produced to clarify the nature of faulting and the relationship between seismicity and known/inferred geological structures in these fields. New 1-D velocity models and station corrections were derived and a subset of events relocated using a combination of catalogue and cross-correlation arrival times. It is distinct from the inventoried TVZ-wide high-precision 2007-2023 catalogue: an earlier period, field-specific focus, and different methods.
\par\vspace{2pt}\textbf{Key references:} \cite{Clarke2009}.
\catentryrule
\subsubsection*{\hypertarget{c:56}{56.}~ Rotokawa Geothermal Field Microseismicity Catalogue 2008-2012 (Sherburn et al. 2015)}
{\small\itshape induced $\cdot$ geographical $\cdot$ manual $\cdot$ reviewed}\par\vspace{2pt}
\begin{itemize}[leftmargin=1.3em,itemsep=0.6pt,topsep=1pt,parsep=0pt]
\kv{Also known as}{Sherburn et al. 2015 Rotokawa microseismicity; Microseismicity at Rotokawa geothermal field 2008-2012}
\kv{Coverage}{2008 to 2012}
\kv{Region}{Rotokawa geothermal field, central Taupo Volcanic Zone, North Island, New Zealand}
\kv{Bounding box}{Rotokawa geothermal field, central Taupo Volcanic Zone, North Island, New Zealand (\textasciitilde{}38.62S, 176.20E); most events confined to a \textasciitilde{}1 km\textasciicircum{}2 zone between reinjection and production wells.}
\kv{Depth range}{Shallow reservoir depths; \textgreater{}70\% of events concentrated \textasciitilde{}1.5-3 km depth}
\kv{Magnitude}{ML; \textgreater{}1000 events ML\textgreater{}=0.8, \textasciitilde{}50 events ML\textgreater{}=2, 2 events ML\textgreater{}=3 (largest M3.1). Stated completeness around ML\textasciitilde{}0.8.}
\kv{Events}{\textgreater{}1000 events of ML\textgreater{}=0.8 located mid-2008 to end-2012 (50 of ML\textgreater{}=2; 2 of ML\textgreater{}=3)}
\kv{Content}{hypocentral locations, phase arrival-time picks, quality information}
\kv{Producer / contact}{S. Sherburn (lead author, GNS Science) and co-authors S.M. Sewell, S. Bourguignon, W. Cumming, S. Bannister, C. Bardsley, J. Winick, J. Quinao, I.C. Wallis (GNS Science / Mighty River Power / Contact / consultants).}
\kv{Data source}{Events located by the local Rotokawa geothermal-field seismic network (operator monitoring), 2008-2012; located using local-network arrival-time picks. Interpreted in the context of deep fluid reinjection (e.g. Nga Awa Purua power station) and production records.}
\kv{Availability}{available (DOI) $\cdot$ DOI \href{https://doi.org/10.1016/j.geothermics.2014.11.001}{\nolinkurl{10.1016/j.geothermics.2014.11.001}}}
\end{itemize}
\par\vspace{2pt}A field-specific induced-microseismicity catalogue for the Rotokawa geothermal field (TVZ) covering mid-2008 to end-2012, interpreted as injection-induced via cooling-driven reservoir contraction (reinjected fluid \textasciitilde{}200 C cooler than the reservoir). It documents the spatial-temporal variation of microseismicity and its close correlation with deep reinjection rate. It is distinct from the inventoried Rotokawa matched-filter catalogue (Hopp et al. 2020, 2012-2015): a different time window and detection approach (local-network locations rather than 2012-2015 array template matching).
\par\vspace{2pt}\textbf{Key references:} \cite{Sherburn2015}.
\catentryrule
\subsubsection*{\hypertarget{c:57}{57.}~ Wairakei Geothermal Field Microseismicity Catalogue (Wairakei Seismic Network; Sepulveda et al. 2015)}
{\small\itshape induced $\cdot$ geographical $\cdot$ manual $\cdot$ partially reviewed}\par\vspace{2pt}
\begin{itemize}[leftmargin=1.3em,itemsep=0.6pt,topsep=1pt,parsep=0pt]
\kv{Also known as}{Wairakei Seismic Network (WSN) microseismicity; Spatial-temporal Characteristics of Microseismicity (2009-2014) of the Wairakei Geothermal Field (Sepulveda et al. 2015, WGC); Wairakei-Tauhara downhole microseismic catalogue}
\kv{Coverage}{2009-03 to ongoing}
\kv{Region}{Wairakei-Tauhara geothermal field, central Taupo Volcanic Zone, North Island, New Zealand}
\kv{Bounding box}{Wairakei-Tauhara geothermal field, central Taupo Volcanic Zone, North Island, New Zealand (\textasciitilde{}38.62S, 176.10E).}
\kv{Depth range}{Shallow crustal; \textasciitilde{}0 to \textasciitilde{}8.5 km, with 95\% of seismicity reported above \textasciitilde{}6.5 km and 99\% above \textasciitilde{}8.5 km, constraining the reservoir base}
\kv{Magnitude}{Microseismicity (low magnitude)}
\kv{Events}{5,649 located events (March 2009 - April 2013), per Sepulveda et al. (2015) and Sepulveda et al. (2013); spatial-temporal characterization extended to 2009-2014}
\kv{Content}{hypocentral locations, quality information}
\kv{Producer / contact}{Fabian R. Sepulveda (Contact Energy Ltd, Wairakei) and co-authors (J. Andrews/GNS Science; J. Kim; C. Siega; S.F. Milloy). Operator: Contact Energy Ltd.}
\kv{Data source}{Local Wairakei Seismic Network of 9-13 downhole/borehole seismometers (depths \textasciitilde{}65 m to \textasciitilde{}1400 m) operated by Contact Energy, commissioned from 2009; events located from local-network arrival times.}
\kv{Availability}{available on request}
\end{itemize}
\par\vspace{2pt}A dedicated operator-run downhole borehole microseismic monitoring catalogue for the Wairakei (Wairakei-Tauhara) geothermal field, used to constrain the deep field structure, reservoir base, fluid-circulation depth and the modern (microseismicity-constrained) field boundary at depth. It is not represented in the existing inventory. Recorded by a 9-13 station borehole network commissioned from early 2009.
\par\vspace{2pt}\textbf{Key references:} \cite{Sepulveda2015,Sepulveda2013}.
\catentryrule
\subsubsection*{\hypertarget{c:58}{58.}~ Mt Ruapehu 2022 Volcanic Unrest Drumbeat / Tremor Seismicity Catalogue (Bramwell et al. 2025)}
{\small\itshape background $\cdot$ particular-event-type $\cdot$ manual $\cdot$ reviewed}\par\vspace{2pt}
\begin{itemize}[leftmargin=1.3em,itemsep=0.6pt,topsep=1pt,parsep=0pt]
\kv{Also known as}{Drumbeat and low frequency catalogue for 2022 unrest at Ruapehu volcano (Zenodo 13910455); Bramwell et al. 2025 Ruapehu drumbeat/tremor catalogue}
\kv{Coverage}{2022-03 to 2022-07}
\kv{Region}{Mt Ruapehu, Tongariro National Park, southern Taupo Volcanic Zone, North Island, New Zealand}
\kv{Bounding box}{Mt Ruapehu crater/summit area, southern Taupo Volcanic Zone, North Island, New Zealand (\textasciitilde{}39.28S, 175.57E)}
\kv{Depth range}{Shallow volcanic / crater-lake hydrothermal system}
\kv{Magnitude}{Low-frequency drumbeat earthquakes and tremor (not characterized by standard ML magnitude/Mc); analysed in time, amplitude and frequency domains}
\kv{Events}{42,951 LF drumbeats; \textasciitilde{}144 h of overlapping drumbeats; \textgreater{}2000 h of tremor over the 121-day unrest period (7 March - 6 July 2022)}
\kv{Content}{other (low-frequency drumbeat earthquakes and volcanic tremor; source dynamics), quality information}
\kv{Producer / contact}{Liam A. Bramwell (lead author, Victoria University of Wellington); co-authors Finnigan Illsley-Kemp, Ery C. Hughes, Sophie Butcher, Oliver D. Lamb, Yannik Behr. Waveforms: GeoNet / GNS Science.}
\kv{Data source}{Continuous seismic waveforms from the GeoNet permanent network around Ruapehu (19 stations within 26 km; primary station MAVZ \textasciitilde{}1.45 km N of the crater lake). Events manually picked using the Pyrocko Snuffler toolbox. Full catalogue archived on Zenodo.}
\kv{Availability}{available (DOI) $\cdot$ DOI \href{https://doi.org/10.5281/zenodo.13910455}{\nolinkurl{10.5281/zenodo.13910455}}}
\end{itemize}
\par\vspace{2pt}A manually picked catalogue of low-frequency (LF) drumbeat earthquakes and volcanic tremor produced for Mt Ruapehu's 2022 volcanic unrest, used to characterize the source dynamics (magma ascent into the shallow system) alongside crater-lake temperature and gas data. It is a distinct event-type catalogue (LF drumbeats + tremor) not captured by the inventoried Mitchinson et al. 2024 Ruapehu ML swarm catalogue (which clusters GeoNet's standard volcano-tectonic catalogue, 1990-2023).
\par\vspace{2pt}\textbf{Key references:} \cite{Bramwell2025,Bramwell2024}.
\catentryrule
\subsubsection*{\hypertarget{c:59}{59.}~ Auckland Volcanic Field Matched-Filter Earthquake Catalogue (van Wijk et al. 2021, COVID-19 lockdown study)}
{\small\itshape background $\cdot$ geographical $\cdot$ template-matching $\cdot$ reviewed}\par\vspace{2pt}
\begin{itemize}[leftmargin=1.3em,itemsep=0.6pt,topsep=1pt,parsep=0pt]
\kv{Also known as}{AVF COVID-19 lockdown matched-filter catalogue; van Wijk et al. 2021 Auckland Volcanic Field seismicity; Auckland Volcanic Seismic Network (AVSN) detections}
\kv{Coverage}{2019-11 to 2020-06}
\kv{Region}{Auckland Volcanic Field, North Island, New Zealand}
\kv{Bounding box}{Auckland Volcanic Field, greater Auckland region, North Island, New Zealand (\textasciitilde{}36.9S, 174.8E)}
\kv{Depth range}{Shallow crustal}
\kv{Magnitude}{Small local earthquakes}
\kv{Events}{35 local earthquakes (1 Nov 2019 - 15 Jun 2020), expanded from 5 known GeoNet events (30 new detections)}
\kv{Content}{hypocentral locations, phase arrival-time picks}
\kv{Producer / contact}{Kasper van Wijk (lead author, University of Auckland); co-authors Calum J. Chamberlain (VUW), Thomas Lecocq (Royal Observatory of Belgium), Koen Van Noten (Royal Observatory of Belgium). Waveforms: GeoNet FDSN.}
\kv{Data source}{Continuous seismic data from the Auckland Volcanic Seismic Network (AVSN), available via the GeoNet FDSN service. 59 known-earthquake templates used in a matched-filtering (template matching) scheme with EQcorrscan to detect additional local events; starting GeoNet catalogue of 5 known local events.}
\kv{Availability}{available on request $\cdot$ DOI \href{https://doi.org/10.5194/se-12-363-2021}{\nolinkurl{10.5194/se-12-363-2021}}}
\end{itemize}
\par\vspace{2pt}A small but genuinely distinct local catalogue for the Auckland Volcanic Field (a monogenetic volcanic field not otherwise represented in the inventory). Template matched-filtering on the AVSN identified 35 local earthquakes for 1 Nov 2019-15 Jun 2020 (30 new detections beyond the 5 GeoNet events), exploiting reduced urban seismic noise during the COVID-19 lockdown. Produced to test detection capability and seismic monitoring of the AVF during the lockdown 'quieting'.
\par\vspace{2pt}\textbf{Key references:} \cite{vanWijk2021,vanWijk2020}.
\catentryrule
\subsubsection*{\hypertarget{c:60}{60.}~ Monowai Submarine Volcano Long-Range Hydroacoustic (T-phase) Activity Catalogue}
{\small\itshape induced $\cdot$ particular-event-type $\cdot$ other $\cdot$ reviewed}\par\vspace{2pt}
\begin{itemize}[leftmargin=1.3em,itemsep=0.6pt,topsep=1pt,parsep=0pt]
\kv{Also known as}{Metz et al. (2018) Monowai T-phase catalogue; Tracking Submarine Volcanic Activity at Monowai}
\kv{Coverage}{2003-07 to 2017-01}
\kv{Region}{Northern Kermadec Arc (Monowai volcanic centre), offshore New Zealand / Kermadec territory, SW Pacific}
\kv{Bounding box}{Point source: Monowai volcanic centre \textasciitilde{}25.9 deg S, 177.2 deg W (northern Kermadec Arc, SW Pacific). Detected via IMS hydrophone array near Juan Fernandez Islands.}
\kv{Magnitude}{Detection threshold reported as \textasciitilde{}2.2 mb in the Kermadec Arc region (about one order of magnitude better than regional land networks). The '1.4-4.2 mb' arrival range from the proposal brief was NOT corroborated in the sources accessed; treat as unverified.}
\kv{Events}{82 discrete activity episodes (Jul 2003-Mar 2004 and Apr 2014-Jan 2017)}
\kv{Content}{other (long-range hydroacoustic T-phase detections / activity episodes), quality information, amplitude picks}
\kv{Producer / contact}{D. Metz (University of Oxford); A. B. Watts (University of Oxford); I. Grevemeyer (GEOMAR Helmholtz Centre for Ocean Research Kiel)}
\kv{Data source}{Long-range hydroacoustic (T-phase) recordings from a CTBTO International Monitoring System (IMS) hydrophone array located near the Juan Fernandez Islands (SE Pacific, \textasciitilde{}17,000 km distant); direction-of-arrival and density-based spatial clustering analysis of 3.5 years of recordings to track Monowai volcanic centre.}
\kv{Availability}{available (DOI) $\cdot$ DOI \href{https://doi.org/10.1029/2018JB015888}{\nolinkurl{10.1029/2018JB015888}}}
\end{itemize}
\par\vspace{2pt}A hydroacoustic event catalogue tracking eruptive/seismic activity at the Monowai submarine volcanic centre, northern Kermadec Arc (within New Zealand's Kermadec territory). Produced to demonstrate that ultra-long-range SOFAR-channel T-phase recordings from distant IMS hydrophone arrays can resolve submarine volcanic activity roughly one order of magnitude smaller than land-based seismic networks. 82 discrete episodes of activity were identified across two recording windows. Distinct from the Brothers volcano local OBH T-wave catalogue (different volcano; remote IMS hydrophone array rather than local ocean-bottom instruments).
\par\vspace{2pt}\textbf{Key references:} \cite{Metz2018,Metz2016}.
\catentryrule
\subsubsection*{\hypertarget{c:61}{61.}~ Tongariro / Te Maari Relocated Microseismicity Catalogue 2010-2021 (Heise et al. 2024)}
{\small\itshape background $\cdot$ geographical $\cdot$ machine-learning $\cdot$ partially reviewed}\par\vspace{2pt}
\begin{itemize}[leftmargin=1.3em,itemsep=0.6pt,topsep=1pt,parsep=0pt]
\kv{Also known as}{Te Maari relocated seismicity; Mt Tongariro tomoDD relocated microseismicity 2010-2021}
\kv{Coverage}{2010 to 2021}
\kv{Region}{Tongariro/Ruapehu, central North Island (Taupo Volcanic Zone), New Zealand}
\kv{Bounding box}{Approx. Mt Tongariro / Te Maari area, central North Island, NZ (\textasciitilde{}39.1S, 175.6-175.7E); exact box not stated in profile}
\kv{Depth range}{Shallow crustal; focus on events shallower than \textasciitilde{}10 km (analysis spans \textasciitilde{}2.5-12 km)}
\kv{Magnitude}{No magnitude threshold applied; microseismicity (small magnitudes); explicit Mc/range not stated in the article}
\kv{Events}{\textasciitilde{}1435 relocated hypocentres (events with \textgreater{}5 phase observations)}
\kv{Content}{hypocentral locations, relative relocations, phase arrival-time picks, quality information}
\kv{Producer / contact}{Wiebke Heise, GNS Science (Earth Sciences New Zealand), Lower Hutt, NZ}
\kv{Data source}{Continuous GeoNet seismic data (Tongariro/Ruapehu network), processed by GNS Science}
\kv{Availability}{available on request $\cdot$ DOI \href{https://doi.org/10.21420/0ky8-mx13}{\nolinkurl{10.21420/0ky8-mx13}}}
\end{itemize}
\par\vspace{2pt}A relocated microseismicity catalogue for Mt Tongariro (Te Maari / Upper Te Maari area) covering 2010-2021, produced as part of a joint magnetotellurics + seismicity study of the 2012 Te Maari phreatic eruption sequence. The catalogue was built from \textasciitilde{}12 years of continuous GeoNet data using the EQTransformer machine-learning phase picker, initial absolute locations from NonLinLoc, and high-precision relative relocation with the double-difference tomoDD algorithm, yielding \textasciitilde{}1435 relocated hypocentres (events with \textgreater{}5 phase observations). Its purpose is to image the relationship between the shallow hydrothermal system and the underlying magmatic system that primed the phreatic eruptions.
\par\vspace{2pt}\textbf{Key references:} \cite{Heise2024}.
\catentryrule
\subsubsection*{\hypertarget{c:62}{62.}~ Auckland Volcanic Field Self-Consistent Machine-Learning Earthquake Catalogue (Chamberlain et al. 2025; NHC report 114 'Sensing unrest in New Zealand's largest city')}
{\small\itshape background $\cdot$ geographical $\cdot$ machine-learning $\cdot$ partially reviewed}\par\vspace{2pt}
\begin{itemize}[leftmargin=1.3em,itemsep=0.6pt,topsep=1pt,parsep=0pt]
\kv{Also known as}{Sensing unrest in New Zealand's largest city: detailed mapping of seismicity in Auckland; AVF machine-learning earthquake catalogue; O'Hagan (2025) AVF catalogue; NHC Toka Tu Ake report no. 114; Auckland Volcanic Field seismic catalogues (Zenodo 15825886)}
\kv{Coverage}{2011 to 2022}
\kv{Region}{Auckland Volcanic Field, Auckland, North Island, Aotearoa New Zealand (urban monogenetic volcanic field)}
\kv{Bounding box}{Region of interest approx. 174.4-175.3 E, 37.2-36.5 S (located events filtered to 174.0-175.5 E, 37.5-36.0 S; station retrieval bounds 174.5-175.25 E, 37.25-36.4 S)}
\kv{Depth range}{Approx. 0-50+ km; majority shallow (concentration near \textasciitilde{}5 km, likely a location artefact); 11 events with nominal negative depths (within uncertainty); deeper events down-weighted/required more picks. Average 68\% depth uncertainty 3.8 km, horizontal 3.7 km.}
\kv{Magnitude}{Local magnitude (MLNZ20 scale, Christophersen et al. 2022); range approx. -0.5 to 3.9; magnitude of completeness Mc \textasciitilde{}0.6 (maximum-curvature and goodness-of-fit); Gutenberg-Richter b-value \textasciitilde{}0.8. For comparison GeoNet reports 65 events (M 0.4-3.9) over the same region/period.}
\kv{Events}{368 earthquakes within the AVF region of interest (611 true earthquakes total before regional filtering); separate quarry-blast and unclassified-event catalogues also released; \textgreater{}5x the GeoNet count for the same region/period}
\kv{Content}{hypocentral locations; phase arrival-time picks; amplitude picks; quality information; other (also separate quarry-blast and unclassified-event catalogues; focal mechanisms recommended but NOT yet computed)}
\kv{Producer / contact}{Calum John Chamberlain (Victoria University of Wellington), calum.chamberlain@\allowbreak{}vuw.ac.nz; co-authors Jenni Hopkins, John Townend (VUW), Kasper van Wijk (University of Auckland), Nathanial Lindsey (FiberSense); dataset prepared by MSc student Sarah O'Hagan (VUW)}
\kv{Data source}{GeoNet seismic waveform archive (Auckland Volcano Seismic Network plus 5 GeoNet broadband sites within 200 km; \textasciitilde{}16 stations) downloaded from the GeoNet AWS public S3 bucket, 01/01/2011-01/01/2022. Re-detected/relocated independently of the routine GeoNet catalogue. Funded by Natural Hazards Commission Toka Tu Ake biennial research fund; Chamberlain supported by a Rutherford Discovery Fellowship.}
\kv{Availability}{available (DOI) $\cdot$ DOI \href{https://doi.org/10.5281/zenodo.15825886}{\nolinkurl{10.5281/zenodo.15825886}}}
\end{itemize}
\par\vspace{2pt}First long-duration, self-consistent located-earthquake catalogue for the Auckland Volcanic Field (AVF), built to characterise background seismicity so that future deviations (possible volcanic unrest) can be recognised. A machine-learning-enhanced, semi-automated workflow (EQTransformer picking via seisbench -\textgreater{} PyOcto phase association -\textgreater{} NonLinLoc location with the NZ3D velocity model -\textgreater{} manual quality control -\textgreater{} MLNZ20 local magnitudes) was applied to \textasciitilde{}11 years of GeoNet waveforms, detecting and locating more than 5x as many earthquakes as the GeoNet catalogue. The final catalogue contains 368 earthquakes within the AVF region of interest (611 true earthquakes from 19,800 associated events before regional filtering), and resolves nine sequences of \textgreater{}5 events occurring without volcanic unrest. The report also tested DAS earthquake detection on a Wellington telecom fibre, but only for detection-method development; the AVF event list is the in-scope located-earthquake product.
\par\vspace{2pt}\textbf{Key references:} \cite{Chamberlain2025,OHagan2025,OHagan2025a}.
\catentryrule
\subsection{Slow-slip, tremor, LFE and moment-tensor catalogues}
\textit{Catalogues of the subduction slow-earthquake family (slow-slip events, tectonic tremor, low-frequency and repeating earthquakes) together with regional moment-tensor solutions.}\par\vspace{4pt}
\subsubsection*{\hypertarget{c:63}{63.}~ GeoNet New Zealand Earthquake Moment Tensor Solutions (MwNZ)}
{\small\itshape background $\cdot$ geographical $\cdot$ automatic $\cdot$ partially reviewed}\par\vspace{2pt}
\begin{itemize}[leftmargin=1.3em,itemsep=0.6pt,topsep=1pt,parsep=0pt]
\kv{Also known as}{MwNZ; GeoNet CMT; regional moment tensor (RMT) catalogue; Ristau RMT catalogue}
\kv{Coverage}{2003-08 to ongoing}
\kv{Region}{New Zealand and adjacent offshore regions}
\kv{Bounding box}{New Zealand and adjacent offshore regions (approx. NZ landmass plus Hikurangi/Kermadec and Puysegur margins)}
\kv{Magnitude}{Approximately Mw 3.0-8.0; solutions routinely computed for moderate (M \textgreater{} \textasciitilde{}4) earthquakes}
\kv{Events}{\textasciitilde{}1900 (grows monthly; exact count not stated on landing pages)}
\kv{Content}{focal mechanisms; hypocentral locations; quality information}
\kv{Producer / contact}{GNS Science / GeoNet (John Ristau, GNS Science)}
\kv{Data source}{Regional broadband seismic waveform data from the GeoNet national network (and adjacent stations), inverted for regional centroid moment tensor solutions; routine production since August 2003.}
\kv{Availability}{available (DOI) $\cdot$ DOI \href{https://doi.org/10.21420/MMJ9-CZ67}{\nolinkurl{10.21420/MMJ9-CZ67}}}
\end{itemize}
\par\vspace{2pt}National regional centroid moment tensor (focal mechanism) catalogue for New Zealand and adjacent offshore earthquakes, produced routinely since August 2003 to provide consistent Mw and fault-geometry information for moderate-to-large NZ earthquakes. It established the MwNZ regional moment magnitude scale used across NZ seismology and hazard work. Two inversion approaches have been used over time (Dreger/Berkeley code 2003-2020; a Penn State inversion code from June 2020).
\par\vspace{2pt}\textbf{Key references:} \cite{GNSScience2006,Ristau2008,Ristau2009}.
\catentryrule
\subsubsection*{\hypertarget{c:64}{64.}~ Hikurangi Slow Slip Event Inventory (Wallace \& Beavan 2010)}
{\small\itshape compilation $\cdot$ particular-event-type $\cdot$ other $\cdot$ reviewed}\par\vspace{2pt}
\begin{itemize}[leftmargin=1.3em,itemsep=0.6pt,topsep=1pt,parsep=0pt]
\kv{Also known as}{Hikurangi SSE catalogue; Hikurangi slow-slip inventory; Wallace \& Beavan 2010}
\kv{Coverage}{2002 to 2010}
\kv{Region}{Hikurangi margin}
\kv{Bounding box}{North Island, New Zealand / Hikurangi subduction margin (eastern North Island)}
\kv{Depth range}{shallow (north/central) to \textgreater{}25-30 km (deep southern SSEs)}
\kv{Magnitude}{Equivalent moment magnitudes approximately Mw 6.3-7.2}
\kv{Events}{\textasciitilde{}15 slow slip events}
\kv{Content}{fault-geometry; other}
\kv{Producer / contact}{Laura M. Wallace and John Beavan, GNS Science}
\kv{Data source}{Continuous GPS (cGPS) time series from the GeoNet continuous GPS network across the North Island / Hikurangi margin; elastic dislocation modelling of surface displacements.}
\kv{Availability}{available on request $\cdot$ DOI \href{https://doi.org/10.1029/2010JB007717}{\nolinkurl{10.1029/2010JB007717}}}
\end{itemize}
\par\vspace{2pt}Foundational inventory of Hikurangi slow slip events documented with continuous GPS, revealing marked diversity in SSE behaviour: shallow, frequent events in the north and central margin versus deep (\textgreater{}25-30 km), roughly 5-yearly events in the south. It underpins nearly all subsequent NZ SSE studies and characterised durations (days to \textgreater{}1 year), equivalent moment release (Mw \textasciitilde{}6.3-7.2) and recurrence intervals (2 to \textgreater{}5 years).
\par\vspace{2pt}\textbf{Key references:} \cite{Wallace2010}.
\catentryrule
\subsubsection*{\hypertarget{c:65}{65.}~ Northern Hikurangi Tectonic Tremor Catalogue 2010-2015 (Todd \& Schwartz 2016)}
{\small\itshape compilation $\cdot$ particular-event-type $\cdot$ other $\cdot$ reviewed}\par\vspace{2pt}
\begin{itemize}[leftmargin=1.3em,itemsep=0.6pt,topsep=1pt,parsep=0pt]
\kv{Also known as}{Todd \& Schwartz 2016 tremor catalogue; northern Hikurangi tremor catalogue}
\kv{Coverage}{2010 to 2015}
\kv{Region}{Hikurangi margin (Raukumara Peninsula)}
\kv{Bounding box}{Northern Hikurangi margin / Raukumara Peninsula, North Island, New Zealand}
\kv{Content}{other}
\kv{Producer / contact}{Erin K. Todd and Susan Y. Schwartz (University of California, Santa Cruz)}
\kv{Data source}{Densified onshore GeoNet seismic network (northern Hikurangi / Raukumara Peninsula); envelope cross-correlation of continuous seismic data to detect and locate tectonic tremor.}
\kv{Availability}{available on request $\cdot$ DOI \href{https://doi.org/10.1002/2016JB013480}{\nolinkurl{10.1002/2016JB013480}}}
\end{itemize}
\par\vspace{2pt}First systematic tectonic tremor catalogue for the northern Hikurangi margin (2010-2015), associating tremor with geodetic transients off the Raukumara Peninsula and helping identify slow slip events. It used envelope cross-correlation on the densified onshore GeoNet network. This is a tremor catalogue rather than a hypocentral earthquake catalogue.
\par\vspace{2pt}\textbf{Key references:} \cite{Todd2016}.
\catentryrule
\subsubsection*{\hypertarget{c:66}{66.}~ Ambient Tectonic Tremor Catalogue: Manawatu, Cape Turnagain, Marlborough \& Puysegur (Romanet \& Ide 2019)}
{\small\itshape compilation $\cdot$ particular-event-type $\cdot$ other $\cdot$ reviewed}\par\vspace{2pt}
\begin{itemize}[leftmargin=1.3em,itemsep=0.6pt,topsep=1pt,parsep=0pt]
\kv{Also known as}{Romanet \& Ide 2019 ambient tremor catalogue}
\kv{Coverage}{2005 to 2016}
\kv{Region}{Manawatu, Cape Turnagain, Marlborough, Puysegur (Fiordland)}
\kv{Bounding box}{Four NZ regions: Manawatu, Cape Turnagain (lower North Island), Marlborough (northern South Island), Puysegur/Fiordland (SW South Island)}
\kv{Content}{other}
\kv{Producer / contact}{Pierre Romanet and Satoshi Ide (University of Tokyo)}
\kv{Data source}{GeoNet continuous seismic data across four NZ regions; tremor detection using duration-energy scaling combined with visual inspection.}
\kv{Availability}{available on request $\cdot$ DOI \href{https://doi.org/10.1186/s40623-019-1039-1}{\nolinkurl{10.1186/s40623-019-1039-1}}}
\end{itemize}
\par\vspace{2pt}Multi-region ambient tectonic tremor catalogue across four New Zealand areas (Manawatu, Cape Turnagain, Marlborough and Puysegur), notably documenting tremor in the Puysegur/Fiordland subduction system in addition to the Hikurangi-related regions. Tremor was detected using a duration-energy scaling approach plus visual inspection of GeoNet data over roughly 2005-2016. It is a tremor catalogue rather than a hypocentral earthquake catalogue.
\par{\small\itshape Related: The Manawatu portion overlaps the deep slow-slip-and-tremor catalogue of Fasola et al. (2023); this catalogue spans several regions (Manawatu, Cape Turnagain, Marlborough, Puysegur).}
\par\vspace{2pt}\textbf{Key references:} \cite{Romanet2019}.
\catentryrule
\subsubsection*{\hypertarget{c:67}{67.}~ Alpine Fault Low-Frequency Earthquake Catalogue (Chamberlain et al. 2014)}
{\small\itshape background $\cdot$ particular-event-type $\cdot$ template-matching $\cdot$ reviewed}\par\vspace{2pt}
\begin{itemize}[leftmargin=1.3em,itemsep=0.6pt,topsep=1pt,parsep=0pt]
\kv{Also known as}{Alpine Fault LFE catalogue 2014; Chamberlain et al. 2014}
\kv{Coverage}{2009 to 2012}
\kv{Region}{Alpine Fault / Southern Alps}
\kv{Bounding box}{Central Alpine Fault, Southern Alps, South Island, New Zealand}
\kv{Depth range}{\textasciitilde{}20-30 km (deep extent of the Alpine Fault)}
\kv{Events}{\textasciitilde{}8,760 LFEs in 14 families}
\kv{Content}{hypocentral locations; phase arrival-time picks; relative relocations; other}
\kv{Producer / contact}{Calum J. Chamberlain (with David R. Shelly, John Townend, Tim A. Stern), Victoria University of Wellington / USGS}
\kv{Data source}{GeoNet plus SAMBA (Southern Alps Microearthquake Borehole Array) continuous seismic data; matched-filter (template-matching) detection of low-frequency earthquakes within tremor.}
\kv{Availability}{available on request $\cdot$ DOI \href{https://doi.org/10.1002/2014GC005436}{\nolinkurl{10.1002/2014GC005436}}}
\end{itemize}
\par\vspace{2pt}First low-frequency earthquake (LFE) catalogue for the deep Alpine Fault: approximately 8,760 LFEs in 14 families over 36 months (2009-2012), revealing punctuated slow slip at roughly 20-30 km depth. It was built from GeoNet and SAMBA continuous data using matched-filter detection within tremor and is the earliest NZ LFE catalogue, serving as a precursor to Baratin et al. (2018).
\par{\small\itshape Related: Extended by Baratin et al. (2018), which adds focal mechanisms and additional years to these LFE families.}
\par\vspace{2pt}\textbf{Key references:} \cite{Chamberlain2014}.
\catentryrule
\subsubsection*{\hypertarget{c:68}{68.}~ Alpine Fault Low-Frequency Earthquake Catalogue with Focal Mechanisms (Baratin et al. 2018)}
{\small\itshape background $\cdot$ particular-event-type $\cdot$ template-matching $\cdot$ reviewed}\par\vspace{2pt}
\begin{itemize}[leftmargin=1.3em,itemsep=0.6pt,topsep=1pt,parsep=0pt]
\kv{Also known as}{Baratin et al. 2018 Alpine Fault LFE catalogue}
\kv{Coverage}{2009-03 to 2016}
\kv{Region}{Alpine Fault / Southern Alps}
\kv{Bounding box}{Central Alpine Fault, Southern Alps, South Island, New Zealand}
\kv{Depth range}{\textasciitilde{}20-30 km (deep extent of the Alpine Fault)}
\kv{Events}{\textasciitilde{}10,000 LFEs in 14 families (first Alpine Fault LFE focal mechanisms)}
\kv{Content}{hypocentral locations; phase arrival-time picks; relative relocations; focal mechanisms; other}
\kv{Producer / contact}{Laura-May Baratin (with Calum J. Chamberlain, John Townend, Martha K. Savage), Victoria University of Wellington}
\kv{Data source}{GeoNet and SAMBA continuous seismic data; matched-filter (template-matching) detection of LFEs, extended templates relative to Chamberlain et al. (2014), with LFE focal mechanism determination.}
\kv{Availability}{available on request $\cdot$ DOI \href{https://doi.org/10.1016/j.epsl.2017.12.021}{\nolinkurl{10.1016/j.epsl.2017.12.021}}}
\end{itemize}
\par\vspace{2pt}Extended, longer-duration deep Alpine Fault LFE catalogue (2009-2017) that adds the first LFE focal mechanisms for the region. It shows that increased LFE rates coincide with tremor and regional earthquakes, interpreted as quasi-continuous deformation and triggered slow slip on the deep Alpine Fault. It builds on and extends the Chamberlain et al. (2014) template set.
\par{\small\itshape Related: Builds on and extends the Alpine Fault LFE catalogue of Chamberlain et al. (2014), adding focal mechanisms.}
\par\vspace{2pt}\textbf{Key references:} \cite{Baratin2018}.
\catentryrule
\subsubsection*{\hypertarget{c:69}{69.}~ Kaimanawa (Deep Hikurangi) Low-Frequency Earthquake Catalogue (Aden-Antoniow et al. 2024)}
{\small\itshape background $\cdot$ particular-event-type $\cdot$ template-matching $\cdot$ reviewed}\par\vspace{2pt}
\begin{itemize}[leftmargin=1.3em,itemsep=0.6pt,topsep=1pt,parsep=0pt]
\kv{Also known as}{Aden-Antoniow et al. 2024 Hikurangi LFE catalogue; Kaimanawa LFE catalogue}
\kv{Coverage}{2010 to 2022}
\kv{Region}{Hikurangi margin (Kaimanawa Range)}
\kv{Bounding box}{Beneath the Kaimanawa Range, central North Island, New Zealand (Hikurangi subduction zone)}
\kv{Depth range}{\textasciitilde{}50 km}
\kv{Events}{\textasciitilde{}21,000 events from 36 LFE sources}
\kv{Content}{hypocentral locations; relative relocations; other}
\kv{Producer / contact}{Florent Aden-Antoniow (with William B. Frank, Calum J. Chamberlain, John Townend, Laura M. Wallace, Stephen Bannister)}
\kv{Data source}{GeoNet regional continuous seismic data; two-iteration matched-filter (template-matching) detection using templates derived from prior tectonic tremor observations.}
\kv{Availability}{available on request $\cdot$ DOI \href{https://doi.org/10.1029/2023JB027971}{\nolinkurl{10.1029/2023JB027971}}}
\end{itemize}
\par\vspace{2pt}First low-frequency earthquake catalogue for the Hikurangi subduction zone proper, beneath the Kaimanawa Range at about 50 km depth, downdip of regularly recurring deep \textasciitilde{}M7 slow slip events. It comprises 36 LFE sources producing almost 21,000 events over more than a decade (2010-2022), detected with a two-iteration matched-filter on GeoNet data. It is distinct from the Alpine Fault LFE catalogues.
\par\vspace{2pt}\textbf{Key references:} \cite{AdenAntoniow2024}.
\catentryrule
\subsubsection*{\hypertarget{c:70}{70.}~ Raukumara Peninsula Repeating Earthquake Catalogue 2003-2020 (Hughes et al. 2021)}
{\small\itshape background $\cdot$ particular-event-type $\cdot$ template-matching $\cdot$ reviewed}\par\vspace{2pt}
\begin{itemize}[leftmargin=1.3em,itemsep=0.6pt,topsep=1pt,parsep=0pt]
\kv{Also known as}{Hughes et al. 2021 Raukumara repeating-earthquake catalogue}
\kv{Coverage}{2003 to 2020}
\kv{Region}{Hikurangi margin (Raukumara Peninsula)}
\kv{Bounding box}{Raukumara Peninsula, northern Hikurangi subduction margin, North Island, New Zealand}
\kv{Events}{61 repeating-earthquake families}
\kv{Content}{hypocentral locations; relative relocations; focal mechanisms; other}
\kv{Producer / contact}{Laura Hughes (with Calum J. Chamberlain, John Townend, Amanda M. Thomas), Victoria University of Wellington}
\kv{Data source}{GeoNet continuous seismic data, northern Hikurangi / Raukumara Peninsula; waveform cross-correlation to identify repeating-earthquake families.}
\kv{Availability}{available on request $\cdot$ DOI \href{https://doi.org/10.1029/2021GC009670}{\nolinkurl{10.1029/2021GC009670}}}
\end{itemize}
\par\vspace{2pt}Catalogue of 61 repeating-earthquake families on the northern Hikurangi margin (Raukumara Peninsula, 2003-2020) used as an interface 'creepmeter'. The families coincide spatially with slow-slip patches and include both interplate and intraplate focal mechanisms. It was built from GeoNet continuous data using waveform cross-correlation.
\par\vspace{2pt}\textbf{Key references:} \cite{Hughes2021}.
\catentryrule
\subsubsection*{\hypertarget{c:71}{71.}~ New Zealand-Wide Earthquake Focal Mechanism Catalogue for Tectonic Stress (Townend et al. 2012)}
{\small\itshape compilation $\cdot$ geographical $\cdot$ other $\cdot$ reviewed}\par\vspace{2pt}
\begin{itemize}[leftmargin=1.3em,itemsep=0.6pt,topsep=1pt,parsep=0pt]
\kv{Also known as}{Townend et al. (2012) NZ focal mechanism compilation; central NZ stress-field focal mechanism dataset}
\kv{Coverage}{2004-01 to 2011-02}
\kv{Region}{Central New Zealand, Hikurangi subduction margin (North Island, including Hawke's Bay area) and South Island, along the Australia-Pacific plate boundary}
\kv{Events}{\textasciitilde{}3424 focal mechanisms}
\kv{Content}{focal mechanisms; hypocentral locations; quality information}
\kv{Producer / contact}{John Townend (lead author), Victoria University of Wellington / GNS Science; co-authors Steven Sherburn (GNS Science), Richard Arnold (VUW), Carolin Boese, Louise Woods}
\kv{Data source}{Compiled first-motion / polarity focal mechanisms derived from earthquakes recorded by the GeoNet national seismograph network (and related temporary deployments); a purpose-built NZ-wide focal-mechanism catalogue assembled for stress-tensor inversion. Distinct from the GeoNet moment-tensor (MwNZ / Ristau RMT) product.}
\kv{Availability}{available on request $\cdot$ DOI \href{https://doi.org/10.1016/j.epsl.2012.08.003}{\nolinkurl{10.1016/j.epsl.2012.08.003}}}
\end{itemize}
\par\vspace{2pt}A New Zealand-wide compilation of \textasciitilde{}3424 earthquake focal mechanisms assembled to estimate the present-day tectonic stress field at \textasciitilde{}100 locations across central New Zealand and map the orientation of maximum horizontal compressive stress (SHmax) along the Australia-Pacific plate boundary. Produced to characterise three-dimensional variations in stress along the Hikurangi subduction margin and South Island, and to relate stress orientation to along-strike variations in subduction-thrust coupling. It was the largest study of the NZ plate-boundary stress field at the time of publication.
\par\vspace{2pt}\textbf{Key references:} \cite{Townend2012}.
\catentryrule
\subsubsection*{\hypertarget{c:72}{72.}~ Southeastern South Island / Otago Earthquake Focal Mechanism Catalogue (Williams et al. 2026)}
{\small\itshape compilation $\cdot$ geographical $\cdot$ other $\cdot$ reviewed}\par\vspace{2pt}
\begin{itemize}[leftmargin=1.3em,itemsep=0.6pt,topsep=1pt,parsep=0pt]
\kv{Also known as}{Williams et al. (2026) southeastern South Island focal mechanisms; Otago/Southland focal-mechanism catalogue (OtagoNet + SOSA + COSA)}
\kv{Coverage}{2001 to 2020}
\kv{Region}{Southeastern South Island, Aotearoa New Zealand, Otago and Southland}
\kv{Magnitude}{MLV 1.3-4.3}
\kv{Events}{126 combined focal mechanisms (91 rated C-quality or better, of which 57 strike-slip); constituent solutions: 96 OtagoNet + 30 SOSA + 11 COSA (167 array solutions plus GeoNet RMT entries)}
\kv{Content}{focal mechanisms; hypocentral locations; quality information}
\kv{Producer / contact}{Jack Williams (lead author); co-authors Donna Eberhart-Phillips, Sandra Bourguignon, Mark Stirling, Martin Reyners, Phaedra Upton (University of Otago / GNS Science)}
\kv{Data source}{First-motion / polarity focal mechanisms computed from earthquakes recorded by temporary regional networks: OtagoNet (Otago temporary broadband network, Reyners 2014), the Southland Otago Seismic Array (SOSA, Williams et al. 2022), and the Central Otago Seismic Array (COSA); supplemented by GeoNet regional moment-tensor (RMT) catalogue entries.}
\kv{Availability}{available (DOI) $\cdot$ DOI \href{https://doi.org/10.5281/zenodo.17228270}{\nolinkurl{10.5281/zenodo.17228270}}}
\end{itemize}
\par\vspace{2pt}A regional focal-mechanism catalogue for the southeastern South Island (Otago/Southland) of Aotearoa New Zealand, compiling 126 combined nodal-plane solutions for small earthquakes (MLV 1.3-4.3) to study how oblique Australia-Pacific plate motion is accommodated. The mechanisms reveal a strike-slip stress state with a WNW-trending maximum compressive stress and are used to argue for scale-dependent partitioning of transpressional strain controlled by inherited structures. Data are deposited on Zenodo under CC-BY-4.0.
\par\vspace{2pt}\textbf{Key references:} \cite{Williams2026a,Williams2025a}.
\catentryrule
\subsubsection*{\hypertarget{c:73}{73.}~ HOBITSS Ocean-Bottom Focal-Mechanism / Stress-Tensor Cycling Catalogue, Northern Hikurangi (Warren-Smith et al. 2019)}
{\small\itshape background $\cdot$ particular-event-type $\cdot$ other $\cdot$ reviewed}\par\vspace{2pt}
\begin{itemize}[leftmargin=1.3em,itemsep=0.6pt,topsep=1pt,parsep=0pt]
\kv{Also known as}{Warren-Smith et al. (2019) HOBITSS focal mechanisms; northern Hikurangi intraslab focal-mechanism / stress-ratio cycling dataset}
\kv{Coverage}{2014-05 to 2015-06}
\kv{Region}{Northern Hikurangi subduction margin, offshore Gisborne, North Island, New Zealand}
\kv{Depth range}{\textasciitilde{}13-15 km thick intraslab seismicity zone within the subducting Pacific oceanic crust; subduction interface seismogenic depths roughly 15-35 km (per-event depth range not tabulated in the paper)}
\kv{Magnitude}{Earthquakes with Mw \textgreater{} 1.8 (lower magnitude threshold for mechanism calculation); upper range and Mc not separately documented}
\kv{Events}{278 intraslab earthquakes analysed (subset of Mw\textgreater{}1.8 GeoNet events); \textasciitilde{}140 focal mechanisms at SSE initiation and \textasciitilde{}80 at SSE shutdown used for stress inversions}
\kv{Content}{focal mechanisms; hypocentral locations; quality information; other}
\kv{Producer / contact}{Emily Warren-Smith (lead author), GNS Science, Lower Hutt; co-authors B. Fry, L. Wallace, E. Chon, S. Henrys, A. Sheehan, K. Mochizuki, S. Schwartz, S. Webb, S. Lebedev}
\kv{Data source}{Earthquake first-motion polarity focal mechanisms computed for Mw\textgreater{}1.8 events from the GeoNet national seismicity catalogue, using the 2014-2015 HOBITSS (Hikurangi Ocean Bottom Investigation of Tremor and Slow Slip) ocean-bottom seismometer network (10 broadband + 5 short-period OBS, plus 24 absolute pressure gauges) combined with onshore GeoNet stations, offshore Gisborne.}
\kv{Availability}{available on request $\cdot$ DOI \href{https://doi.org/10.1038/s41561-019-0367-x}{\nolinkurl{10.1038/s41561-019-0367-x}}}
\end{itemize}
\par\vspace{2pt}A focal-mechanism catalogue for microearthquakes in the subducting Pacific oceanic crust at the northern Hikurangi margin offshore Gisborne, derived from the 2014-2015 HOBITSS ocean-bottom seismometer deployment. Used to perform time-varying stress-tensor inversions and track the stress shape ratio (R) as a proxy for fluid pressure cycling before, during, and after four shallow slow slip events. The product is distinct in content (focal mechanisms plus a stress-ratio time series) from the HOBITSS hypocentral microearthquake catalogues (Yarce et al. 2019, 2023).
\par{\small\itshape Related: Derived from the same 2014-2015 HOBITSS deployment as the manual (Yarce et al. 2019) and automatic-detection (Yarce et al. 2023) earthquake catalogues; this product provides focal mechanisms and stress-tensor estimates.}
\par\vspace{2pt}\textbf{Key references:} \cite{WarrenSmith2019}.
\catentryrule
\subsubsection*{\hypertarget{c:74}{74.}~ Hikurangi Trench Earthquake-Swarm Catalogue (ETAS-based, Nishikawa et al. 2021)}
{\small\itshape compilation $\cdot$ particular-event-type $\cdot$ other $\cdot$ non-reviewed}\par\vspace{2pt}
\begin{itemize}[leftmargin=1.3em,itemsep=0.6pt,topsep=1pt,parsep=0pt]
\kv{Also known as}{Nishikawa Hikurangi swarm catalog; ETAS swarm catalogue, Hikurangi Trench (1997-2015)}
\kv{Coverage}{1997 to 2015}
\kv{Region}{Hikurangi Trench / Hikurangi subduction margin, North Island, New Zealand (offshore eastern margin)}
\kv{Magnitude}{M\textgreater{}=3; magnitude of completeness (Mc) estimated to vary in time: 3.0 (1997-2001), 2.7 (2002-2006), 2.5 (2007-2011), 3.1 (2012-2015), attributed to changes in the GeoNet location algorithm/procedure in 2012}
\kv{Content}{hypocentral locations; quality information; other (statistical swarm classification / background-rate anomalies)}
\kv{Producer / contact}{Tomoaki Nishikawa (lead author); at time of publication Graduate School of Science / Disaster Prevention Research Institute, Kyoto University, Japan. Co-authors Takuya Nishimura and Yutaro Okada.}
\kv{Data source}{Derived by applying the epidemic-type aftershock-sequence (ETAS) model to the GeoNet New Zealand earthquake catalogue (M\textgreater{}=3, 1997-2015) to statistically separate background/swarm seismicity from aftershock sequences along the Hikurangi Trench.}
\kv{Availability}{DOI \href{https://doi.org/10.1029/2020JB020618}{\nolinkurl{10.1029/2020JB020618}}}
\end{itemize}
\par\vspace{2pt}A derived, declustering-type catalogue produced to study the spatiotemporal relationship between seismicity and slow slip events (SSEs) along the Hikurangi subduction margin. The authors fitted the ETAS model to GeoNet data to identify anomalous earthquake swarms (statistically elevated background rate) and built a new swarm catalogue, finding many swarms occur as intraplate events lagging SSEs by days or more. The focus is on isolating SSE-associated swarm activity rather than cataloguing all seismicity.
\par\vspace{2pt}\textbf{Key references:} \cite{Nishikawa2021}.
\catentryrule
\subsubsection*{\hypertarget{c:75}{75.}~ Manawatu Deep Short-Term Slow-Slip and Tremor Catalogue (Fasola et al. 2023)}
{\small\itshape background $\cdot$ particular-event-type $\cdot$ other $\cdot$ reviewed}\par\vspace{2pt}
\begin{itemize}[leftmargin=1.3em,itemsep=0.6pt,topsep=1pt,parsep=0pt]
\kv{Also known as}{Fasola et al. (2023) Manawatu deep tremor / short-term ETS-like SSE; Manawatu/Kaimanawa deep tremor-associated slip model}
\kv{Coverage}{2005 to 2016}
\kv{Region}{Manawatu and Kaimanawa region, deep central Hikurangi subduction zone, lower North Island, Aotearoa New Zealand}
\kv{Events}{\textasciitilde{}354 tremor events (input Romanet \& Ide 2019 Manawatu catalogue), grouped into 26 tremor clusters (durations 1-80 days, \textasciitilde{}75\% of tremor); no separate event count published for the derived slip catalogue}
\kv{Content}{hypocentral locations, fault-geometry, other}
\kv{Producer / contact}{Shannon L. Fasola (lead author; Univ. of Kansas / KU Crustal Deformation group), co-authors Noel M. Jackson (formerly Noel M. Bartlow, Univ. of Kansas) and Charles A. Williams (GNS Science, New Zealand)}
\kv{Data source}{Tremor locations/times reused from the Romanet \& Ide (2019) ambient-tremor catalogue (envelope cross-correlation on GeoNet permanent seismometers; \textasciitilde{}354 Manawatu tremor events, 2005-2016). Continuous GNSS daily time-series from the GeoNet Aotearoa NZ Continuous GNSS Network (GNS Science). New products derived in this paper: time-averaged tremor-associated surface displacements/velocities and a static slip inversion (tremor-associated slip distribution) on the Hikurangi plate interface, plus identification of deep short-term ETS-like slow slip.}
\kv{Availability}{available on request $\cdot$ DOI \href{https://doi.org/10.1029/2023GL105428}{\nolinkurl{10.1029/2023GL105428}}}
\end{itemize}
\par\vspace{2pt}Investigates whether small, short-term ETS-like slow-slip events occur at the deep Manawatu/Kaimanawa tremor source, where tectonic tremor lies adjacent to (not co-located with) the deep long-term Manawatu and Kaimanawa slow-slip events on the central Hikurangi subduction interface. The authors cluster existing tremor, decompose GNSS at tremor times to recover tremor-associated displacement offsets, and invert these for slip on the plate interface, concluding the interface below deep long-term SSEs likely slips frequently in small ETS-like SSEs too small to detect individually with geodesy. Distinct in focus (deep Manawatu tremor-associated slip modeling) from the northern Hikurangi tremor catalogue (Todd \& Schwartz 2016) and the source ambient-tremor catalogue (Romanet \& Ide 2019).
\par{\small\itshape Related: The Manawatu tremor here also features in the multi-region ambient-tremor catalogue of Romanet \& Ide (2019); this entry focuses on deep short-term slow slip and tremor.}
\par\vspace{2pt}\textbf{Key references:} \cite{Fasola2023,Romanet2019}.
\catentryrule
\subsubsection*{\hypertarget{c:76}{76.}~ Global Centroid Moment Tensor (GCMT) catalogue (New Zealand component)}
{\small\itshape compilation $\cdot$ particular-event-type $\cdot$ other $\cdot$ reviewed}\par\vspace{2pt}
\begin{itemize}[leftmargin=1.3em,itemsep=0.6pt,topsep=1pt,parsep=0pt]
\kv{Also known as}{Global CMT; GCMT; formerly Harvard CMT (HRV) catalog}
\kv{Coverage}{1976 to ongoing}
\kv{Region}{Global; New Zealand component covers onshore NZ, Hikurangi-Kermadec subduction margin, and subantarctic region}
\kv{Bounding box}{Global (NZ component covers approx. -50 to -33 lat, 165 to -176 lon)}
\kv{Depth range}{0 to \textasciitilde{}700 km (centroid depths; global product spanning shallow crustal to deep subduction)}
\kv{Magnitude}{Mw \textgreater{}= \textasciitilde{}5.0 globally (systematic threshold); the NZ component contains moderate-to-large events at or above this threshold}
\kv{Content}{focal mechanisms, hypocentral locations, other}
\kv{Producer / contact}{Global CMT Project, Lamont-Doherty Earth Observatory (LDEO), Columbia University (PI: Göran Ekström; Co-PI: Meredith Nettles). Originally founded at Harvard University (1982); relocated to LDEO in 2006. NSF-funded.}
\kv{Data source}{Long-period and broadband digital seismic waveform data from global networks (e.g. GSN/FDSN), inverted independently at LDEO/Columbia to derive centroid moment tensor solutions. Not a regional NZ network product; the NZ component is the subset of GCMT solutions located within/around New Zealand.}
\kv{Availability}{available (DOI) $\cdot$ DOI \href{https://doi.org/10.1016/j.pepi.2012.04.002}{\nolinkurl{10.1016/j.pepi.2012.04.002}}}
\end{itemize}
\par\vspace{2pt}The GCMT (formerly Harvard CMT) Project systematically determines centroid moment tensors and focal mechanisms for global earthquakes of M\textgreater{}=\textasciitilde{}5.0 since 1976, with more than 25,000 solutions. Its NZ component provides an independent, long-standing record of focal mechanisms and centroid locations for moderate-to-large events along the Hikurangi-Kermadec margin, crustal NZ, and the subantarctic region. It is distinct from the regional GeoNet/Ristau MwNZ moment-tensor catalogue, which is NZ-operated and extends to smaller events.
\par\vspace{2pt}\textbf{Key references:} \cite{Ekstrm2012,Dziewonski1981}.
\catentryrule
\subsection{Derived, homogenised and synthetic / hazard-model catalogues}
\textit{Magnitude-homogenised, declustered, and physics-based synthetic catalogues prepared for the National Seismic Hazard Model.}\par\vspace{4pt}
\subsubsection*{\hypertarget{c:77}{77.}~ NSHM 2010 Earthquake Catalogue (Stirling et al. 2012)}
{\small\itshape background $\cdot$ geographical $\cdot$ other $\cdot$ reviewed}\par\vspace{2pt}
\begin{itemize}[leftmargin=1.3em,itemsep=0.6pt,topsep=1pt,parsep=0pt]
\kv{Also known as}{NSHM 2010 catalogue; Stirling et al. 2012 catalogue}
\kv{Coverage}{1840 to 2009}
\kv{Region}{New Zealand (nationwide)}
\kv{Bounding box}{New Zealand and its offshore margins (approx. 165E-180E / 178W, 48S-33S)}
\kv{Content}{hypocentral locations; quality information; historical observations; other}
\kv{Producer / contact}{GNS Science (Mark W. Stirling et al.); parent network GeoNet}
\kv{Data source}{Derived backbone catalogue built from the GeoNet/national instrumental earthquake catalogue plus the \textasciitilde{}170-year historical NZ earthquake record; homogenised to Mw and declustered, with duplicates, explosions and mining-induced events removed.}
\kv{Availability}{available (DOI) $\cdot$ DOI \href{https://doi.org/10.1785/0120110170}{\nolinkurl{10.1785/0120110170}}}
\end{itemize}
\par\vspace{2pt}Homogenised (Mw), declustered backbone earthquake catalogue underpinning the distributed-seismicity component of the 2010 New Zealand National Seismic Hazard Model. It was produced to give a temporally consistent magnitude basis (instrumental + \textasciitilde{}170-yr historical record) for probabilistic seismic hazard analysis across New Zealand. The catalogue is documented within and alongside the NSHM 2010 update.
\par\vspace{2pt}\textbf{Key references:} \cite{Stirling2012}.
\catentryrule
\subsubsection*{\hypertarget{c:78}{78.}~ Mw-Homogenised NZ Earthquake Catalogue for NSHM 2022 (Christophersen et al. 2022)}
{\small\itshape compilation $\cdot$ parameter-testing $\cdot$ other $\cdot$ reviewed}\par\vspace{2pt}
\begin{itemize}[leftmargin=1.3em,itemsep=0.6pt,topsep=1pt,parsep=0pt]
\kv{Also known as}{Christophersen et al. 2022 'consistent magnitudes' catalogue; GNS Science report 2021/42}
\kv{Coverage}{1931 to 2020}
\kv{Region}{New Zealand (nationwide)}
\kv{Bounding box}{New Zealand and offshore margins}
\kv{Magnitude}{Recomputes magnitudes and converts to a standardized Mw proxy (MwNZ); introduces MLNZ20 local-magnitude scale; standardization roughly halves the count of events with M\textgreater{}=4.95 and corrects 2011-2012 step changes.}
\kv{Events}{\textgreater{}250,000 earthquakes (magnitudes recomputed for \textgreater{}99\% of catalogue)}
\kv{Content}{quality information; other}
\kv{Producer / contact}{GNS Science (Annemarie Christophersen et al.); parent network GeoNet}
\kv{Data source}{GNS reprocessing of \textgreater{}250,000 NZ earthquakes using \textgreater{}2 million individual GeoNet waveforms; derives a new local magnitude scale (MLNZ20) and regressions to convert historical magnitudes to a Mw proxy for consistency over time.}
\kv{Availability}{available (DOI) $\cdot$ DOI \href{https://doi.org/10.21420/A2SN-XM76}{\nolinkurl{10.21420/A2SN-XM76}}}
\end{itemize}
\par\vspace{2pt}A GNS Science reprocessing that recomputes earthquake magnitudes for the NZ catalogue and standardises them to moment magnitude (Mw) so that earthquake size is consistent over the full catalogue period. It was produced to provide a temporally consistent magnitude backbone for the 2022 National Seismic Hazard Model, correcting step changes (e.g. the 2012 SeisComP transition) and the systematic ML-vs-Mw offset. It is the magnitude basis later folded into the Rollins integrated catalogue and the NSHM 2022 seismicity rate model.
\par{\small\itshape Related: A companion NSHM 2022 input catalogue; the augmented / integrated catalogue (Rollins et al. 2022) adds event classifications and depth-distribution models.}
\par\vspace{2pt}\textbf{Key references:} \cite{Christophersen2022,Christophersen2024}.
\catentryrule
\subsubsection*{\hypertarget{c:79}{79.}~ Integrated / Augmented Earthquake Catalogue for Aotearoa New Zealand (NSHM 2022; Rollins et al.)}
{\small\itshape compilation $\cdot$ geographical $\cdot$ other $\cdot$ reviewed}\par\vspace{2pt}
\begin{itemize}[leftmargin=1.3em,itemsep=0.6pt,topsep=1pt,parsep=0pt]
\kv{Also known as}{Augmented NZ earthquake catalogue; Rollins et al. catalogue; GNS Science report 2021/58; Integrated Earthquake Catalog Version 1; NSHM22 Mw-standardized catalogue}
\kv{Coverage}{1840 to 2020}
\kv{Region}{New Zealand (nationwide) incl. Hikurangi-Kermadec and Puysegur margins}
\kv{Bounding box}{Greater New Zealand region incl. Hikurangi-Kermadec and Puysegur subduction zones}
\kv{Magnitude}{Uses Christophersen et al. 2022 standardized Mw (MwNZ proxy) magnitudes; provides regional magnitude-frequency distributions and Mc in the companion paper (Rollins et al. 2024).}
\kv{Content}{hypocentral locations; relative relocations; focal mechanisms; quality information; historical observations; other}
\kv{Producer / contact}{GNS Science (Chris Rollins, Annemarie Christophersen, Kiran Kumar S. Thingbaijam, Matthew C. Gerstenberger et al.); parent network GeoNet}
\kv{Data source}{GeoNet national operational seismic catalogue, with event parameters (depths, uncertainties, focal mechanisms, locations) overwritten using refined estimates from regional relocation studies, the literature, the NZ Centroid Moment Tensor (CMT) database and global catalogues; magnitudes from the Christophersen et al. 2022 Mw standardization.}
\kv{Availability}{available (DOI) $\cdot$ DOI \href{https://doi.org/10.21420/XT4Y-WY45}{\nolinkurl{10.21420/XT4Y-WY45}}}
\end{itemize}
\par\vspace{2pt}An augmented, integrated NZ earthquake catalogue built for the 2022 National Seismic Hazard Model. It merges higher-quality depths, uncertainties, focal mechanisms and relocated hypocentres into the GeoNet catalogue, applies the Christophersen 2022 standardized Mw magnitudes, and classifies events as upper-plate (crustal), subduction-interface, or intraslab to feed the separate components of the Seismicity Rate Model. Revised parameters were imported for \textasciitilde{}60\% of events (incl. \textasciitilde{}92\% of 2000-2020 events).
\par{\small\itshape Related: Builds on the Mw-homogenised NSHM 2022 catalogue (Christophersen et al. 2022), adding event classifications and depth-distribution models.}
\par\vspace{2pt}\textbf{Key references:} \cite{Rollins2022,Rollins2025,Rollins2024}.
\catentryrule
\subsubsection*{\hypertarget{c:80}{80.}~ RSQSim Synthetic Earthquake Catalogue for New Zealand (NSHM2012 fault system, 276 kyr)}
{\small\itshape background $\cdot$ parameter-testing $\cdot$ other $\cdot$ non-reviewed}\par\vspace{2pt}
\begin{itemize}[leftmargin=1.3em,itemsep=0.6pt,topsep=1pt,parsep=0pt]
\kv{Also known as}{RSQSim Simulated Earthquake Catalog 5091; rundir5091; Shaw et al. 2022}
\kv{Coverage}{synthetic (276,000-year simulated span; published 2021) to synthetic (276,000-year simulated span)}
\kv{Region}{New Zealand (nationwide fault system)}
\kv{Bounding box}{New Zealand (NSHM 2012 fault model domain)}
\kv{Magnitude}{Simulated event magnitudes derived from rupture areas/slip on the fault model; full magnitude distribution over the 276 kyr simulation (no observational Mc).}
\kv{Content}{hypocentral locations; fault-geometry; other}
\kv{Producer / contact}{Bruce E. Shaw (Lamont-Doherty Earth Observatory, Columbia University); collaborators Bill Fry, Andrew Nicol, Andrew Howell, Matthew Gerstenberger (GNS Science)}
\kv{Data source}{Synthetic catalogue generated by the RSQSim rate-and-state earthquake simulator run on the New Zealand NSHM 2012 fault model (fault geometries, slip rates, rupture parameters); not based on observed events.}
\kv{Availability}{available (DOI) $\cdot$ DOI \href{https://doi.org/10.5281/zenodo.5534462}{\nolinkurl{10.5281/zenodo.5534462}}}
\end{itemize}
\par\vspace{2pt}A \textasciitilde{}276,000-year synthetic earthquake catalogue produced by the RSQSim rate-and-state simulator on the NZ NSHM 2012 fault system, generating deterministic finite-fault ruptures for physics-based seismic hazard. It was produced to test physics-based earthquake simulation as an alternative/complement to empirical hazard models for New Zealand. Distinct from the later NZNSHM22-fault RSQSim version (rundir5566).
\par\vspace{2pt}\textbf{Key references:} \cite{Shaw2021,Shaw2022}.
\catentryrule
\subsubsection*{\hypertarget{c:81}{81.}~ RSQSim Synthetic Earthquake Catalogue + Ground Motions for New Zealand (NZNSHM22 fault system)}
{\small\itshape background $\cdot$ parameter-testing $\cdot$ other $\cdot$ non-reviewed}\par\vspace{2pt}
\begin{itemize}[leftmargin=1.3em,itemsep=0.6pt,topsep=1pt,parsep=0pt]
\kv{Also known as}{RSQSim rundir5566; Shaw \& Milner 2025 NZ catalog}
\kv{Coverage}{synthetic (RSQSim simulated span; published 2025) to synthetic (RSQSim simulated span)}
\kv{Region}{New Zealand (nationwide fault system); plus California in the bundled record}
\kv{Bounding box}{New Zealand (NZNSHM22 fault model domain); record also includes California}
\kv{Magnitude}{Simulated event magnitudes from rupture area/slip on 2 km triangular fault patches; full simulated magnitude distribution.}
\kv{Content}{hypocentral locations; fault-geometry; other}
\kv{Producer / contact}{Bruce E. Shaw (Columbia University) \& Kevin R. Milner (USGS); Christine Goulet (USGS) co-author on the companion paper}
\kv{Data source}{Synthetic catalogue, rupture slip-time histories and simulated broadband ground motions generated by the RSQSim simulator on the 2022 NZ NSHM (NZNSHM22) fault system discretized into 2 km triangular patches; bundled in the same Zenodo record as a California (UCERF3, rundir5652) catalogue.}
\kv{Availability}{available (DOI) $\cdot$ DOI \href{https://doi.org/10.5281/zenodo.14532399}{\nolinkurl{10.5281/zenodo.14532399}}}
\end{itemize}
\par\vspace{2pt}An updated physics-based synthetic earthquake catalogue with rupture slip-time histories and simulated broadband ground motions on the 2022 NZ NSHM fault system (NZ rundir5566). It extends the earlier NSHM2012-fault RSQSim NZ catalogue to the current hazard-model fault geometry and adds ground-motion simulations, supporting comparison of physics-based simulators against empirical ground-motion models. The Zenodo record bundles both the NZ (rundir5566) and a California (rundir5652) catalogue.
\par\vspace{2pt}\textbf{Key references:} \cite{Shaw2025,Shaw2025a}.
\catentryrule
\section{Conclusion}\label{sec:conclusion}
This report inventories 81 earthquake catalogues developed for, or covering, Aotearoa New Zealand with coverage spanning 1960--2026, profiled to a common submission schema. The exercise confirms both the strength of the national record and the scale of the discoverability gap below it: temporary deployments and regional catalogues, often the most detailed local datasets available, are also the hardest to find, compare, and cite. Populating a lightweight, openly versioned registry with descriptors such as these is a low-cost, near-term step that would make the country's earthquake science measurably more findable, comparable, and reusable. This inventory is offered as a starting population of that registry and as a reproducible reference that can be extended as the long tail is filled in.
\appendix
\section{The Catalogue Submission Profile (schema)}\label{app:schema}
Each entry is described by the minimum-attributes submission profile proposed for the national \textquotedblleft catalogue of catalogues\textquotedblright. Required fields are kept to a small core to lower the barrier to contribution; controlled drop-downs are preferred over free text.
\begin{itemize}[leftmargin=1.4em,itemsep=1pt]
\item \textbf{Identification and contact:} catalogue name; contact person(s); data source; brief description; additional notes.
\item \textbf{Coverage:} start date; end date (or \emph{continuous/ongoing}); bounding box (lat/long); depth range; region of interest; magnitude information (range and/or completeness).
\item \textbf{Content present} (tick all): hypocentral locations; relative relocations; phase arrival-time picks; amplitude picks; fault-geometry; historical observations; quality information; other.
\item \textbf{Classification} (controlled): catalogue type (background / aftershock / induced / historical / compilation / large-events-only); focus (geographical / parameter-testing / particular-event-type); detection method (manual / automatic / template-matching / machine-learning / felt-reports / other); review status (non-reviewed / partially reviewed / reviewed).
\item \textbf{Availability:} available (DOI) / available on request / not available; plus DOI or dataset link where one exists.
\end{itemize}
This profile complements the richer event-level metadata profile (aligned with QuakeML and FDSN conventions) used for catalogues ingested in full.

\addcontentsline{toc}{section}{References}
\bibliographystyle{plainnat}
\bibliography{nz_catalogues}

\end{document}